% Supporting decision-making in organisations — Computer Science Upper Sixth — Cours signé OnBuch+
% Official MINESEC syllabus. Compiler avec : tectonic CSC13-supporting-decision-making-in-organisations.tex
\newcommand{\DOCMATIERE}{Computer Science}
\newcommand{\DOCNIVEAU}{Upper Sixth}
\newcommand{\DOCMODULE}{Upper Sixth — Computer Science}
\newcommand{\DOCLECON}{13}
\newcommand{\DOCTITRE}{Supporting decision-making in organisations}
% OnBuch+ — préambule « cours signés » ENRICHI (compilé avec tectonic / XeLaTeX)
% Système visuel « Pop » d'OnBuch+ : fond crème (jamais blanc), bordures encre
% épaisses, ombres « dures » décalées, titres Archivo Black en capitales.
% Version MATHÉMATIQUES : graphes (pgfplots/tikz), boîtes pédagogiques
% (définition, propriété, méthode, attention, le savais-tu, exercice résolu,
% exercice, TP, bilan), page de garde et signature OnBuch+ ancrée en bas.
%
% Macros attendues dans main.tex AVANT \input{preamble} :
%   \DOCMATIERE  \DOCNIVEAU  \DOCMODULE  \DOCLECON (numéro)  \DOCTITRE
\documentclass[11pt]{article}

\usepackage{fontspec}
\usepackage[english]{babel}
\usepackage[a4paper,margin=2cm,top=3.4cm,bottom=2.8cm,headheight=1.4cm,footskip=1.3cm]{geometry}
\usepackage{amsmath,amssymb,amsfonts}
\usepackage{mathtools} % \xmapsto, \coloneqq... souvent utilisés par les modèles
\usepackage{cancel} % simplifications barrées (bilans de forces)
% \abs{z} (module, valeur absolue) et \norm{u} (norme), utilisés par les modèles
\providecommand{\abs}{}\let\abs\relax\DeclarePairedDelimiter{\abs}{\lvert}{\rvert}
\providecommand{\norm}{}\let\norm\relax\DeclarePairedDelimiter{\norm}{\lVert}{\rVert}
\usepackage[table]{xcolor} % [table] : fournit aussi \rowcolors (alternance de lignes)
\usepackage{pagecolor}
\usepackage{tikz}
\usetikzlibrary{arrows.meta,positioning,calc,shapes.geometric,shapes.misc,shapes.arrows,decorations.pathmorphing,decorations.pathreplacing,decorations.markings,patterns,shadows,mindmap,trees,fit,backgrounds,matrix,intersections,angles,quotes}
\usepackage{pgfplots}
\usepackage[version=4]{mhchem} % équations nucléaires \ce{^{235}_{92}U}
\usepackage[european,siunitx]{circuitikz} % schémas électriques
\pgfplotsset{compat=1.18}
\usepgfplotslibrary{fillbetween,groupplots} % aires entre courbes (calcul intégral)
\usepackage{enumitem}
\usepackage{fancyhdr}
% [most] charge toutes les bibliothèques tcolorbox (listings, minted, raster,
% documentation...) et gonfle inutilement la mémoire TeX sur un document long
% (~300 boîtes) : on ne charge que ce qui est réellement utilisé.
\usepackage[skins,breakable]{tcolorbox}
\usepackage{graphicx}
\usepackage{colortbl}
\usepackage{booktabs}
\usepackage{tabularx}
\usepackage{diagbox} % case d'en-tête diagonale des tableaux à double entrée
% Colonnes extensibles centrée / gauche / droite (C, L, R), souvent utilisées par les modèles
\newcolumntype{C}{>{\centering\arraybackslash}X}
\newcolumntype{L}{>{\raggedright\arraybackslash}X}
\newcolumntype{R}{>{\raggedleft\arraybackslash}X}
\usepackage{array}
\usepackage{multicol}
\usepackage{siunitx}
% parse-numbers=false : accepte aussi \SI{2,5 \times 10^{-3}}{...}
\sisetup{locale=UK,output-decimal-marker={.},group-minimum-digits=4,parse-numbers=false}
\DeclareSIUnit\ohm{\ensuremath{\symup{\Omega}}} % U+2126 absent de la police
\DeclareSIUnit\division{div} % réglages d'oscilloscope (ms/div, V/div)
\DeclareSIUnit\tr{tr} % tours (vitesse de rotation, ex. \SI{1500}{\tr\per\minute})
\usepackage{qrcode}
% \degree : commande fréquemment utilisée par les modèles pour un angle ou une
% température en dehors de \SI{}{} (non fournie par les paquets chargés).
\providecommand{\degree}{^{\circ}}
% \qed (fin de démonstration), utilisé par les modèles sans amsthm
\providecommand{\qed}{\hfill$\square$}
% \barre : double barre des valeurs interdites dans un tableau de variations (tkz-tab)
\providecommand{\barre}{\ensuremath{\|}}
% environnement proof (amsthm non chargé) utilisé par les modèles
\ifdefined\proof\else\newenvironment{proof}[1][Proof]{\par\noindent\textit{#1.}\ }{\qed\par}\fi
\usepackage{lastpage}
\usepackage{hyperref}
\hypersetup{hidelinks}

% ============================================================
% Palette « Pop » OnBuch+ — mêmes hex que la classe P (plus_ui.dart)
% ============================================================
\definecolor{poporange}{HTML}{F59321}
\definecolor{popdark}{HTML}{E07A0C}
\definecolor{popcream}{HTML}{FFF6E8}
\definecolor{popink}{HTML}{1C1714}
\definecolor{poppurple}{HTML}{7A5AE0}
\definecolor{popgreen}{HTML}{2E9E6A}
\definecolor{poppink}{HTML}{E0457B}
\definecolor{popblue}{HTML}{2F7DE1}
\definecolor{popgold}{HTML}{C98A12}
\definecolor{popmuted}{HTML}{8A7F76}
\definecolor{popline}{HTML}{EADFD2}
\definecolor{popgray}{HTML}{8A7F76}
\colorlet{popgrey}{popgray}
% Alias tolérants (le modèle nomme parfois les couleurs différemment)
\colorlet{popwhite}{white}
\colorlet{popblack}{popink}
\colorlet{popred}{poppink}
\colorlet{popyellow}{popgold}
% Teintes claires pour fonds de tableaux / zones
\colorlet{popblueL}{popblue!10!white}
\colorlet{poporangeL}{poporange!14!white}
\colorlet{popgreenL}{popgreen!12!white}
\colorlet{poppurpleL}{poppurple!12!white}
\colorlet{poppinkL}{poppink!12!white}
\colorlet{popgoldL}{popgold!14!white}

\pagecolor{popcream}

% ============================================================
% Typographie
% ============================================================
\defaultfontfeatures{Ligatures=TeX}
\setmainfont{PlusJakartaSans-Medium.ttf}[Path=../../../fonts/,BoldFont=PlusJakartaSans-Bold.ttf]
% Maths en sans-serif (Fira Math) avec chiffres et lettres droites de Plus
% Jakarta, pour que les formules se fondent dans le texte.
\usepackage[math-style=ISO,bold-style=ISO]{unicode-math}
\setmathfont{FiraMath-Regular.otf}[Path=../../../fonts/]
\setmathfont{PlusJakartaSans-Medium.ttf}[Path=../../../fonts/,range={up/num,up/{Latin,latin},\mathup}]
% (icomma retiré : décimales avec point en anglais)
% Caractères mathématiques Unicode absents des polices de texte (Plus Jakarta,
% Archivo Black) : rendus par leur équivalent mathématique, en texte comme en
% formule — sinon ils s'affichent en carré vide (ex. « Arithmétique dans ℤ »).
\usepackage{newunicodechar}
\newunicodechar{ℝ}{\ensuremath{\mathbb{R}}}
\newunicodechar{ℤ}{\ensuremath{\mathbb{Z}}}
\newunicodechar{ℂ}{\ensuremath{\mathbb{C}}}
\newunicodechar{ℕ}{\ensuremath{\mathbb{N}}}
\newunicodechar{ℚ}{\ensuremath{\mathbb{Q}}}
\newunicodechar{𝔻}{\ensuremath{\mathbb{D}}}
\newunicodechar{α}{\ensuremath{\alpha}}
\newunicodechar{β}{\ensuremath{\beta}}
\newunicodechar{γ}{\ensuremath{\gamma}}
\newunicodechar{δ}{\ensuremath{\delta}}
\newunicodechar{ε}{\ensuremath{\varepsilon}}
\newunicodechar{θ}{\ensuremath{\theta}}
\newunicodechar{λ}{\ensuremath{\lambda}}
\newunicodechar{μ}{\ensuremath{\mu}}
\newunicodechar{σ}{\ensuremath{\sigma}}
\newunicodechar{τ}{\ensuremath{\tau}}
\newunicodechar{φ}{\ensuremath{\varphi}}
\newunicodechar{ψ}{\ensuremath{\psi}}
\newunicodechar{ω}{\ensuremath{\omega}}
\newunicodechar{Σ}{\ensuremath{\Sigma}}
% majuscules produites par \MakeUppercase dans les titres (α-aminés -> Α)
\newunicodechar{Α}{\ensuremath{\alpha}}
\newunicodechar{Β}{\ensuremath{\beta}}
\newunicodechar{Γ}{\ensuremath{\Gamma}}
\newunicodechar{Δ}{\ensuremath{\Delta}}
\newunicodechar{Ε}{\ensuremath{\varepsilon}}
\newunicodechar{Θ}{\ensuremath{\Theta}}
\newunicodechar{Λ}{\ensuremath{\Lambda}}
\newunicodechar{Μ}{\ensuremath{\mu}}
\newunicodechar{Τ}{\ensuremath{\tau}}
\newunicodechar{Φ}{\ensuremath{\Phi}}
\newunicodechar{Ψ}{\ensuremath{\Psi}}
\newunicodechar{Ω}{\ensuremath{\Omega}}
\newunicodechar{↦}{\ensuremath{\mapsto}}
\newunicodechar{∈}{\ensuremath{\in}}
\newunicodechar{∉}{\ensuremath{\notin}}
\newunicodechar{∧}{\ensuremath{\wedge}}
\newunicodechar{⇔}{\ensuremath{\Leftrightarrow}}
\newunicodechar{⟺}{\ensuremath{\Longleftrightarrow}}
\newunicodechar{⇒}{\ensuremath{\Rightarrow}}
\newunicodechar{⟹}{\ensuremath{\Longrightarrow}}
\newunicodechar{∘}{\ensuremath{\circ}}
\newunicodechar{∪}{\ensuremath{\cup}}
\newunicodechar{∩}{\ensuremath{\cap}}
\newunicodechar{≡}{\ensuremath{\equiv}}
\newunicodechar{⊂}{\ensuremath{\subset}}
\newunicodechar{⊥}{\ensuremath{\perp}}
\newunicodechar{∃}{\ensuremath{\exists}}
\newunicodechar{∀}{\ensuremath{\forall}}
\newunicodechar{∛}{\ensuremath{\sqrt[3]{\,}}}
\newunicodechar{∜}{\ensuremath{\sqrt[4]{\,}}}
\newunicodechar{⃗}{\ensuremath{{}^{\to}}}
\newunicodechar{ⁿ}{\ifmmode^{n}\else\textsuperscript{\ensuremath{n}}\fi}
\newunicodechar{ˣ}{\ifmmode^{x}\else\textsuperscript{\ensuremath{x}}\fi}
\newunicodechar{ᵇ}{\ifmmode^{b}\else\textsuperscript{\ensuremath{b}}\fi}
\newunicodechar{ʳ}{\ifmmode^{r}\else\textsuperscript{\ensuremath{r}}\fi}
\newunicodechar{ᵘ}{\ifmmode^{u}\else\textsuperscript{\ensuremath{u}}\fi}
\newunicodechar{ʸ}{\ifmmode^{y}\else\textsuperscript{\ensuremath{y}}\fi}
\newunicodechar{ᵃ}{\ifmmode^{a}\else\textsuperscript{\ensuremath{a}}\fi}
\newunicodechar{ᵉ}{\ifmmode^{e}\else\textsuperscript{\ensuremath{e}}\fi}
\newunicodechar{ᵏ}{\ifmmode^{k}\else\textsuperscript{\ensuremath{k}}\fi}
\newunicodechar{ᵖ}{\ifmmode^{p}\else\textsuperscript{\ensuremath{p}}\fi}
\newunicodechar{ᴺ}{\ifmmode^{N}\else\textsuperscript{\ensuremath{N}}\fi}
\newunicodechar{ᶜ}{\ifmmode^{c}\else\textsuperscript{\ensuremath{c}}\fi}
\newunicodechar{ʰ}{\ifmmode^{h}\else\textsuperscript{\ensuremath{h}}\fi}
\newunicodechar{ᵠ}{\ifmmode^{q}\else\textsuperscript{\ensuremath{q}}\fi}
\newunicodechar{ₙ}{\ifmmode_{n}\else\textsubscript{\ensuremath{n}}\fi}
\newunicodechar{ₐ}{\ifmmode_{a}\else\textsubscript{\ensuremath{a}}\fi}
\newunicodechar{ᵢ}{\ifmmode_{i}\else\textsubscript{\ensuremath{i}}\fi}
\newunicodechar{ᵣ}{\ifmmode_{r}\else\textsubscript{\ensuremath{r}}\fi}
\newunicodechar{ₖ}{\ifmmode_{k}\else\textsubscript{\ensuremath{k}}\fi}
\newunicodechar{ᵦ}{\ifmmode_{\beta}\else\textsubscript{\ensuremath{\beta}}\fi}
\newunicodechar{ₓ}{\ifmmode_{x}\else\textsubscript{\ensuremath{x}}\fi}
\newfontfamily\popdisplay{ArchivoBlack-Regular.ttf}[Path=../../../fonts/]
\newfontfamily\poplogo{SpaceGrotesk-Bold.ttf}[Path=../../../fonts/]
\newfontfamily\popsemi{PlusJakartaSans-SemiBold.ttf}[Path=../../../fonts/]
\newfontfamily\popxbold{PlusJakartaSans-ExtraBold.ttf}[Path=../../../fonts/]
\newfontfamily\popxboldls{PlusJakartaSans-ExtraBold.ttf}[Path=../../../fonts/,LetterSpace=3.0]

\setlist[enumerate]{leftmargin=1.7em,itemsep=5pt,topsep=4pt}
\setlist[itemize]{leftmargin=1.7em,itemsep=4pt,topsep=4pt,label={\color{poporange}\textbullet}}
\setlength{\parindent}{0pt}
\setlength{\parskip}{5pt}
\renewcommand{\arraystretch}{1.25}

% pgfplots : style Pop par défaut
\pgfplotsset{
  every axis/.append style={axis line style={popink,line width=1pt},tick style={popink},
    grid style={popline,line width=0.5pt},label style={font=\small},tick label style={font=\footnotesize}},
  popcurve/.style={poporange,line width=1.8pt,no markers,smooth},
}
% Même style disponible dans un \draw TikZ simple (hors pgfplots)
\tikzset{popcurve/.style={poporange,line width=1.8pt}}
\tikzset{popgrid/.style={popline,very thin}}
\colorlet{popcurve}{poporange} % le nom du style est parfois utilisé comme couleur

% ============================================================
% Pill / squiggle
% ============================================================
\newcommand{\poppill}[3]{%
  \tikz[baseline=-0.55ex]\node[draw=popink,line width=1.2pt,rounded corners=2.6pt,%
    fill=#2,inner xsep=6.5pt,inner ysep=3pt]{%
    \popxboldls\fontsize{7}{7}\selectfont\color{#3}\MakeUppercase{#1}};%
}
\newcommand{\popsquiggle}[1]{%
  \tikz[baseline=-0.4ex]{
    \draw[#1,line width=2.6pt,line cap=round]
      plot [smooth] coordinates {(0,0.09) (0.34,0.01) (0.68,0.21) (1.02,0.01) (1.36,0.10)};
  }%
}
% Mot-clé surligné (marqueur orange sous le texte)
\newcommand{\cle}[1]{{\popxbold\color{popdark}#1}}

% ============================================================
% En-tête / pied
% ============================================================
\pagestyle{fancy}
\fancyhf{}
\renewcommand{\headrulewidth}{0pt}
\renewcommand{\footrulewidth}{0pt}
\lhead{\raisebox{-0.15ex}{\poplogo\fontsize{13}{13}\selectfont\color{popink}OnBuch\color{poporange}+}}
\rhead{\poppill{\DOCMATIERE\ \textbullet\ \DOCNIVEAU\ \textbullet\ Chapter \DOCLECON}{white}{popink}}
\lfoot{\popxbold\fontsize{7.6}{7.6}\selectfont\color{popmuted}\MakeUppercase{Signed course OnBuch+}\ \textbullet\ {\popsemi\DOCTITRE}}
\rfoot{\tikz[baseline=-0.6ex]\node[rounded corners=8.5pt,fill=popink,minimum height=17pt,minimum width=17pt,inner xsep=5pt,inner sep=0]{\popxbold\fontsize{8}{8}\selectfont\color{popcream}\thepage\,/\,\pageref*{LastPage}};}

% ============================================================
% Titres
% ============================================================
\newcommand{\chaptitre}[2]{%
  \begin{tcolorbox}[enhanced,colback=poporange,colframe=popink,boxrule=2.6pt,arc=11pt,
      drop shadow=popink,left=15pt,right=15pt,top=12pt,bottom=11pt,before skip=3pt,after skip=18pt]
    \poppill{#2}{popink}{popcream}\par\vspace{9pt}
    {\raggedright\hyphenpenalty=10000\exhyphenpenalty=10000\popdisplay\fontsize{22}{24}\selectfont\color{popcream}\MakeUppercase{#1}\par}
  \end{tcolorbox}
}
% Section numérotée : pastille orange avec le numéro + titre Archivo
\newcounter{popsec}
\newcommand{\coursec}[1]{%
  \stepcounter{popsec}\setcounter{popsub}{0}%
  \par\vspace{16pt}\noindent
  \begin{minipage}{\linewidth}
  \tikz[baseline=-0.7ex]\node[draw=popink,line width=1.6pt,fill=poporange,rounded corners=4pt,minimum size=22pt,inner sep=0]{\popdisplay\fontsize{12}{12}\selectfont\color{popcream}\thepopsec};%
  \hspace{8pt}{\popdisplay\fontsize{15}{17}\selectfont\color{popink}\MakeUppercase{#1}}\par
  \vspace{4pt}\noindent{\color{poporange}\rule{36pt}{3.6pt}}
  \end{minipage}\par\nopagebreak\vspace{8pt}
}
\newcounter{popsub}
\newcommand{\courssub}[1]{%
  \stepcounter{popsub}%
  \par\vspace{9pt}\noindent{\popxbold\fontsize{11.5}{13}\selectfont\color{popink}{\color{poporange}\thepopsec.\thepopsub}\hspace{6pt}#1}\par\nopagebreak\vspace{4pt}
}

% ============================================================
% Boîtes pédagogiques — même recette : fond blanc, bordure épaisse colorée,
% ombre dure encre, badge plein. Titre optionnel [#1].
% ============================================================
\tcbset{popbase/.style={breakable,enhanced,colback=white,boxrule=2.3pt,arc=9pt,
  shadow={3pt}{-3pt}{0pt}{popink},left=14pt,right=14pt,top=11pt,bottom=11pt,before skip=11pt,after skip=15pt,
  fontupper=\popsemi\fontsize{10.6}{14.5}\selectfont\color{popink},
  fontlower=\popsemi\fontsize{10.6}{14.5}\selectfont\color{popink}}}
% Coche dessinée (le glyphe ✓ n'existe pas dans les polices du document)
\newcommand{\popcheck}[1]{\tikz[baseline=-0.5ex]\draw[#1,line width=1.6pt,line cap=round,line join=round](-0.13,0.0)--(-0.04,-0.09)--(0.13,0.1);}
\providecommand{\coche}{\popcheck{popgreen}}
\providecommand{\resultat}[1]{\par\smallskip\noindent\fcolorbox{popgreen}{popgreenL}{\parbox{\dimexpr\linewidth-2\fboxsep-2\fboxrule\relax}{\textbf{Result.} #1}}\par\smallskip}
\newcommand{\popboxhead}[3]{\poppill{#1}{#2}{white}\ifx\relax#3\relax\else\hspace{6pt}{\popxbold\fontsize{10.6}{12}\selectfont\color{popink}#3}\fi\par\vspace{7pt}}

\newtcolorbox{aretenir}[1][]{popbase,colframe=popgreen,before upper={\popboxhead{Key points}{popgreen}{#1}}}
\newtcolorbox{exemplebox}[1][]{popbase,colframe=poporange,before upper={\popboxhead{Exemple}{poporange}{#1}}}
\newtcolorbox{definition}[1][]{popbase,colframe=popblue,before upper={\popboxhead{Definition}{popblue}{#1}}}
\newtcolorbox{propriete}[1][]{popbase,colframe=poppurple,before upper={\popboxhead{Property}{poppurple}{#1}}}
\newtcolorbox{methode}[1][]{popbase,colframe=popgold,before upper={\popboxhead{Method}{popgold}{#1}}}
\newtcolorbox{attention}[1][]{popbase,colframe=poppink,before upper={\popboxhead{Warning}{poppink}{#1}}}
\newtcolorbox{experience}[1][]{popbase,colframe=popblue,colback=popblueL,before upper={\popboxhead{Experiment}{popblue}{#1}}}
% « Le savais-tu ? » : carte violette pleine (contexte, culture, Cameroun)
\newtcolorbox{savaistu}[1][]{popbase,colback=poppurple,colframe=popink,
  fontupper=\popsemi\fontsize{10.4}{14.2}\selectfont\color{white},
  before upper={\poppill{Did you know?}{white}{poppurple}\ifx\relax#1\relax\else\hspace{6pt}{\popxbold\color{white}#1}\fi\par\vspace{7pt}}}
% Exercice résolu : énoncé en haut, \tcblower puis la solution sur fond crème
\newcounter{popexo}\newcounter{popexr}
\newtcolorbox{exoresolu}[1][]{popbase,colframe=poporange,colbacklower=popcream,
  segmentation style={popink,line width=1.2pt,dashed},
  before upper={\stepcounter{popexr}\popboxhead{Worked example \thepopexr}{poporange}{#1}},
  before lower={\poppill{Solution}{popink}{popcream}\par\vspace{6pt}}}
% Exercice (non résolu en place) — la correction est dans la partie Corrigés
\newtcolorbox{exercice}[1][]{popbase,colframe=popink,
  before upper={\stepcounter{popexo}\popboxhead{Exercise \thepopexo}{popink}{#1}}}
% Corrigé
\newtcolorbox{corrige}[1][]{popbase,colframe=popgreen,colback=popgreenL,
  before upper={\popboxhead{Solution}{popgreen}{#1}}}
% Objectifs / prérequis / bilan
\newtcolorbox{objectifs}{popbase,colframe=popink,colback=poporangeL,
  before upper={\popboxhead{Chapter objectives}{popink}{}}}
\newtcolorbox{prerequis}{popbase,colframe=popink,colback=popblueL,
  before upper={\popboxhead{Prerequisites}{popblue}{}}}
\newtcolorbox{bilan}{popbase,colframe=popink,colback=white,boxrule=2.8pt,
  before upper={\popboxhead{Fiche bilan}{poporange}{L'essentiel à connaître pour l'examen}}}
% Figure encadrée avec légende
\newtcolorbox{popfigure}[1][]{enhanced,breakable=false,colback=white,colframe=popline,boxrule=1.4pt,arc=8pt,
  left=10pt,right=10pt,top=10pt,bottom=8pt,before skip=10pt,after skip=12pt,halign=center,
  lower separated=false,halign lower=center,
  fontlower=\popsemi\fontsize{9}{11.5}\selectfont\color{popmuted},
  #1}
\newcommand{\legende}[1]{\par\vspace{6pt}{\popsemi\fontsize{9}{11.5}\selectfont\color{popmuted}#1}}

% ============================================================
% Signature OnBuch+
% ============================================================
\newcommand{\poplogoinline}[1]{{\poplogo\fontsize{#1}{#1}\selectfont\color{popink}OnBuch\color{poporange}+}}
\newcommand{\signatureonbuch}{%
  \begin{tcolorbox}[enhanced,colback=popink,colframe=popink,boxrule=2.2pt,arc=10pt,
      drop shadow=poporange,left=14pt,right=14pt,top=10pt,bottom=10pt,before skip=0pt,after skip=0pt]
    \tikz[baseline=-0.7ex]\node[circle,fill=popgreen,draw=popcream,line width=1.3pt,minimum size=22pt,inner sep=0]{\popcheck{white}};%
    \hspace{9pt}\begin{minipage}[c]{0.62\linewidth}
      {\popxbold\fontsize{12}{13}\selectfont\color{popcream}Signed\ \textbullet\ {\poplogo OnBuch\color{poporange}+}}\par\vspace{2pt}
      {\raggedright\popsemi\fontsize{8}{10.5}\selectfont\color{popline}\MakeUppercase{OnBuch+ teaching team}\\\MakeUppercase{Follows the official MINESEC syllabus}\par}
    \end{minipage}\hfill
    {\poplogo\fontsize{22}{22}\selectfont\color{popcream}OnBuch\color{poporange}+}
  \end{tcolorbox}%
}

% ============================================================
% Page de garde
% #1 titre · #2 matière · #3 niveau · #4 module · #5 n° leçon · #6 durée ·
% #7 description / objectifs (texte ou itemize)
% ============================================================
\newcommand{\pagegarde}[7]{%
  \thispagestyle{empty}
  \newgeometry{a4paper,margin=1.7cm,top=1.6cm,bottom=1.4cm}%
  \begin{tikzpicture}[remember picture,overlay]
    \fill[poporange!18] ([xshift=-3.2cm,yshift=-3.6cm]current page.north east) circle (4.6cm);
    \fill[poppurple!14] ([xshift=2.4cm,yshift=7.6cm]current page.south west) circle (3.8cm);
  \end{tikzpicture}%
  {\poplogo\fontsize{30}{30}\selectfont\color{popink}OnBuch\color{poporange}+}\hfill
  \raisebox{0.6ex}{\poppill{Signed course \textbullet\ Premium}{popink}{popcream}}\par
  \vspace{1pt}\popsquiggle{poppurple}\par
  \vspace{8pt}
  \begin{tcolorbox}[enhanced,colback=poporange,colframe=popink,boxrule=2.8pt,arc=13pt,
      drop shadow=popink,left=18pt,right=18pt,top=12pt,bottom=13pt,after skip=10pt]
    \poppill{#2 \textbullet\ #3}{popink}{popcream}\hspace{5pt}\poppill{#4}{white}{popink}\par\vspace{8pt}
    {\popdisplay\fontsize{14}{14}\selectfont\color{popink}CHAPTER #5}\par\vspace{4pt}
    {\raggedright\hyphenpenalty=10000\exhyphenpenalty=10000\popdisplay\fontsize{25}{27}\selectfont\color{popcream}\MakeUppercase{#1}\par}
    \vspace{8pt}
    {\popxbold\fontsize{9.5}{9.5}\selectfont\color{popink}\MakeUppercase{Suggested time: #6}}
  \end{tcolorbox}
  \begin{tcolorbox}[enhanced,colback=white,colframe=popink,boxrule=2.2pt,arc=10pt,
      drop shadow=popink,left=16pt,right=16pt,top=10pt,bottom=10pt,after skip=8pt]
    \poppill{In this chapter}{poppurple}{white}\par\vspace{5pt}
    \popsemi\fontsize{9.3}{12.6}\selectfont\color{popink}#7
  \end{tcolorbox}
  \noindent\begin{minipage}[t]{0.487\linewidth}
    \begin{tcolorbox}[enhanced,colback=popgreen,colframe=popink,boxrule=2.2pt,arc=10pt,
        drop shadow=popink,left=11pt,right=11pt,top=10pt,bottom=10pt,height=3.1cm,valign=center]
      \tcbox[colback=white,colframe=popink,boxrule=1.3pt,arc=5pt,boxsep=0pt,nobeforeafter,
        left=3pt,right=3pt,top=3pt,bottom=3pt,valign=center]{\qrcode[height=1.9cm]{https://play.google.com/store/apps/details?id=com.luvvix.onbuch&hl=en}}%
      \hspace{8pt}\begin{minipage}[c]{0.5\linewidth}
        {\popdisplay\fontsize{11}{12.5}\selectfont\color{white}\MakeUppercase{Download OnBuch}}\par\vspace{4pt}
        {\raggedright\popsemi\fontsize{8.4}{11}\selectfont\color{white}Free \textbullet\ Google Play\\AI tutor, past papers, results\par}
      \end{minipage}
    \end{tcolorbox}
  \end{minipage}\hfill
  \begin{minipage}[t]{0.487\linewidth}
    \begin{tcolorbox}[enhanced,colback=poppurple,colframe=popink,boxrule=2.2pt,arc=10pt,
        drop shadow=popink,left=11pt,right=11pt,top=10pt,bottom=10pt,height=3.1cm,valign=center]
      \tikz[baseline=-0.7ex]\node[circle,fill=white,draw=popink,line width=1.3pt,minimum size=1.6cm,inner sep=0]{\popdisplay\fontsize{24}{24}\selectfont\color{poppurple}+};%
      \hspace{8pt}\begin{minipage}[c]{0.6\linewidth}
        {\popdisplay\fontsize{11}{12.5}\selectfont\color{white}\MakeUppercase{Go OnBuch+}}\par\vspace{4pt}
        {\raggedright\popsemi\fontsize{8.4}{11}\selectfont\color{white}All signed courses, unlimited AI tutor, PDF sheets — OnBuch+ tab in the app\par}
      \end{minipage}
    \end{tcolorbox}
  \end{minipage}
  \vspace{8pt}
  \popcontenu
  \vfill
  \signatureonbuch
  \restoregeometry
  \newpage
}

% Rangée « Ce cours contient » (page de garde)
\newcommand{\poptile}[3]{% #1 couleur #2 gros texte #3 légende
  \begin{tcolorbox}[enhanced,colback=#1,colframe=popink,boxrule=2pt,arc=9pt,shadow={2.5pt}{-2.5pt}{0pt}{popink},
      left=6pt,right=6pt,top=9pt,bottom=9pt,halign=center,valign=center,height=1.75cm,width=\linewidth,nobeforeafter]
    {\popdisplay\fontsize{12.5}{13}\selectfont\color{white}\MakeUppercase{#2}}\par\vspace{3pt}
    {\centering\popsemi\fontsize{7.8}{9.5}\selectfont\color{white}#3\par}
  \end{tcolorbox}}
\newcommand{\popcontenu}{%
  \par\noindent{\popxboldls\fontsize{7.5}{7.5}\selectfont\color{popmuted}THIS SIGNED COURSE CONTAINS}\par\vspace{7pt}
  \noindent\begin{minipage}[t]{0.235\linewidth}\poptile{popblue}{Course}{complete, illustrated, step by step}\end{minipage}\hfill
  \begin{minipage}[t]{0.235\linewidth}\poptile{poporange}{Methods}{and worked examples}\end{minipage}\hfill
  \begin{minipage}[t]{0.235\linewidth}\poptile{poppink}{Exercises}{exam style + solutions}\end{minipage}\hfill
  \begin{minipage}[t]{0.235\linewidth}\poptile{popgreen}{Summary sheet}{the essentials to remember}\end{minipage}\par}

% Sommaire
\newtcolorbox{sommaire}{enhanced,colback=white,colframe=popink,boxrule=2.2pt,arc=10pt,
  drop shadow=popink,left=18pt,right=18pt,top=13pt,bottom=13pt,before skip=6pt,after skip=20pt,
  before upper={\poppill{Contents}{poppurple}{white}\par\vspace{9pt}\popsemi\fontsize{10.6}{14.5}\selectfont\color{popink}}}

% Signature de fin de cours — poussée tout en bas de la dernière page
\newcommand{\findecours}{%
  \par\vfill\nopagebreak
  \begin{center}\popsquiggle{poporange}\end{center}\vspace{4pt}
  \signatureonbuch
}
\AtBeginDocument{\def\llbracket{\mathopen{[\mkern-3.5mu[}}\def\rrbracket{\mathclose{]\mkern-3.5mu]}}} % intervalles d'entiers (glyphe ⟦ absent de la police)
% Glyphes absents de Fira Math : redéfinis à partir de glyphes disponibles
\AtBeginDocument{%
  \def\setminus{\mathbin{\backslash}}\let\smallsetminus\setminus
  \def\vdots{\mathord{\vbox{\baselineskip3.2pt\lineskiplimit0pt\kern4pt\hbox{.}\hbox{.}\hbox{.}}}}%
  \def\ddots{\mathinner{\mkern1mu\raise6.5pt\hbox{.}\mkern2mu\raise3.6pt\hbox{.}\mkern2mu\raise0.7pt\hbox{.}\mkern1mu}}%
  \def\triangle{\mathord{\scalebox{0.9}{$\symup{\Delta}$}}}%
  \def\nabla{\mathord{\raisebox{\depth}{\scalebox{1}[-1]{$\symup{\Delta}$}}}}%
  \def\checkmark{\popcheck{popdark}}%
  \def\star{\mathbin{\ast}}\def\bigstar{\mathord{\ast}}%
  \def\diamond{\mathbin{\scalebox{0.6}{$\Diamond$}}}%
  \def\bigcup{\mathop{\mathchoice{\raisebox{-0.3em}{\scalebox{1.7}{$\cup$}}}{\scalebox{1.25}{$\cup$}}{\cup}{\cup}}\displaylimits}%
  \def\bigcap{\mathop{\mathchoice{\raisebox{-0.3em}{\scalebox{1.7}{$\cap$}}}{\scalebox{1.25}{$\cap$}}{\cap}{\cap}}\displaylimits}%
  \def\lesseqqgtr{\mathrel{\lessgtr}}\def\bigoplus{\mathop{\oplus}\displaylimits}%
}
\providecommand{\ssi}{if and only if} % abréviation fréquente
\AtBeginDocument{\providecommand{\wideparen}{\overparen}} % arc de cercle

% Fonctions usuelles que les modèles utilisent sans que amsmath les fournisse
\providecommand{\arsinh}{\operatorname{arsinh}}\providecommand{\arcosh}{\operatorname{arcosh}}
\providecommand{\artanh}{\operatorname{artanh}}\providecommand{\arsech}{\operatorname{arsech}}
\providecommand{\arcsch}{\operatorname{arcsch}}\providecommand{\arcoth}{\operatorname{arcoth}}
\providecommand{\arcsinh}{\operatorname{arsinh}}\providecommand{\arccosh}{\operatorname{arcosh}}
\providecommand{\arctanh}{\operatorname{artanh}}\providecommand{\sech}{\operatorname{sech}}
\providecommand{\csch}{\operatorname{csch}}\providecommand{\arcsec}{\operatorname{arcsec}}
\providecommand{\arccsc}{\operatorname{arccsc}}\providecommand{\arccot}{\operatorname{arccot}}
\providecommand{\sgn}{\operatorname{sgn}}\providecommand{\lcm}{\operatorname{lcm}}
\providecommand{\Var}{\operatorname{Var}}\providecommand{\Cov}{\operatorname{Cov}}

%% FIGFIX-BEGIN v5 — correctifs globaux des figures (docs/onbuch-generation/tools/figfix)
\usepackage{adjustbox}\usepackage{etoolbox}\tcbuselibrary{raster}
% toute figure TikZ/pgfplots plus large que la ligne est réduite à la largeur disponible
\BeforeBeginEnvironment{tikzpicture}{\begin{adjustbox}{max width=\linewidth}}
\AfterEndEnvironment{tikzpicture}{\end{adjustbox}}
% légendes pgfplots lisibles par-dessus les courbes
\pgfplotsset{every axis/.append style={legend style={fill=white,fill opacity=0.92,text opacity=1,draw=black!15,inner sep=3pt}}}
% étiquettes d'arêtes (syntaxe edge["..."]) avec fond blanc
\tikzset{every edge quotes/.append style={fill=white,inner sep=1.2pt}}
%% FIGFIX-END
\begin{document}

\pagegarde{Supporting decision-making in organisations}{Computer Science}{Upper Sixth}{Upper Sixth — Computer Science}{13}{5 periods}{This chapter introduces information systems as the nervous system of modern organisations. It shows how data is captured, verified and validated, and how different information systems support decision-making at the operational, tactical and strategic levels of an organisation.
\vspace{4pt}
\begin{multicols}{2}\small
\begin{itemize}
\item By the end of this chapter I can define an information system and distinguish data, information and knowledge.
\item By the end of this chapter I can describe the components (hardware, software, data, people, procedures) and functions of an information system.
\item By the end of this chapter I can identify and compare manual and automated data capture methods (keyboard, OMR, OCR, MICR, barcodes, RFID, biometrics, sensors).
\item By the end of this chapter I can choose an appropriate data capture method for a given real-life situation and justify the choice.
\item By the end of this chapter I can explain the difference between data verification and data validation and give examples of each.
\item By the end of this chapter I can apply validation checks (presence, range, type, format, length, check digit) to detect errors in a data set.
\end{itemize}\end{multicols}}

\begin{sommaire}
\begin{multicols}{2}
\begin{enumerate}
\item Information Systems: Concepts and Components
\item Data Capture
\item Data Verification and Validation
\item Organisational Levels and Information Systems for Decision-Making
\item Integration activity
\item Methods and worked examples
\item Exercises
\item Solutions to the exercises
\item Summary sheet
\end{enumerate}
\end{multicols}
\end{sommaire}

\begin{objectifs}
\begin{itemize}
\item By the end of this chapter I can define an information system and distinguish data, information and knowledge.
\item By the end of this chapter I can describe the components (hardware, software, data, people, procedures) and functions of an information system.
\item By the end of this chapter I can identify and compare manual and automated data capture methods (keyboard, OMR, OCR, MICR, barcodes, RFID, biometrics, sensors).
\item By the end of this chapter I can choose an appropriate data capture method for a given real-life situation and justify the choice.
\item By the end of this chapter I can explain the difference between data verification and data validation and give examples of each.
\item By the end of this chapter I can apply validation checks (presence, range, type, format, length, check digit) to detect errors in a data set.
\item By the end of this chapter I can describe the three organisational levels (operational, tactical, strategic) and the type of decisions taken at each.
\item By the end of this chapter I can match information systems (TPS, MIS, DSS, ESS) to organisational levels and explain their roles in decision-making.
\item By the end of this chapter I can analyse a simple organisation and propose an information system architecture supporting its decisions.
\end{itemize}
\end{objectifs}

\begin{prerequis}
\begin{itemize}
\item Basic computer literacy: hardware, software and data (Lower Sixth)
\item Notions of databases: records, fields, files (Lower Sixth)
\item Spreadsheet skills: entering, sorting and filtering data (Lower Sixth)
\item General knowledge of how Cameroonian businesses, schools and administrations operate
\end{itemize}
\end{prerequis}

\newpage

\coursec{Information Systems: Concepts and Components}

Every morning at the Mfoundi market in Yaoundé, a wholesaler must decide how many bags of plantain to order: too few and customers leave disappointed, too many and the stock rots by evening. Meanwhile, a few streets away, the regional hospital director must decide how to allocate beds, and the Minister of Finance must decide next year's budget. All of them, from the trader to the minister, rely on one thing: timely, accurate information. In the 21st century, that information increasingly flows through computerised information systems. How do organisations capture data, make sure it is correct, and turn it into decisions at every level? That is the question of this chapter.

\courssub{What is an Information System?}

\begin{definition}[Information System (IS)]
An \cle{information system} is an organised set of resources — people, hardware, software, data, and procedures — that collects, stores, processes, and distributes information to support decision-making, coordination, control, analysis, and visualisation in an organisation.
\end{definition}

An IS is not merely a computer program; it is the \emph{entire socio-technical ensemble} that turns raw facts into actionable insight. Consider the \emph{school report-card system} used in many Cameroonian secondary schools:
\begin{itemize}
    \item \textbf{People:} teachers (enter marks), administrators (approve), students/parents (receive reports).
    \item \textbf{Hardware:} school office computers, printers, the server hosting the database.
    \item \textbf{Software:} the mark-entry application, the database engine, the report generator.
    \item \textbf{Data:} student names, subject codes, raw marks, computed averages, rankings.
    \item \textbf{Procedures:} ``Marks must be entered by Friday 5\,pm; the head of department verifies; the principal signs off before printing.''
\end{itemize}
Remove any one component and the system fails: without procedures, teachers enter marks at different times causing inconsistency; without people, the software sits idle.

\begin{popfigure}
\begin{tikzpicture}[node distance=1.8cm, >=Stealth, font=\small]
    \node[draw, rounded corners, fill=popblueL, thick] (input) {Input / Capture};
    \node[draw, rounded corners, fill=popgreenL, thick, right of=input] (process) {Processing};
    \node[draw, rounded corners, fill=poporange!20, thick, right of=process] (storage) {Storage};
    \node[draw, rounded corners, fill=poppink!20, thick, right of=storage] (output) {Output / Dissemination};
    \draw[->, thick] (input) -- (process);
    \draw[->, thick] (process) -- (storage);
    \draw[->, thick] (storage) -- (output);
    % Feedback loop
    \draw[->, thick, poppurple] (output) -- ++(0,1.5) -| node[above, midway, text width=2.5cm, align=center] {Feedback \& Control} (input);
    % Data flowing into input
    \node[left of=input, xshift=-1.2cm, text width=2cm, align=center] (raw) {Raw Data};
    \draw[->, thick] (raw) -- (input);
    % Information flowing out of output
    \node[right of=output, xshift=1.2cm, text width=2cm, align=center] (info) {Information};
    \draw[->, thick] (output) -- (info);
\end{tikzpicture}
\legende{The core functions of an information system: Input $\to$ Processing $\to$ Storage $\to$ Output, with a Feedback loop for control.}
\end{popfigure}

\begin{propriete}[Five Core Functions of an IS]
Every information system performs five fundamental functions:
\begin{enumerate}
    \item \textbf{Input / Capture:} gathering raw data from internal or external sources (keyboards, scanners, sensors, mobile money USSD, web forms).
    \item \textbf{Processing:} transforming data into information (calculations, sorting, filtering, aggregation, applying business rules).
    \item \textbf{Storage:} holding data and information for future use (databases, files, cloud, backups).
    \item \textbf{Output / Dissemination:} presenting information to users in useful forms (reports, dashboards, SMS alerts, printed receipts).
    \item \textbf{Feedback and Control:} comparing output with goals, detecting errors, and adjusting input or processing (e.g. a variance report triggers a re-order).
\end{enumerate}
These functions form a continuous cycle; the feedback loop is what makes the system \emph{adaptive} rather than a one-shot data processor.
\end{propriete}

\courssub{Data, Information, and Knowledge: The Hierarchy}

\begin{definition}[Data]
\cle{Data} are raw, unprocessed facts, symbols, or signals without context. Examples: \texttt{37}, \texttt{Douala}, \texttt{2025-06-12}, \texttt{15000}.
\end{definition}

\begin{definition}[Information]
\cle{Information} is data that has been processed, organised, or structured so that it is meaningful and useful to the recipient. Example: \texttt{37°C recorded in Douala on 12 June 2025 at 14:00}.
\end{definition}

\begin{definition}[Knowledge]
\cle{Knowledge} is the understanding, experience, and insight that enables a person to interpret information and act effectively. Example: \emph{``When temperature exceeds 35°C in Douala, demand for bottled water rises by 20\%; we should increase stock now.''}
\end{definition}

\begin{popfigure}
\begin{tikzpicture}[scale=0.9, every node/.style={font=\small}]
    % Pyramid using three triangles stacked
    \coordinate (base) at (0,0);
    \coordinate (mid) at (0,3);
    \coordinate (top) at (0,5);
    % Base: Data
    \draw[fill=popblueL, thick] (-3,0) -- (3,0) -- (0,3) -- cycle;
    \node[text width=5cm, align=center, text=white] at (0,1) {\textbf{DATA}\\[2pt] Raw facts, no context\\[2pt] \textit{e.g. } 37, Douala, 15000};
    % Middle: Information
    \draw[fill=popgreen, thick] (-2,3) -- (2,3) -- (0,5) -- cycle;
    \node[text width=4cm, align=center, text=white] at (0,3.8) {\textbf{INFORMATION}\\[2pt] Processed, meaningful\\[2pt] \textit{e.g. } 37°C in Douala today};
    % Top: Knowledge
    \draw[fill=poppurple, thick] (-1,5) -- (1,5) -- (0,6.5) -- cycle;
    \node[text width=3cm, align=center, text=white] at (0,5.5) {\textbf{KNOWLEDGE}\\[2pt] Applied understanding\\[2pt] \textit{e.g. } Order more water stock};
    % Arrows on the side
    \draw[->, thick, popdark] (3.5,1) -- (3.5,2.5) node[midway, right] {Processing};
    \draw[->, thick, popdark] (2.5,4) -- (2.5,5) node[midway, right] {Interpretation};
\end{tikzpicture}
\legende{The Data–Information–Knowledge pyramid. Each level adds context, meaning, and actionability.}
\end{popfigure}

\begin{attention}[Confusing Data with Information]
In GCE exams, a frequent mistake is to define information as ``processed data'' without giving a concrete example, or to list raw figures as information. \textbf{Always pair a definition with a Cameroonian illustration:} ``Data: \texttt{15000}; Information: \texttt{15000 FCFA is the price of a bag of cement in Bamenda today}.''
\end{attention}

\begin{exemplebox}[Mobile Money Transaction]
When you dial \texttt{*150\#} on MTN MoMo or \texttt{*126\#} on Orange Money:
\begin{itemize}
    \item \textbf{Data captured:} your PIN, recipient number, amount \texttt{5000}.
    \item \textbf{Processing:} check balance $\ge$ 5000, verify recipient exists, apply transaction fee.
    \item \textbf{Information produced:} ``Transfer of 5000 FCFA to 6XX XX XX XX successful. New balance: 42,300 FCFA. Fee: 150 FCFA.''
    \item \textbf{Knowledge applied by user:} ``I still have enough for the electricity bill; I'll pay it now.''
\end{itemize}
\end{exemplebox}

\courssub{The Five Components of an Information System}

\begin{propriete}[The Five Components and Their Roles]
\label{prop:five-components}
\begin{center}
\begin{tabularx}{\linewidth}{|l|X|X|}\hline
\rowcolor{popblueL}
\textbf{Component} & \textbf{Role} & \textbf{Cameroonian Example} \\\hline
\textbf{Hardware} & Physical devices that capture, process, store, output & School server, POS terminal at a supermarket, smartphone for MoMo \\\hline
\textbf{Software} & Programs that instruct hardware: OS, DBMS, applications & Mark-entry app, \texttt{Excel} for budgets, \texttt{MySQL} for student records \\\hline
\textbf{Data} & Raw material: facts stored in databases, files, spreadsheets & Student IDs, marks, fee balances, inventory quantities \\\hline
\textbf{People} & Users (operational, managerial, strategic) + IT staff (admins, developers) & Teachers, bursar, principal, network technician \\\hline
\textbf{Procedures} & Rules, policies, step-by-step instructions for using the IS & ``Enter marks by Friday'', ``Backup database daily at 22:00'' \\\hline
\end{tabularx}
\end{center}
\end{propriete}

\begin{popfigure}
\begin{tikzpicture}[node distance=2.2cm, >=Stealth, font=\small]
    \node[draw, circle, fill=popblue, text=white, thick, minimum size=2.2cm, align=center] (center){Information\\ System};
    \node[draw, rounded corners, fill=popblueL, thick, above of=center] (hw) {Hardware};
    \node[draw, rounded corners, fill=popgreenL, thick, right of=center] (sw) {Software};
    \node[draw, rounded corners, fill=poporange!20, thick, below of=center] (data) {Data};
    \node[draw, rounded corners, fill=poppink!20, thick, left of=center] (people) {People};
    \node[draw, rounded corners, fill=poppurple!20, thick, above left of=center, yshift=-0.5cm, xshift=-0.5cm] (proc) {Procedures};
    \foreach \comp in {hw, sw, data, people, proc}
        \draw[<->, thick, popdark] (center) -- (\comp);
    % Interactions between components
    \draw[->, dashed, popmuted] (hw) to[bend left=20] node[above, sloped, font=\tiny] {runs} (sw);
    \draw[->, dashed, popmuted] (sw) to[bend left=20] node[above, sloped, font=\tiny] {manages} (data);
    \draw[->, dashed, popmuted] (people) to[bend left=20] node[below, sloped, font=\tiny] {follow} (proc);
    \draw[->, dashed, popmuted] (proc) to[bend left=20] node[above, sloped, font=\tiny] {govern} (data);
\end{tikzpicture}
\legende{The five components of an IS interact continuously. Procedures govern how people use hardware and software to manage data.}
\end{popfigure}

\begin{methode}[Analysing Any Organisation's IS]
To describe or evaluate an information system in an exam or case study, ask these five questions:
\begin{enumerate}
    \item \textbf{Who captures what data?} (Identify people, input devices, source documents)
    \item \textbf{Where and how is it stored?} (Database type, file structure, backup strategy)
    \item \textbf{What processing occurs?} (Calculations, validations, reports generated)
    \item \textbf{Who receives what output, and in what form?} (Screen, print, SMS, dashboard; operational vs. managerial users)
    \item \textbf{What procedures govern the cycle?} (Schedules, authorisation levels, error-handling, audit trail)
\end{enumerate}
Apply this checklist to the school report-card system, a hospital patient-record system, or a customs clearance system at Douala port.
\end{methode}

\courssub{Types of Information Systems — Overview}

While Section 4 will treat each type in depth, here is the \emph{big picture}:

\begin{center}
\begin{tabularx}{\linewidth}{|l|X|X|}\hline
\rowcolor{popblueL}
\textbf{Type} & \textbf{Primary Users} & \textbf{Typical Decisions Supported} \\\hline
Transaction Processing System (TPS) & Operational staff (clerks, cashiers) & Daily routine: record sales, register students, process payroll \\\hline
Management Information System (MIS) & Middle managers (heads of department, principals) & Periodic: monitor budgets, track attendance trends, stock re-order \\\hline
Decision Support System (DSS) & Senior managers, analysts & Semi-structured: ``What if we increase fees by 10\%?'', site selection \\\hline
Executive Support System (ESS) & Top executives (ministers, directors general) & Strategic: long-term policy, national budget allocation, 5-year plan \\\hline
\end{tabularx}
\end{center}

\begin{savaistu}[From Paper Files to Computerised Customs]
Until the early 2000s, Cameroon's customs service (\emph{Douanes}) relied on paper manifests stored in vast warehouses at Douala and Kribi ports. Clearing a container could take weeks. The introduction of the \emph{ASYCUDA++} system (Automated System for Customs Data) — a UNCTAD-supplied TPS/MIS — reduced clearance to days, cut corruption opportunities, and increased revenue collection. This is a classic exam illustration of \emph{computerisation benefits}.
\end{savaistu}

\courssub{Benefits and Limitations of Computerised IS}

\begin{propriete}[Benefits of Computerised Information Systems]
\begin{itemize}
    \item \textbf{Speed:} Thousands of transactions per second (e.g. MoMo processes millions daily).
    \item \textbf{Accuracy:} Automated calculations eliminate arithmetic errors (if validation is correct).
    \item \textbf{Storage capacity:} Terabytes on a single server vs. rooms of paper files.
    \item \textbf{Sharing and accessibility:} Authorised users access same data simultaneously from different sites (e.g. regional hospitals viewing central drug stock).
    \item \textbf{Better decisions:} Timely, consistent information enables evidence-based choices.
\end{itemize}
\end{propriete}

\begin{propriete}[Limitations and Risks]
\begin{itemize}
    \item \textbf{High initial cost:} Hardware, licences, development, training.
    \item \textbf{Training and change resistance:} Staff must learn new procedures; ``we've always done it on paper'' mindset.
    \item \textbf{Infrastructure dependence:} Power cuts (load-shedding), unstable internet in rural areas.
    \item \textbf{Security risks:} Unauthorised access, data breaches, ransomware, viruses.
    \item \textbf{System failure impact:} A crashed server halts the whole operation; paper fallback is often forgotten.
\end{itemize}
\end{propriete}

\courssub{Information Quality Criteria}

For information to support sound decisions, it must meet quality standards. The GCE syllabus expects you to know and illustrate these criteria.

\begin{definition}[Information Quality Criteria]
\begin{itemize}
    \item \textbf{Accurate:} Free from errors (e.g. a mark of 37 entered as 73 is inaccurate).
    \item \textbf{Complete:} All required fields present (e.g. a student record missing the date of birth).
    \item \textbf{Timely:} Available when needed (e.g. exam results released before university admission deadline).
    \item \textbf{Relevant:} Pertinent to the decision at hand (e.g. giving the principal a list of all students' shoe sizes is irrelevant for budget planning).
    \item \textbf{Understandable:} Clear format, language, and units (e.g. a report mixing FCFA and USD without conversion confuses the bursar).
\end{itemize}
\end{definition}

\begin{exemplebox}[Wrong Mark Entered]
In a school's mark-entry system, a teacher accidentally types \texttt{73} instead of \texttt{37} for a student in Mathematics.
\begin{itemize}
    \item \textbf{Accuracy violated:} The stored data is wrong.
    \item \textbf{Consequence:} The student's average is inflated, possibly affecting ranking, prize allocation, and even university admission.
    \item \textbf{Control needed:} Validation rule ``mark must be between 0 and 100'' would not catch this; a \emph{verification} step (double-entry or teacher review of a printed mark sheet) is required.
\end{itemize}
\end{exemplebox}

\courssub{Worked Example}

\begin{exoresolu}[Analysing a Rural Health Centre IS]
\textbf{Statement:} A rural health centre in the North Region uses a computerised system to manage patient consultations. Each morning, a nurse records patient details (name, age, village, symptoms) on paper forms. In the afternoon, a data-entry clerk types the forms into a desktop application connected to a local MySQL database. The doctor views a daily summary on screen. At month-end, the clerk prints a report of the top 10 diseases and sends it to the District Health Office via email. Power cuts are frequent; the centre has a small petrol generator but no UPS.

\textbf{Questions:}
\begin{enumerate}
    \item Identify the five components of this IS.
    \item List the five IS functions and give an example from the scenario for each.
    \item State two information quality criteria that are at risk because of the power cuts, and explain why.
    \item Propose one procedure to mitigate the risk of data loss during power cuts.
\end{enumerate}
\tcblower
\textbf{Detailed Solution:}
\begin{enumerate}
    \item \textbf{Five components:}
        \begin{itemize}
            \item Hardware: desktop PC(s), printer, generator, (missing UPS).
            \item Software: MySQL DBMS, data-entry application, OS (likely Windows/Linux).
            \item Data: patient records (name, age, village, symptoms), monthly disease summary.
            \item People: nurse (captures), data-entry clerk (inputs), doctor (uses output), District Health Office (receives report).
            \item Procedures: ``Morning: nurse fills paper forms; Afternoon: clerk enters data; Month-end: clerk prints and emails top-10 report.''
        \end{itemize}
    \item \textbf{Five functions with examples:}
        \begin{itemize}
            \item Input/Capture: Nurse writing patient details on paper forms.
            \item Processing: Application calculates age from date of birth, aggregates disease counts for monthly report.
            \item Storage: MySQL database holds all patient records.
            \item Output/Dissemination: Doctor's daily summary on screen; printed/emailed monthly top-10 report.
            \item Feedback/Control: Doctor notices a surge in malaria cases from the daily summary and requests extra rapid-test kits; District Office uses monthly report to allocate drugs.
        \end{itemize}
    \item \textbf{Quality criteria at risk due to power cuts:}
        \begin{itemize}
            \item \textbf{Completeness:} If power fails during data entry, the clerk's unsaved transaction is lost $\to$ patient record incomplete or missing.
            \item \textbf{Timeliness:} Frequent interruptions delay entry; the daily summary may be late or based on partial data.
        \end{itemize}
    \item \textbf{Mitigation procedure:} ``After every 10 records entered, the clerk must click `Save' and the system must auto-backup to a USB drive. At end of each session, the clerk verifies record count against paper forms before shutting down.'' (A UPS would be a hardware solution; the question asks for a \emph{procedure}.)
\end{enumerate}
\end{exoresolu}

\courssub{Exam Tips and Common Pitfalls}

\begin{aretenir}[Key Points for the Exam]
\begin{itemize}
    \item An IS = \textbf{Hardware + Software + Data + People + Procedures}. Never omit \emph{people} or \emph{procedures}.
    \item \textbf{Data $\neq$ Information $\neq$ Knowledge}. Use the pyramid: raw $\to$ meaningful $\to$ actionable.
    \item The five functions form a \textbf{cycle} with feedback; the feedback loop is what enables control.
    \item Quality criteria: \textbf{Accurate, Complete, Timely, Relevant, Understandable} — memorise the list and be ready to illustrate each.
    \item Computerised IS bring speed, accuracy, capacity, sharing, better decisions — but cost, training, power, security, and dependence are real limitations.
\end{itemize}
\end{aretenir}

\begin{attention}[Common Mistakes in GCE Answers]
\begin{itemize}
    \item Writing ``computer'' instead of ``hardware'' as a component (hardware includes servers, networks, printers, sensors…).
    \item Confusing \emph{validation} (checking data at entry) with \emph{verification} (checking data was copied correctly) — covered in Lesson 64.
    \item Listing TPS, MIS, DSS, ESS as ``components'' of an IS; they are \emph{types} of IS, not components.
    \item Forgetting that \textbf{procedures} include backup schedules, access rights, error-reporting steps — not just ``how to click buttons''.
\end{itemize}
\end{attention}

\begin{exercice}[Practice Questions]
\begin{enumerate}
    \item Define the term \emph{information system} and distinguish it from a \emph{computer system}.
    \item A small business uses a spreadsheet to track sales. Identify the five IS components in this context.
    \item Explain, with an example, the difference between data and information.
    \item State three criteria of quality information and describe a situation where each criterion is violated.
    \item Why is the feedback loop essential in an information system? Give a Cameroonian organisational example.
\end{enumerate}
\end{exercice}

\begin{corrige}[Exercise 1]
An information system is an organised set of resources (people, hardware, software, data, procedures) that collects, stores, processes, and distributes information to support decision-making and control. A computer system refers only to the hardware and software (the technical platform); it does not include the people, data, or procedures that give the system its organisational purpose.
\end{corrige}

\begin{corrige}[Exercise 2]
\begin{itemize}
    \item Hardware: the laptop/desktop running the spreadsheet.
    \item Software: the spreadsheet application (e.g. Excel, LibreOffice Calc) and the OS.
    \item Data: the sales figures, dates, product names, customer IDs entered in the cells.
    \item People: the owner/manager who enters data and reads charts; possibly an accountant who audits it.
    \item Procedures: ``Update the sheet every evening'', ``Save a backup copy on Google Drive every Friday'', ``Lock the file with a password''.
\end{itemize}
\end{corrige}

\begin{corrige}[Exercise 3]
Data: raw fact \texttt{2500}. Information: \texttt{2500 FCFA is the price of a kilogram of rice in Maroua market on 15 July 2025}. The data alone has no meaning; the information adds context (currency, item, location, date) making it useful for a buyer's decision.
\end{corrige}

\begin{corrige}[Exercise 4]
\begin{itemize}
    \item \textbf{Accurate:} violated when a mark of 42 is recorded as 24.
    \item \textbf{Timely:} violated when the monthly sales report arrives after the re-order deadline.
    \item \textbf{Relevant:} violated when the principal receives a detailed list of every student's blood type for a budget meeting.
\end{itemize}
\end{corrige}

\begin{corrige}[Exercise 5]
The feedback loop compares output with objectives and triggers corrective action. Without it, the system is open-loop and errors accumulate. Example: In a regional hospital pharmacy, the monthly stock report (output) shows anti-malarial drugs below minimum. The feedback loop prompts the pharmacist to place an emergency order (corrective input), preventing stock-out.
\end{corrige}

\coursec{Data Capture}

In the previous section we saw that an information system is a cycle: \cle{data} is collected, processed, stored and finally delivered as \cle{information} to support decision-making. Everything depends on the very first step. If the data entering the system is wrong, incomplete or late, then no matter how powerful the computer is, the output will be useless --- computer scientists summarise this with the famous expression \emph{garbage in, garbage out} (GIGO). This section studies that first step in detail: how raw data gets from the real world into the computer.

\courssub{What is Data Capture?}

Imagine the GCE Board after the Advanced Level examinations. Hundreds of thousands of answer sheets must become marks in a results database. Or think of a supermarket in Douala at closing time: every item sold during the day must be recorded so that the manager knows what to reorder. In both cases, data exists in the physical world (shaded boxes on a sheet, products on shelves) and must be transferred into a computer system. That transfer is called \cle{data capture}.

\begin{definition}[Data Capture]
\cle{Data capture} is the process of collecting raw data from a source and entering it into a computer system so that it can be processed and stored. It is the \textbf{input stage} of the information system cycle.
\end{definition}

We distinguish two broad families of capture methods:

\begin{itemize}
\item \textbf{Manual data capture}: a human being reads the data and types, clicks or dictates it into the system (keyboard, mouse, touchscreen, microphone).
\item \textbf{Automated data capture}: a device reads the data directly from a document, a card, a tag or the environment, with little or no human typing (scanners, card readers, RFID, sensors).
\end{itemize}

\begin{popfigure}
\begin{center}
\begin{tikzpicture}[
  box/.style={draw=popblue, thick, rounded corners, fill=popblueL, text width=4.6cm, align=center, minimum height=1cm, font=\small},
  leaf/.style={draw=popmuted, rounded corners, fill=white, text width=3.4cm, align=center, font=\footnotesize, minimum height=0.8cm},
  arr/.style={-{Stealth}, thick, popdark}]
\node[box, fill=popblue, text=white, font=\small\bfseries] (root) at (0,0) {Data Capture Methods};
\node[box, align=center] (man) at (-5.2,-2){Manual\\(human enters data)};
\node[box, align=center] (aut) at (5.2,-2){Automated\\(device reads data)};
\node[leaf] (m1) at (-8.6,-4) {Keyboard entry};
\node[leaf] (m2) at (-5.2,-4) {Mouse / touchscreen};
\node[leaf] (m3) at (-1.8,-4) {Voice dictation};
\node[leaf] (a1) at (1.8,-4) {Optical: OMR, OCR, barcode, QR};
\node[leaf] (a2) at (5.2,-4) {Magnetic / chip: stripe, MICR, smart card};
\node[leaf] (a3) at (8.6,-4) {RFID, biometrics, sensors};
\draw[arr] (root) -- (man);
\draw[arr] (root) -- (aut);
\draw[arr, popmuted] (man) -- (m1);
\draw[arr, popmuted] (man) -- (m2);
\draw[arr, popmuted] (man) -- (m3);
\draw[arr, popmuted] (aut) -- (a1);
\draw[arr, popmuted] (aut) -- (a2);
\draw[arr, popmuted] (aut) -- (a3);
\end{tikzpicture}
\end{center}
\legende{The two families of data capture methods, with typical devices.}
\end{popfigure}

\courssub{Manual Data Capture Methods}

\textbf{Keyboard entry.} The oldest and most flexible method: an operator reads a source document (a registration form, an invoice) and types the data. This is how a secretary enters students' names into a school database, or how a cashier at a small shop in Bamenda types prices on a till.

\textbf{Mouse, touchscreen and graphics input.} The user selects options, ticks boxes or draws directly on a screen. Mobile money agents in Cameroon capture transactions this way: they tap menus and amounts on a phone rather than typing long texts.

\textbf{Voice dictation (speech recognition).} The user speaks and software converts speech to text. Useful when hands are busy (a doctor dictating notes) or for users with disabilities.

\begin{center}
\begin{tabularx}{\linewidth}{|l|X|X|}
\hline
\rowcolor{popblueL} \textbf{Method} & \textbf{Advantages} & \textbf{Drawbacks} \\
\hline
Manual capture & Very flexible (any kind of data); cheap to start (a keyboard costs little); humans can interpret unclear handwriting. & Slow for large volumes; tiring; \textbf{error-prone} (typing mistakes); labour cost grows with volume. \\
\hline
Automated capture & Very fast; consistent accuracy; low cost per item at high volume; works 24/7 (sensors). & High initial equipment cost; needs standardised documents or tags; misreads still occur. \\
\hline
\end{tabularx}
\end{center}

\courssub{Automated Capture --- Optical Methods}

Optical methods ``read'' marks, characters or printed patterns using light.

\textbf{OMR (Optical Mark Recognition).} An OMR scanner detects the \emph{position of shaded marks} on a pre-printed form. It does not read letters; it only checks whether specific boxes are darkened. This is exactly how GCE multiple-choice answer sheets are marked: the candidate shades bubbles with a pencil, the scanner detects which bubbles are shaded, and software converts positions into answers and then into marks.

\textbf{OCR (Optical Character Recognition).} OCR goes further: software analyses the \emph{shapes} of printed or handwritten characters and converts them into editable text. Scanning a printed document and obtaining a Word file you can edit is OCR at work. Banks use OCR to read amounts on deposit slips.

\textbf{Barcode readers.} A barcode is a pattern of parallel black bars of varying widths encoding a number (usually the product code). A laser scanner at a supermarket checkout reads the code, the till looks up the price and description in the stock database, and the sale is recorded automatically. Libraries also barcode every book to speed up lending.

\textbf{QR codes.} A QR (Quick Response) code is a two-dimensional square pattern that stores much more data than a barcode (a web address, a payment reference). Mobile payment services and event tickets increasingly use QR codes scanned by phone cameras.

\begin{popfigure}
\begin{center}
\begin{tikzpicture}[
  stage/.style={draw=popgreen, thick, rounded corners, fill=popgreenL, text width=3.1cm, align=center, minimum height=1.1cm, font=\small},
  arr/.style={-{Stealth}, very thick, popgreen},
  note/.style={font=\footnotesize\itshape, text=popmuted, align=center}]
\node[stage] (sheet) at (0,0) {Candidate shades bubbles on GCE answer sheet};
\node[stage] (scan) at (4.6,0) {OMR scanner detects shaded positions};
\node[stage] (marks) at (9.2,0) {Software converts positions to answers and marks};
\node[stage] (res) at (13.4,0) {Results file stored in database};
\draw[arr] (sheet) -- (scan);
\draw[arr] (scan) -- (marks);
\draw[arr] (marks) -- (res);
\node[note] at (2.3,-1.1) {paper document};
\node[note] at (6.9,-1.1) {light sensors};
\node[note] at (11.3,-1.1) {digital data};
\end{tikzpicture}
\end{center}
\legende{OMR processing of a GCE multiple-choice answer sheet.}
\end{popfigure}

\begin{popfigure}
\begin{center}
\begin{tikzpicture}[font=\small]
% barcode
\foreach \x/\w in {0/0.06, 0.12/0.03, 0.2/0.09, 0.34/0.03, 0.42/0.06, 0.55/0.12, 0.72/0.03, 0.8/0.06, 0.92/0.03, 1.0/0.09, 1.15/0.03, 1.24/0.06, 1.36/0.03, 1.45/0.09, 1.6/0.03, 1.68/0.06, 1.8/0.03, 1.88/0.12, 2.05/0.03, 2.13/0.06, 2.25/0.03, 2.33/0.09, 2.48/0.03, 2.56/0.06}{
  \fill[popink] (\x,0) rectangle ++(\w,1.4);}
\node[below, font=\footnotesize] at (1.3,-0.05) {6 001234 567890};
\node[above, font=\footnotesize\itshape, text=popmuted] at (1.3,1.45) {Barcode on a product};
% flow
\node[draw=poporange, thick, rounded corners, fill=white, text width=2.6cm, align=center, minimum height=0.9cm] (reader) at (4.6,0.7) {Laser scanner reads code};
\node[draw=poporange, thick, rounded corners, fill=white, text width=2.6cm, align=center, minimum height=0.9cm] (db) at (8.2,0.7) {Stock database: price \& description};
\node[draw=poporange, thick, rounded corners, fill=white, text width=2.6cm, align=center, minimum height=0.9cm] (till) at (11.8,0.7) {Till: receipt + stock updated};
\draw[-{Stealth}, very thick, poporange] (2.7,0.7) -- (reader);
\draw[-{Stealth}, very thick, poporange] (reader) -- (db);
\draw[-{Stealth}, very thick, poporange] (db) -- (till);
\end{tikzpicture}
\end{center}
\legende{Structure of a barcode and the data flow at a supermarket checkout.}
\end{popfigure}

\courssub{Automated Capture --- Magnetic and Chip Methods}

\textbf{Magnetic stripe cards.} The brown stripe on the back of a bank card stores the account number as a magnetic pattern, read when the card is swiped. Cheap, but the stripe wears out and is easy to copy.

\textbf{MICR (Magnetic Ink Character Recognition).} Bank cheques carry a line of characters printed in special magnetic ink along the bottom (bank code, branch, account number, cheque number). MICR readers read these characters magnetically, which works even if the cheque is stamped or slightly dirty --- a key reason banks still trust it for clearing cheques.

\textbf{Smart cards (chip cards).} A smart card contains an embedded microchip that can store and even process data securely: SIM cards in mobile phones, modern bank cards (chip and PIN), and some student ID cards. The chip is far harder to copy than a magnetic stripe.

\courssub{Automated Capture --- Radio and Biometric Methods}

\textbf{RFID (Radio Frequency Identification).} An RFID tag is a tiny chip with an antenna that transmits its identity when it passes near a reader --- no contact, no line of sight needed. Uses include tracking stock in warehouses, electronic toll collection on motorways, and anti-theft tags on goods.

\textbf{Biometric capture.} Biometrics capture physical characteristics of a person: fingerprints, facial features, iris patterns. In Cameroon, biometric capture is used during national identity card enrolment and voter registration, where fingerprints and a photograph are captured so that each person is registered only once.

\courssub{Automated Capture --- Sensors and Data Logging}

\textbf{Sensors} measure physical quantities (temperature, humidity, pressure, position) and convert them directly into digital data, often at regular intervals --- this is called \cle{automatic data logging}. Examples: weather stations recording rainfall and temperature every hour; GPS trackers on inter-urban transport buses sending their position so the agency can monitor journeys; cold-chain sensors checking that vaccines in a health centre fridge stay at the right temperature. Sensors capture data continuously, without fatigue and without human intervention.

\courssub{Choosing a Capture Method}

No single method is ``best''; the choice depends on the situation.

\begin{methode}[Choosing a Data Capture Technique]
Evaluate each candidate method against five criteria:
\begin{enumerate}
\item \textbf{Volume of data}: thousands of forms per day demand automation; a few entries per day can stay manual.
\item \textbf{Speed required}: a checkout queue or an election night needs near-instant capture.
\item \textbf{Accuracy}: how costly is an error? Bank cheques justify MICR; a suggestion box does not.
\item \textbf{Cost}: equipment, tags, training, maintenance --- balanced against labour saved.
\item \textbf{Environment}: dust, rain, darkness, movement (a barcode needs a clean, visible label; RFID works through dirt; sensors suit remote sites).
\end{enumerate}
\end{methode}

\begin{propriete}[OMR vs OCR vs MICR]
The three ``recognition'' technologies read different things and serve different purposes (learn this distinction --- it is a classic examination question):
\begin{itemize}
\item \textbf{OMR} reads the \emph{position of shaded marks} on pre-printed forms $\rightarrow$ multiple-choice answer sheets, surveys, lotteries.
\item \textbf{OCR} reads the \emph{shapes of printed or handwritten characters} $\rightarrow$ scanning documents into editable text, reading forms.
\item \textbf{MICR} reads characters printed in \emph{magnetic ink} $\rightarrow$ the code line on bank cheques.
\end{itemize}
\end{propriete}

\begin{exemplebox}[Matching Methods to Situations]
\begin{itemize}
\item Marking 50\,000 GCE MCQ sheets: \textbf{OMR} (huge volume, standard forms, high accuracy needed).
\item Checkout in a Yaound\'e supermarket: \textbf{barcode reader} (fast, cheap, products already labelled).
\item Paying a bus fare with a contactless card: \textbf{RFID / smart card}.
\item Recording rainfall on Mount Cameroon: \textbf{sensors with data logging} (remote, continuous).
\item Entering a one-off letter into a computer: \textbf{keyboard} (tiny volume, irregular data).
\end{itemize}
\end{exemplebox}

\courssub{Sources of Error at Capture}

Even well-designed capture systems produce errors. Two error types are so common in examinations that they have names:

\begin{definition}[Transcription Error]
A \cle{transcription error} occurs when data is copied incorrectly from the source to the system, e.g. reading a handwritten ``7'' as a ``1'', or typing \texttt{Manka} instead of \texttt{Mankaa}.
\end{definition}

\begin{definition}[Transposition Error]
A \cle{transposition error} occurs when two characters or digits are accidentally swapped during entry, e.g. typing \texttt{54} instead of \texttt{45}, or \texttt{Fotso} as \texttt{Ftoso}.
\end{definition}

Automated methods add their own error type: the \cle{misread}. A barcode scanner fails on a crumpled label; an OMR scanner misreads a sheet where a candidate shaded two bubbles or used a pen instead of a pencil; a fingerprint reader rejects a wet finger.

\begin{attention}[Automated Does Not Mean Error-Free]
A very common mistake is to write that automated capture ``eliminates errors''. It does not: it \emph{reduces} human typing errors but introduces misreads, damaged tags, poorly shaded forms and faulty sensors. This is precisely why the next stage --- \textbf{verification and validation} (Section 3) --- exists: captured data must be checked before it is trusted.
\end{attention}

\begin{exoresolu}[Choosing and Justifying a Capture Method]
A company runs 200 buses between Douala and Yaound\'e. Management wants to know, in real time, where every bus is, and wants drivers to stop filling in paper logbooks. Propose a capture method and justify it using the five choice criteria.
\tcblower
\textbf{Proposed solution:} GPS trackers (sensors with automatic data logging) on each bus, transmitting positions to the central system.
\begin{itemize}
\item \textbf{Volume}: 200 buses reporting every minute is far too much for manual entry --- automation is essential.
\item \textbf{Speed}: ``real time'' monitoring requires instant, continuous capture; only sensors deliver this.
\item \textbf{Accuracy}: GPS positions are machine-generated, removing transcription and transposition errors from handwritten logbooks.
\item \textbf{Cost}: trackers and a subscription cost money, but far less than employing clerks to type 200 logbooks daily.
\item \textbf{Environment}: buses move through rain, dust and night; sensors work in all these conditions, unlike paper forms or barcode scans.
\end{itemize}
\textbf{Residual risk:} devices can fail or lose signal, so logged data should still be validated (e.g. rejecting impossible positions) --- the subject of the next section.
\end{exoresolu}

\begin{savaistu}[Turnaround Documents (GCE Extra)]
A \cle{turnaround document} is a document printed by a computer, sent out, and later returned with additional data so it can be re-captured by machine. Classic example: an electricity or water bill printed with your customer number and amount as a barcode or OCR line; when you pay, the cashier scans the same bill and the payment is captured instantly and matched to your account. Examination answer sheets are also turnaround documents: the candidate's details are pre-printed, the candidate adds marks, and the sheet returns for OMR capture.
\end{savaistu}

\begin{exercice}[Revision Questions]
\begin{enumerate}
\item Define \emph{data capture} and distinguish manual from automated capture.
\item For each situation, name the most suitable capture method and justify: (a) marking a census questionnaire with tick boxes; (b) clearing bank cheques; (c) registering voters so nobody registers twice; (d) monitoring the temperature of a vaccine fridge.
\item A clerk types a mark of \texttt{84} as \texttt{48}. What type of error is this? How does it differ from a transcription error?
\item State one advantage and one limitation of RFID compared with barcodes.
\item Explain why captured data must still be verified and validated even when automated methods are used.
\end{enumerate}
\end{exercice}

\begin{corrige}[Exercise 1 --- hints]
\begin{enumerate}
\item Data capture = collecting raw data and entering it into a computer system; manual = human keys it in, automated = a device reads it directly.
\item (a) OMR; (b) MICR; (c) biometric capture (fingerprints); (d) temperature sensor with data logging --- justify with volume, speed, accuracy, cost, environment.
\item Transposition error (digits swapped); a transcription error is copying a wrong value without swapping.
\item Advantage: no line of sight, works through dirt, many tags read at once. Limitation: higher cost per tag, limited range, possible radio interference.
\item Automated capture still produces misreads and device faults, so checks (verification, validation) remain necessary before processing.
\end{enumerate}
\end{corrige}

\begin{aretenir}[Key Points]
\begin{itemize}
\item Data capture is the input stage of the information system: raw data enters the computer, manually or automatically.
\item Manual methods (keyboard, mouse/touchscreen, voice) are flexible and cheap but slow and error-prone.
\item Automated methods: optical (OMR, OCR, barcode, QR), magnetic/chip (stripe, MICR, smart card), radio/biometric (RFID, fingerprints, facial recognition) and sensors (data logging).
\item Choose using five criteria: volume, speed, accuracy, cost, environment.
\item Typical capture errors: transcription, transposition, misreads --- hence the need for verification and validation.
\item \textbf{Exam tip:} learn at least one concrete example per capture method (OMR $\rightarrow$ GCE answer sheets, MICR $\rightarrow$ cheques, RFID $\rightarrow$ tolls, biometrics $\rightarrow$ national ID). Examiners reward precise, correctly matched examples.
\end{itemize}
\end{aretenir}

\coursec{Data Verification and Validation}

In the previous section we studied \cle{data capture}: how raw data is collected from source documents, sensors, keyboards or scanners and fed into the computer. But capturing data is not enough. A single slip of the finger --- typing $15$ instead of $5$ for a child's medicine dosage, or entering $02$ instead of $20$ for a GCE mark --- can turn a perfectly designed information system into a dangerous machine. This section presents the two complementary defences that every information system uses before data is stored: \cle{verification} and \cle{validation}.

\courssub{Why Errors Matter: Garbage In, Garbage Out}

\textbf{Observation.} Imagine the results office of a school in Yaoundé preparing GCE transcripts. A clerk captures the mark of candidate $0451$ as $09/20$ instead of $19/20$. The software is excellent: it computes averages, ranks candidates and prints beautiful transcripts. But every calculation built on the wrong mark is wrong. The candidate may lose a distinction, a scholarship, or a place at university --- not because the computer failed, but because the \emph{input} was wrong.

\begin{definition}[GIGO --- Garbage In, Garbage Out]
\cle{GIGO} (\emph{Garbage In, Garbage Out}) is the principle that the quality of the output of an information system can never exceed the quality of its input. If incorrect, incomplete or unreasonable data enters the system, then all processing, reports and decisions based on that data will be unreliable, no matter how powerful the software is.
\end{definition}

\begin{exemplebox}[Consequences of bad data for decisions]
\begin{itemize}
\item \textbf{Hospital:} a nurse keys in a dosage of \SI{25}{mg} instead of \SI{2.5}{mg}. The pharmacy system, trusting the value, validates the order: the patient receives ten times the correct dose.
\item \textbf{Examinations:} a mark of $92/20$ is accepted by a badly designed form. The class average is inflated and the ranking of every student is distorted.
\item \textbf{Bank:} one wrong digit in an account number sends a worker's salary of $150\,000$ FCFA to a stranger's account.
\item \textbf{Stock management:} a captured quantity of $1000$ instead of $100$ bags of cement makes the system believe the warehouse is full; no new order is placed and the construction site stops.
\end{itemize}
\end{exemplebox}

The lesson is clear: data must be \emph{checked at the point of entry}, before it is stored. Two families of checks exist, and the GCE examiner expects you to distinguish them perfectly.

\courssub{Verification: Does the Entered Data Match the Source?}

\textbf{Intuition.} You copy a phone number from an admission letter into a form. How do you make sure you copied it \emph{exactly}? You read it again, character by character, comparing what you typed with the letter. That is verification: a comparison between the captured data and the original source document.

\begin{definition}[Verification]
\cle{Verification} is the process of checking that the data entered into a computer system is \textbf{identical to the data on the original source document}. Verification detects \emph{transcription errors} (mistakes made while copying or keying). It is usually performed by a human being, sometimes assisted by the computer.
\end{definition}

\textbf{The three standard verification methods}

\begin{enumerate}
\item \textbf{Double entry (keying the data twice).} The same data is typed twice, either by the same operator or by two different operators, often into two different files. The computer then \emph{compares the two versions} and reports any difference. If the two copies disagree, at least one of them is wrong and the operator returns to the source document. This is how examination marks are often captured: two clerks key the same scripts independently.
\item \textbf{Visual check / proofreading.} The operator (or a second person) carefully reads what is on the screen or on a printed list and compares it, field by field, with the source document. Simple and cheap, but tiring and unreliable for large volumes.
\item \textbf{Echo-back on screen.} As data is typed (or read by a scanner), it is immediately redisplayed so the operator can confirm it before pressing \emph{Enter}. When you type a password and the system asks you to retype it, or when a website shows ``You entered: 6\,77\,45\,23\,10 --- confirm?'', that is echo-back verification.
\end{enumerate}

\begin{attention}[Verification does not check meaning]
Verification only guarantees that the computer holds \emph{what was written on the source}. If the source document itself is wrong (a teacher wrote $09$ instead of $19$ on the mark sheet), verification will faithfully reproduce the error. Verification answers the question: \emph{``Was it copied correctly?''} --- never \emph{``Is it sensible?''}
\end{attention}

\courssub{Validation: Is the Data Reasonable and Acceptable?}

\textbf{Intuition.} A student registration form asks for a mark out of $20$. You type $25$. The form refuses instantly: ``Mark must be between 0 and 20.'' Nobody compared your typing with any document --- the \emph{software itself} rejected the value because it breaks a rule. That is validation.

\begin{definition}[Validation]
\cle{Validation} is the \textbf{automatic checking performed by software} to ensure that entered data is \emph{reasonable}, \emph{complete} and \emph{within acceptable limits}, according to a set of predefined rules (validation rules). Data that fails a validation check is rejected and must be re-entered.
\end{definition}

\begin{attention}[Validation does NOT prove correctness]
This is the most frequent exam trap. Validation only proves that data is \textbf{plausible}, not that it is \textbf{true}. A mark of $18/20$ passes the range check $0$--$20$ perfectly --- but if the real mark was $8/20$, the validated value is still wrong. Validation answers: \emph{``Could this value be acceptable?''} Only verification (comparison with the source) can answer: \emph{``Is this the right value?''}
\end{attention}

\courssub{The Six Standard Validation Checks}

\begin{methode}[The six standard validation checks and when to apply each]
\begin{enumerate}
\item \textbf{Presence check} --- the field must not be left empty. \emph{Use for:} compulsory fields such as candidate name, date of birth. \emph{Rejects:} a blank ``Name'' field.
\item \textbf{Data type check} --- the value must be of the expected type (numeric, alphabetic, date\ldots). \emph{Use for:} quantities, marks, phone numbers. \emph{Rejects:} \texttt{ABC} entered in the ``Mark'' field.
\item \textbf{Format check} --- the value must follow a required pattern. \emph{Use for:} dates (\texttt{DD/MM/YYYY}), matricule numbers (\texttt{AB-1234}), e-mail addresses. \emph{Rejects:} \texttt{2024-31-01} in a \texttt{DD/MM/YYYY} field.
\item \textbf{Length check} --- the value must have an exact or maximum number of characters. \emph{Use for:} phone numbers (9 digits in Cameroon), IDs, passwords. \emph{Rejects:} a 7-digit phone number.
\item \textbf{Range check} --- a numeric value (or date) must lie between a minimum and a maximum. \emph{Use for:} marks ($0$--$20$), ages, months ($1$--$12$). \emph{Rejects:} a mark of $23/20$ or a month of $13$.
\item \textbf{Limit check} --- a special case of the range check with only \emph{one} boundary (a minimum \emph{or} a maximum). \emph{Use for:} a quantity ordered that must not exceed stock; an age that must be at least $15$. \emph{Rejects:} a withdrawal greater than the account balance.
\end{enumerate}
A seventh, more advanced check is often required at GCE level:
\begin{enumerate}
\setcounter{enumi}{6}
\item \textbf{Consistency (cross-field) check} --- two or more fields must be logically compatible. \emph{Use for:} date of birth versus class level (a $5$-year-old cannot be in Upper Sixth); end date must be after start date; ``sex = M'' versus ``title = Mrs''. \emph{Rejects:} a student born in 2015 registered in Terminale.
\end{enumerate}
\end{methode}

\courssub{Check Digits}

\textbf{Intuition.} Long numbers --- bank account numbers, barcode numbers, student IDs --- are extremely easy to mistype, and a simple range check cannot help: any 5-digit number ``looks valid''. The clever solution is to make the number \emph{self-checking}: the last digit is not free, it is \emph{calculated} from the others. Change any digit, and the calculation no longer matches.

\begin{definition}[Check digit]
A \cle{check digit} is an extra digit appended to a number, computed from the other digits by a fixed arithmetic algorithm (very often a weighted sum followed by a modulo operation). When the number is entered, the software recomputes the check digit and compares it with the one typed; any mismatch reveals a transcription error. Check digits are used in bank account numbers, ISBN book codes, barcodes (EAN-13) and national ID numbers.
\end{definition}

\begin{methode}[Computing and checking a check digit (modulo-11 style)]
To compute the check digit of a 4-digit code $d_1 d_2 d_3 d_4$ (digits written left to right):
\begin{enumerate}
\item Assign weights $5, 4, 3, 2$ to the digits from left to right (i.e. $d_1\times 5$, $d_2\times 4$, $d_3\times 3$, $d_4\times 2$).
\item Add the products to obtain the weighted sum $S$.
\item Compute the remainder $r = S \bmod 11$.
\item The check digit is $c = 11 - r$ (if $c = 11$, use $0$; if $c = 10$, use the letter \texttt{X}).
\item \textbf{To check} a full code: recompute $c$ from the first digits and compare it with the last digit typed. Equality $\Rightarrow$ accept; difference $\Rightarrow$ reject and re-enter.
\end{enumerate}
\end{methode}

\begin{exemplebox}[Worked example: student ID 3842]
Weights $5,4,3,2$ applied to $3,8,4,2$:
\begin{align*}
S &= (3\times 5) + (8\times 4) + (4\times 3) + (2\times 2) = 15 + 32 + 12 + 4 = 63\\
r &= 63 \bmod 11 = 8 \qquad (\text{since } 63 = 5\times 11 + 8)\\
c &= 11 - 8 = 3
\end{align*}
The complete student ID is therefore \textbf{38423}. If a clerk types \texttt{38433} (swapping the last two digits of the base number), the system recomputes the check digit for \texttt{3843}:
\begin{align*}
S' &= (3\times 5) + (8\times 4) + (4\times 3) + (3\times 2) = 15 + 32 + 12 + 6 = 65\\
r' &= 65 \bmod 11 = 10 \quad (65 = 5\times 11 + 10)\\
c' &= 11 - 10 = 1
\end{align*}
The typed check digit is $3$, but the computed check digit is $1$; mismatch $\Rightarrow$ entry \emph{rejected}.
\end{exemplebox}

\begin{popfigure}
\begin{center}
\begin{tikzpicture}[font=\small]
% digits boxes
\foreach \x/\d/\w in {0/3/5, 1.6/8/4, 3.2/4/3, 4.8/2/2}{
  \draw[fill=popblueL, draw=popblue, thick, rounded corners=2pt] (\x,0) rectangle (\x+1.1,1);
  \node at (\x+0.55,0.5) {\Large\textbf{\d}};
  \node[below, popmuted] at (\x+0.55,0) {\footnotesize $\times \w$};
}
\node[popmuted] at (2.75,-0.85) {\footnotesize weighted sum $S = 15+32+12+4 = 63$};
\draw[->, thick, poporange] (2.75,-1.1) -- (2.75,-1.7);
\node[draw=poporange, thick, rounded corners=3pt, fill=white, align=center] at (2.75,-2.3) {$63 \bmod 11 = 8$};
\draw[->, thick, poporange] (2.75,-2.7) -- (2.75,-3.3);
\node[draw=popgreen, thick, rounded corners=3pt, fill=popgreenL, align=center] at (2.75,-3.9) {check digit $= 11-8 = \mathbf{3}$};
\draw[->, thick, popgreen] (4.1,-3.9) -- (5.6,-3.9);
\node[draw=popdark, thick, rounded corners=3pt, align=center] at (7.2,-3.9) {full ID: \textbf{3842}\textcolor{poporange}{\textbf{3}}};
\end{tikzpicture}
\end{center}
\legende{Computing the modulo-11 check digit for the student ID 3842.}
\end{popfigure}

\courssub{Parity Bits (Exam-Useful Extra)}

When data \emph{travels} (network, cable, memory), errors are not typing errors but flipped bits: a \texttt{0} becomes a \texttt{1}. The simplest protection is the \cle{parity bit}: an extra bit added to each byte so that the total number of \texttt{1}s is even (\emph{even parity}) or odd (\emph{odd parity}).

\begin{exemplebox}[Even parity]
The byte \texttt{1011001} contains four \texttt{1}s (already even), so the parity bit is \texttt{0}: we transmit \texttt{10110010}. If one bit flips during transmission, the received byte contains an \emph{odd} number of \texttt{1}s and the receiver detects the error and asks for retransmission. Note the limit: if \emph{two} bits flip, the parity looks even again and the error escapes detection --- parity detects single-bit errors only.
\end{exemplebox}

\courssub{Putting It All Together: The Data Entry Pipeline}

\begin{popfigure}
\begin{center}
\begin{tikzpicture}[font=\small, node distance=0.55cm,
  box/.style={draw=popblue, thick, rounded corners=3pt, fill=popblueL, align=center, minimum width=3.1cm, minimum height=0.85cm},
  gate/.style={draw=poporange, thick, rounded corners=3pt, fill=white, align=center, minimum width=3.1cm, minimum height=0.85cm},
  store/.style={draw=popgreen, thick, rounded corners=3pt, fill=popgreenL, align=center, minimum width=3.1cm, minimum height=0.85cm}]
\node[box, align=center] (src){Source document\\ (mark sheet, form)};
\node[box, below=of src, align=center] (key){Data capture\\ (keying / scanning)};
\node[gate, below=of key, align=center] (ver){VERIFICATION\\ double entry, echo-back};
\node[gate, below=of ver, align=center] (val){VALIDATION\\ presence, type, range\ldots};
\node[store, below=of val, align=center] (db){Database\\ (stored record)};
\draw[->, thick] (src) -- (key);
\draw[->, thick] (key) -- (ver);
\draw[->, thick] (ver) -- node[right, popmuted]{\footnotesize matches source? yes} (val);
\draw[->, thick] (val) -- node[right, popmuted]{\footnotesize passes rules? yes} (db);
\draw[->, thick, poppink] (ver.west) -- ++(-1.6,0) node[left, align=center, poppink]{\footnotesize mismatch:\\ re-enter} |- (key.west);
\draw[->, thick, poppink] (val.east) -- ++(1.6,0) node[right, align=center, poppink]{\footnotesize rejected:\\ error message} |- (key.east);
\end{tikzpicture}
\end{center}
\legende{Data entry pipeline: verification and validation act as two successive gates before storage.}
\end{popfigure}

\courssub{Verification versus Validation: The Comparison}

\begin{popfigure}
\begin{center}
\begin{tabularx}{\linewidth}{|l|X|X|}
\hline
\rowcolor{popblueL} \textbf{Criterion} & \textbf{Verification} & \textbf{Validation} \\
\hline
Question answered & Was the data copied exactly from the source? & Is the data reasonable, complete and within limits? \\
\hline
Performed by & Mainly a human (proofreading, double entry), assisted by the computer & Automatically by the software \\
\hline
When & During / just after data entry & At the moment of entry, before storage \\
\hline
Detects & Transcription (copying) errors & Impossible, incomplete or inconsistent values \\
\hline
Typical methods & Double entry, visual check, echo-back & Presence, type, format, length, range, limit, consistency, check digit \\
\hline
Example & Clerk retypes the mark $19$ twice; both copies agree & System rejects a mark of $25/20$ (range check) \\
\hline
Main weakness & Cannot detect an error present on the source document & Cannot prove the value is the true one \\
\hline
\end{tabularx}
\end{center}
\legende{Verification versus validation: two complementary defences.}
\end{popfigure}

\begin{attention}[Do not confuse the two]
A classic GCE trap: ``\emph{The system checks that the mark typed is the same as on the script --- is this validation?}'' \textbf{No}: comparing with the source is \emph{verification}. Conversely, ``\emph{the software rejects a date of birth of 30/02/2015}'' is \emph{validation} (format/range check), not verification. Remember: \textbf{Verification = matches the source; Validation = obeys the rules.}
\end{attention}

\courssub{Designing Validation Rules for a Data-Entry Form}

\begin{methode}[How to answer a ``design the validation'' exam question]
For each field of the given form:
\begin{enumerate}
\item Name the \textbf{appropriate check} (presence, type, format, length, range, limit, consistency, check digit).
\item \textbf{Justify} it in one sentence (what error it prevents).
\item Give a concrete \textbf{rejected example value} --- this is what earns full marks.
\end{enumerate}
\end{methode}

\begin{exoresolu}[Student registration form]
A school uses a registration form with the fields: \emph{Name}, \emph{Date of birth} (DD/MM/YYYY), \emph{Mark at entrance test} (/20), \emph{Phone number} (9 digits), \emph{Student ID} (4 digits + check digit). Specify a suitable validation check for each field, with a rejected example.
\tcblower
\begin{itemize}
\item \textbf{Name:} \emph{presence check} --- the field is compulsory; a blank name is rejected. (Also a type check: \texttt{John3} is rejected.)
\item \textbf{Date of birth:} \emph{format check} DD/MM/YYYY plus a \emph{range check} on the day ($1$--$31$) and month ($1$--$12$); \texttt{31/13/2008} is rejected. A \emph{consistency check} with the class applied for rejects a 4-year-old in Lower Sixth.
\item \textbf{Mark /20:} \emph{data type check} (numeric) and \emph{range check} $0 \le m \le 20$; \texttt{23} and \texttt{AB} are rejected.
\item \textbf{Phone number:} \emph{length check} (exactly 9 digits) and type check (digits only); \texttt{6774512} (7 digits) is rejected.
\item \textbf{Student ID:} \emph{check digit} recomputed by the modulo-11 algorithm; \texttt{38424} is rejected because the correct check digit of $3842$ is $3$.
\end{itemize}
Note that \emph{verification} (double entry or echo-back) is still needed for the mark: validation accepts $18$ even if the true mark is $8$.
\end{exoresolu}

\begin{exercice}[Design your own checks]
A hospital system captures: \emph{patient name}, \emph{weight in kg}, \emph{date of visit} (DD/MM/YYYY), \emph{ward number} (1 to 12), \emph{account number} (8 digits with check digit). For each field, name the validation check(s) you would apply, justify your choice, and give one value that must be rejected.
\end{exercice}

\begin{corrige}[Exercise 1]
\begin{itemize}
\item \textbf{Patient name:} presence check (compulsory field); rejects a blank entry.
\item \textbf{Weight:} data type check (numeric) and range check, e.g.\ $0.5$--$300$ kg; rejects \texttt{-5} and \texttt{ABC}.
\item \textbf{Date of visit:} format check DD/MM/YYYY and range checks on day and month; rejects \texttt{32/01/2025}. A consistency check can reject a visit date before the patient's date of birth.
\item \textbf{Ward number:} range check $1$--$12$; rejects \texttt{0} and \texttt{15}.
\item \textbf{Account number:} check digit verification by recomputation; rejects any number whose last digit does not match the computed one.
\end{itemize}
\end{corrige}

\begin{aretenir}[Key points]
\begin{itemize}
\item \textbf{GIGO}: wrong input $\Rightarrow$ wrong output and wrong decisions, whatever the software quality.
\item \textbf{Verification} = the entered data \emph{matches the source} (double entry, visual check, echo-back).
\item \textbf{Validation} = the software checks the data is \emph{plausible} (presence, type, format, length, range, limit, consistency, check digit). It never proves the data is correct.
\item \textbf{Check digit}: extra digit computed from the others (weighted sum, modulo); detects mistyped account numbers, IDs and barcodes.
\item \textbf{Parity bit}: detects single-bit transmission errors (even/odd parity).
\item In the exam: for each field, name the check \emph{and} give a rejected example value.
\end{itemize}
\end{aretenir}

\coursec{Organisational Levels and Information Systems for Decision-Making}

In the previous section we saw that data must be \cle{verified} and \cle{validated} before it can be trusted. But trustworthy data alone does not run an organisation: it must reach the \emph{right person}, at the \emph{right level of detail}, at the \emph{right time}. A cashier at a microfinance counter in Bamenda does not need the same information as the board of directors in Yaoundé. This section explains how organisations are structured in levels, what kind of decision each level takes, and which type of information system serves each level.

\courssub{The Three Levels of an Organisation}

\begin{experience}[Observation: who needs what?]
Consider a Cameroonian microfinance institution (an \emph{MC2}-type cooperative) with several branches.
\begin{itemize}
\item At the counter, the \textbf{cashier} asks: \emph{``What is the balance of account 04512? Has this member repaid today's instalment?''}
\item At the branch office, the \textbf{branch manager} asks: \emph{``How many loans did we grant this month? Which tellers are behind target?''}
\item At head office, the \textbf{board of directors} asks: \emph{``Should we open three new branches in the North-West next year? Should we launch a mobile-money product?''}
\end{itemize}
Same institution, same underlying data --- but three completely different questions, horizons and levels of detail. This simple observation is the key to the whole section: \textbf{the level of the user determines the kind of information system needed}.
\end{experience}

\begin{definition}[Organisational levels]
Most organisations are analysed as a pyramid of three \cle{levels}:
\begin{itemize}
\item \textbf{Operational level} (bottom): staff who carry out day-to-day activities --- cashiers, tellers, shop assistants, the school bursar recording fees, the secretary registering students.
\item \textbf{Tactical level} (middle management): staff who plan and control in the short to medium term --- branch managers, regional delegates of education, heads of department.
\item \textbf{Strategic level} (top management): those who define long-term direction --- the CEO, the board, a minister, the proprietor of a school group.
\end{itemize}
The pyramid shape reflects two facts: there are \emph{many} operational staff and \emph{few} executives, and the \emph{scope} of each decision widens as one moves up (a cashier's error affects one account; a board's error affects the whole institution).
\end{definition}

\begin{definition}[Types of decision]
A \cle{decision} is a choice among possible actions. Decisions are classified by how well-defined their procedure is:
\begin{itemize}
\item \textbf{Structured decision}: the procedure is known and repetitive; it can be automated by a rule (e.g. \emph{if stock $<$ reorder level, order more}).
\item \textbf{Semi-structured decision}: part of the problem can be modelled, but human judgement is still required (e.g. setting next term's school fees).
\item \textbf{Unstructured decision}: novel, complex, no agreed procedure; relies on experience, intuition and external information (e.g. deciding to merge with another microfinance institution).
\end{itemize}
\end{definition}

\begin{propriete}[Decision type versus organisational level]
The lower the level, the more \cle{structured} the decisions; the higher the level, the more \cle{unstructured} they are:
\begin{itemize}
\item Operational level $\rightarrow$ mostly \textbf{structured} decisions, taken very frequently (many times per day), with a very short time horizon (minutes to days).
\item Tactical level $\rightarrow$ mostly \textbf{semi-structured} decisions, taken periodically (weekly, monthly, termly), horizon of weeks to a year.
\item Strategic level $\rightarrow$ mostly \textbf{unstructured} decisions, taken rarely, horizon of several years, with the greatest consequences.
\end{itemize}
(Not examinable as a proof, but you must be able to \emph{justify} this correspondence with examples in an essay.)
\end{propriete}

\begin{center}
\begin{tabularx}{\linewidth}{|l|X|X|X|}
\hline
\rowcolor{popblueL} \textbf{Feature} & \textbf{Operational} & \textbf{Tactical} & \textbf{Strategic} \\
\hline
Typical user & Cashier, bursar, teacher & Branch manager, regional delegate & CEO, board, minister \\
\hline
Decision type & Structured & Semi-structured & Unstructured \\
\hline
Time horizon & Hours -- days & Weeks -- 1 year & 1 -- 10 years \\
\hline
Frequency & Continuous, daily & Periodic & Rare, ad hoc \\
\hline
Information needed & Detailed, internal, real-time & Summarised, periodic, internal & Aggregated, internal + external, forecasts \\
\hline
Example question & ``Has member X paid today?'' & ``Which branches are below target this quarter?'' & ``Do we expand to a new region?'' \\
\hline
\end{tabularx}
\end{center}

\begin{popfigure}
\begin{center}
\begin{tikzpicture}[scale=1.0, font=\small]
% Pyramid layers
\fill[popblue!80] (0,4.6) -- (1.6,2.9) -- (-1.6,2.9) -- cycle;
\fill[popgreen!75] (-1.6,2.9) -- (1.6,2.9) -- (3.2,1.1) -- (-3.2,1.1) -- cycle;
\fill[poporange!75] (-3.2,1.1) -- (3.2,1.1) -- (4.8,-0.7) -- (-4.8,-0.7) -- cycle;
% Labels inside
\node[white, align=center] at (0,3.55) {\textbf{STRATEGIC}\\ \scriptsize board, CEO, minister};
\node[white, align=center] at (0,2.0) {\textbf{TACTICAL}\ \ \scriptsize branch manager, delegate};
\node[white, align=center] at (0,0.2) {\textbf{OPERATIONAL}\ \ \scriptsize cashier, bursar, teller};
% Right annotations: decisions
\node[right, align=left, text width=4.6cm] at (5.2,3.55) {\scriptsize Unstructured decisions\\ \scriptsize long term, rare\\ \scriptsize \textbf{ESS} (+ DSS)};
\node[right, align=left, text width=4.6cm] at (5.2,2.0) {\scriptsize Semi-structured decisions\\ \scriptsize short/medium term\\ \scriptsize \textbf{MIS} (+ DSS)};
\node[right, align=left, text width=4.6cm] at (5.2,0.2) {\scriptsize Structured decisions\\ \scriptsize daily, repetitive\\ \scriptsize \textbf{TPS}};
\draw[->, popmuted] (4.9,3.55) -- (1.2,3.55);
\draw[->, popmuted] (4.9,2.0) -- (2.6,2.0);
\draw[->, popmuted] (4.9,0.2) -- (4.2,0.2);
% Left arrow: volume vs scope
\draw[<->, thick, popdark] (-5.6,-0.7) -- (-5.6,4.6) node[midway, left, rotate=90, align=center, text width=2.6cm] {\scriptsize fewer people, wider scope};
\end{tikzpicture}
\end{center}
\legende{The organisational pyramid: levels, decision types and the information systems that serve them.}
\end{popfigure}

\courssub{The Four Types of Information System}

Each level is served by a category of information system. Together they form a chain: data enters at the bottom, through the daily transactions, and flows upward in increasingly summarised form. We now define the four classical types, from the bottom of the pyramid to the top.

\begin{definition}[Transaction Processing System (TPS)]
A \cle{Transaction Processing System} is the information system that records the \textbf{daily, routine transactions} of an organisation: sales, deposits, withdrawals, fee payments, enrolments, stock movements. It operates at the \textbf{operational level}, processes data in real time or in batches, and produces documents such as receipts, invoices, pay slips and daily lists. The TPS is the \textbf{backbone} of all other systems: it feeds the organisation's database with raw, validated data.
\end{definition}

\begin{exemplebox}[TPS in our microfinance institution]
Every deposit or loan repayment at the counter is keyed into the TPS, which checks the account number (validation!), updates the member's balance and prints a receipt. At the end of the day it produces the teller's transaction list used for cash reconciliation. In a school, the equivalent is the system recording each fee payment and printing the receipt given to the parent.
\end{exemplebox}

\begin{definition}[Management Information System (MIS)]
A \cle{Management Information System} takes the data recorded by the TPS and produces \textbf{periodic summary reports} for \textbf{middle managers}: weekly, monthly or termly totals, comparisons against targets, and \emph{exception reports} highlighting only abnormal situations (e.g. branches whose loan recovery rate fell below a threshold). MIS serve the \textbf{tactical level} and support semi-structured, short-term decisions. An MIS \textbf{summarises the past}; it does not predict the future.
\end{definition}

\begin{exemplebox}[MIS report]
At the end of each month, the MIS of our microfinance institution produces a report showing, \emph{per branch}: total deposits, loans granted, repayment rate, and comparison with the monthly target. The branch manager of Bafoussam sees at a glance that her recovery rate dropped and investigates --- she does not read the thousands of individual transactions.
\end{exemplebox}

\begin{definition}[Decision Support System (DSS)]
A \cle{Decision Support System} helps managers analyse \textbf{semi-structured} problems using \textbf{models, simulations and ``what-if'' analysis}. The user changes assumptions (interest rate, price, enrolment growth) and the DSS recomputes the projected outcomes. A spreadsheet model is the simplest and most common DSS. The DSS \emph{supports} the decision; the human being still decides.
\end{definition}

\begin{exemplebox}[DSS: what-if pricing]
The credit manager builds a spreadsheet model: \emph{``If we raise the loan interest rate from 2\,\% to 2.5\,\% per month, and 5\,\% of borrowers leave, what happens to yearly revenue?''} She tests several scenarios before recommending a rate to the board. The DSS does not decide --- it \emph{supports} the decision.
\end{exemplebox}

\begin{definition}[Executive Support System (ESS)]
An \cle{Executive Support System} (also called Executive Information System) serves the \textbf{strategic level}. It presents \textbf{highly aggregated key indicators} through easy-to-read \textbf{dashboards} (charts, gauges, trends), combining \textbf{internal} summaries with \textbf{external} data (economic trends, competitor activity, population statistics) to support long-term, unstructured decisions.
\end{definition}

\begin{exemplebox}[ESS: the minister's dashboard]
A national dashboard for the Minister of Secondary Education might show, per region: enrolment trends over five years, examination success rates, teacher--pupil ratios, and demographic forecasts --- a few screens of charts rather than millions of student records. Similarly, our microfinance board watches a dashboard of portfolio growth, default risk and regional economic indicators before deciding to open new branches.
\end{exemplebox}

\begin{popfigure}
\begin{center}
\begin{tikzpicture}[font=\small, node distance=0.55cm,
box/.style={draw=popblue, thick, rounded corners, fill=popblueL, align=center, text width=2.5cm, minimum height=1.15cm},
db/.style={cylinder, draw=popdark, thick, fill=popgreenL, shape border rotate=90, aspect=0.32, minimum width=1.7cm, minimum height=1.25cm, align=center},
lbl/.style={font=\scriptsize\itshape, text=popmuted}]
\node[box, align=center] (tps){\textbf{TPS}\\ \scriptsize daily transactions};
\node[db, right=1.0cm of tps] (db) {\textbf{Database}};
\node[box, right=1.0cm of db, align=center] (mis){\textbf{MIS}\\ \scriptsize periodic reports};
\node[box, right=1.0cm of mis, align=center] (dss){\textbf{DSS}\\ \scriptsize what-if models};
\node[box, above=0.7cm of dss, align=center] (ess){\textbf{ESS}\\ \scriptsize dashboards};
\draw[->, thick, poporange] (tps) -- node[lbl, above]{raw validated data} (db);
\draw[->, thick, poporange] (db) -- node[lbl, above]{summaries} (mis);
\draw[->, thick, poporange] (mis) -- node[lbl, above]{analyses} (dss);
\draw[->, thick, poporange] (dss) -- node[lbl, right]{key indicators} (ess);
\draw[->, thick, poppurple, dashed] ($(ess.north)+(0,0.7)$) node[above, lbl]{external data (economy, demographics)} -- (ess.north);
\node[lbl, below=0.35cm of db] {increasing summarisation $\longrightarrow$};
\end{tikzpicture}
\end{center}
\legende{Data flow: the TPS feeds the database; MIS, DSS and ESS consume increasingly summarised information.}
\end{popfigure}

\begin{popfigure}
\begin{center}
\begin{tikzpicture}[font=\scriptsize]
\node[draw=popdark, thick, fill=popblueL, align=center, text width=2.1cm] (h1) at (0,0) {\textbf{IS type}};
\node[draw=popdark, thick, fill=popblueL, align=center, text width=2.4cm] (h2) at (2.6,0) {\textbf{Main users}};
\node[draw=popdark, thick, fill=popblueL, align=center, text width=2.4cm] (h3) at (5.2,0) {\textbf{Decision type}};
\node[draw=popdark, thick, fill=popblueL, align=center, text width=2.2cm] (h4) at (7.8,0) {\textbf{Time horizon}};
\node[draw=popdark, thick, fill=popblueL, align=center, text width=3.0cm] (h5) at (10.6,0) {\textbf{Example}};
\foreach \y/\a/\b/\c/\d/\e in {
 {-1.0}/{TPS}/{cashiers, clerks}/{structured}/{hours--days}/{receipts, daily lists},
 {-2.0}/{MIS}/{middle managers}/{semi-structured}/{weeks--months}/{monthly sales per region},
 {-3.0}/{DSS}/{managers, analysts}/{semi-structured}/{months--years}/{loan-pricing spreadsheet},
 {-4.0}/{ESS}/{executives, board}/{unstructured}/{years}/{branch-expansion dashboard}}{
 \node[draw=popline, align=center, text width=2.1cm] at (0,\y) {\a};
 \node[draw=popline, align=center, text width=2.4cm] at (2.6,\y) {\b};
 \node[draw=popline, align=center, text width=2.4cm] at (5.2,\y) {\c};
 \node[draw=popline, align=center, text width=2.2cm] at (7.8,\y) {\d};
 \node[draw=popline, align=center, text width=3.0cm] at (10.6,\y) {\e};
}
\end{tikzpicture}
\end{center}
\legende{Matching each information system to its users, decision type, horizon and example.}
\end{popfigure}

\courssub{Case Study: Mapping a Microfinance Institution}

\begin{experience}[Case study: ``Confiance Microfinance'', Cameroon]
\textbf{Situation.} Confiance Microfinance has a head office in Yaoundé and 12 branches. Map its staff and systems.
\begin{itemize}
\item \textbf{Cashiers (operational, TPS):} record deposits, withdrawals and repayments; the TPS validates account numbers and prints receipts.
\item \textbf{Branch managers (tactical, MIS):} receive the monthly MIS report --- deposits, loans granted, recovery rate per teller --- and an exception report flagging accounts in arrears beyond 30 days.
\item \textbf{Credit committee (tactical/strategic, DSS):} uses a spreadsheet model to simulate the effect of changing interest rates or loan ceilings before proposing a policy.
\item \textbf{Board of directors (strategic, ESS):} consults a quarterly dashboard --- portfolio growth, default rate, liquidity, plus external data on the national economy --- to decide on opening branches in two new regions.
\end{itemize}
\textbf{The chain:} every receipt printed by the TPS ends up, summarised three times, as one coloured indicator on the board's dashboard.
\end{experience}

\begin{methode}[Identifying the level and recommending an IS]
Given an exam scenario:
\begin{enumerate}
\item \textbf{Identify the actor and the question} they must answer (who? what decision?).
\item \textbf{Determine the level}: daily routine $\rightarrow$ operational; periodic control and short-term planning $\rightarrow$ tactical; long-term direction $\rightarrow$ strategic.
\item \textbf{Classify the decision}: is there a fixed rule (structured), a partial model (semi-structured), or no procedure at all (unstructured)?
\item \textbf{Match the IS}: operational/structured $\rightarrow$ TPS; tactical/reporting $\rightarrow$ MIS; analysis and what-if $\rightarrow$ DSS; strategic dashboards $\rightarrow$ ESS.
\item \textbf{Justify} your choice in one sentence linking the user's need (detail, period, horizon) to the system's output.
\end{enumerate}
\end{methode}

\begin{exoresolu}[Choosing the right system]
\textbf{Statement.} For each situation, name the organisational level and the most appropriate information system, with a brief justification.
\begin{enumerate}
\item A bursar records each student's fee payment and issues a receipt.
\item A regional delegate compares this term's enrolment with targets across all schools of the region.
\item The board of a bank studies whether to launch mobile banking within five years.
\end{enumerate}
\tcblower
\textbf{Solution.}
\begin{enumerate}
\item \textbf{Operational level --- TPS.} Recording a payment is a daily, structured transaction; the system must update balances immediately and produce a receipt.
\item \textbf{Tactical level --- MIS.} The delegate needs a periodic, summarised comparison against targets drawn from data recorded by schools' TPS --- exactly the role of an MIS report (possibly an exception report listing schools far below target).
\item \textbf{Strategic level --- ESS (supported by DSS).} This is a rare, unstructured, long-term decision requiring aggregated internal indicators plus external data (mobile penetration, competitors). A DSS may first simulate adoption scenarios; the ESS dashboard presents the key indicators to the board.
\end{enumerate}
\end{exoresolu}

\begin{attention}[Common mistakes]
\begin{itemize}
\item Do \textbf{not} attribute \emph{strategic} decisions to a TPS: a TPS only records transactions; it never decides policy.
\item Do \textbf{not} say an ESS handles \emph{daily} transactions: executives never scroll through individual receipts.
\item An MIS does not ``predict the future'' --- it summarises the past; forecasting and what-if analysis belong to the \textbf{DSS}.
\item Remember the direction of flow: TPS $\rightarrow$ database $\rightarrow$ MIS $\rightarrow$ DSS/ESS, with \emph{increasing summarisation}, never the reverse.
\item Do not confuse the \emph{decision type} with the \emph{importance} of the decision: a structured decision (reordering stock) can still be vital; ``unstructured'' only means there is no fixed procedure.
\end{itemize}
\end{attention}

\begin{aretenir}[Exam tip: memorise one complete chain]
Learn \textbf{one full example chain} by heart --- for instance our microfinance: \emph{cashier's receipt (TPS) $\rightarrow$ monthly branch report (MIS) $\rightarrow$ interest-rate simulation (DSS) $\rightarrow$ board expansion dashboard (ESS)}. In an essay question, reproduce this structure and simply substitute the organisation named in the paper (school, hospital, bank). A coherent, fully-worked chain earns far more marks than four disconnected definitions.
\end{aretenir}

\begin{savaistu}[GCE extra: expert systems and AI-assisted decision support]
Beyond the four classical systems, modern organisations increasingly use \cle{expert systems} --- programs that encode the knowledge of human experts as rules (e.g. a loan-eligibility adviser) --- and \cle{AI-assisted decision support}, where machine-learning models detect patterns such as credit-default risk or likely examination failure, and recommend actions to managers. These tools extend the DSS/ESS family, but the final responsibility for unstructured, strategic decisions remains with human leaders.
\end{savaistu}

\begin{exercice}[Practice]
A chain of pharmacies in Douala uses software that (a) records each sale and updates stock, (b) sends the owner a weekly report of the ten best-selling products, (c) simulates the profit impact of opening a new shop in Bonabéri, and (d) shows the owner a dashboard comparing her chain with national market trends. Identify the IS type and organisational level for each of (a)--(d), justifying each answer in one sentence.
\end{exercice}

\begin{corrige}[Practice]
\begin{itemize}
\item (a) \textbf{TPS, operational level}: a structured, real-time daily transaction (sale recorded, stock updated).
\item (b) \textbf{MIS, tactical level}: a periodic summary report derived from TPS data for short-term control.
\item (c) \textbf{DSS, tactical/strategic}: a what-if simulation supporting a semi-structured decision.
\item (d) \textbf{ESS, strategic level}: aggregated key indicators combined with external market data for long-term decisions.
\end{itemize}
\end{corrige}

\newpage

\coursec{Integration activity}

\courssub{Aims of the activity}

This integration activity consolidates the entire chapter \emph{Supporting decision-making in organisations}. It places you in the role of a junior systems analyst tasked with designing a coherent Information System (IS) for a realistic Cameroonian microfinance institution. By the end of this activity, you should be able to:
\begin{itemize}
    \item Identify and distinguish the three organisational levels (operational, tactical, strategic) and the specific decisions taken at each level.
    \item Select and justify appropriate \cle{data capture} methods for different source documents, considering efficiency, accuracy, and the Cameroonian context.
    \item Design comprehensive \cle{validation} rules (range, format, presence, check digit, existence) and choose a suitable \cle{verification} method (double entry, proofreading) for key data entry forms.
    \item Propose a layered IS architecture linking \cle{Transaction Processing Systems (TPS)}, \cle{Management Information Systems (MIS)}, and \cle{Decision Support Systems (DSS)}/\cle{Executive Support Systems (ESS)}.
    \item Represent the flow of data through the organisation using a \cle{Data Flow Diagram (DFD)} (Context and Level 1).
    \item Communicate and defend a technical design in a structured oral presentation.
\end{itemize}

\courssub{Situation: ``Sawa Savings'' Microfinance Cooperative}

\textbf{Sawa Savings} is a growing microfinance cooperative headquartered in \textbf{Buea}, Southwest Region, with three branches: \textbf{Buea Town}, \textbf{Molyko}, and \textbf{Tiko}. It serves over 12,000 members (small traders, farmers, commercial bike riders, and civil servants).

\medskip
\noindent\textbf{Current Manual Operations (The ``As-Is'' Situation):}
\begin{enumerate}
    \item \textbf{Member Registration:} A prospective member fills a paper \emph{Registration Form} (Name, Date of Birth, National ID Number, Phone Number, Next of Kin, Passport Photo). A cashier manually copies details into a large \emph{Member Ledger} (one ledger per branch). Membership cards are handwritten cardboard cards with a sequential number.
    \item \textbf{Daily Transactions (Deposits, Withdrawals, Loan Repayments):} Members queue at the counter. The cashier reads the handwritten card, finds the member's page in the ledger, writes the transaction details (Date, Type, Amount, Balance) by hand, and gives a handwritten receipt from a duplicate book. No biometric or PIN verification exists; identity relies on the cashier's memory or the photo on the cardboard card.
    \item \textbf{Branch Reporting:} Every Friday, the Branch Manager manually totals the ledger pages to produce a \emph{Weekly Summary Report} (Total Deposits, Total Withdrawals, New Loans Issued, Loan Repayments, New Members). This is sent via WhatsApp photo or physical notebook to the Head Office.
    \item \textbf{Head Office \& Board:} The General Manager (GM) compiles the three branch reports into a \emph{Monthly Consolidated Report} using a personal Excel spreadsheet. The Board of Directors meets \textbf{quarterly}. They receive printed reports that are already 3--4 weeks old. Strategic decisions (e.g., opening a new branch in Limbe, adjusting interest rates, launching a mobile money partnership) are based on outdated, aggregated figures. There is no drill-down capability to see which branch or product is underperforming.
\end{enumerate}

\medskip
\noindent\textbf{Key Pain Points Identified by Management:}
\begin{itemize}
    \item \textbf{Errors:} Frequent arithmetic mistakes in ledgers; transposition errors in ID numbers (e.g., writing \texttt{112345678} instead of \texttt{112345768}); duplicate membership numbers across branches.
    \item \textbf{Fraud Risk:} Handwritten receipts are easy to forge; no audit trail for cashier actions; ``ghost members'' were discovered during a recent internal audit.
    \item \textbf{Slow Service:} Transactions take far too long (often more than ten minutes) because of the manual ledger search.
    \item \textbf{Poor Decision Support:} The Board cannot answer questions such as: ``What is the loan default rate \emph{this month} in Molyko branch for bike riders?'' or ``Which teller processed the most suspicious transactions last week?''
\end{itemize}

\medskip
\noindent\textbf{Constraints \& Context:}
\begin{itemize}
    \item Electricity is stable in Buea Town and Molyko; Tiko experiences 2--3 hours of daily outages (a generator is available).
    \item Internet: 4G coverage is good at all branches; fibre optic only at Head Office.
    \item Staff: 3 cashiers per branch (a mix of young graduates and older staff with limited IT skills). Budget is tight (non-profit cooperative).
    \item Regulatory: Must comply with COBAC (Central African Banking Commission) requirements on \cle{Know Your Customer (KYC)} and on data retention (records must be kept for many years, typically at least 10).
\end{itemize}

\courssub{Tasks}

Work in teams of 3--4. You have \textbf{3 hours} to complete the analysis and prepare a 5-minute pitch. Use the templates provided.

\medskip
\noindent\textbf{Task 1 -- Organisational Levels \& Decision Mapping (20 marks)}
\begin{enumerate}
    \item Complete the table below by identifying \textbf{two distinct decisions} taken at each of the three organisational levels within Sawa Savings. Classify each decision as \cle{Structured}, \cle{Semi-structured}, or \cle{Unstructured}. Justify your classification in one sentence.
    \item For each level, name the \textbf{type of Information System} (TPS, MIS, DSS, ESS) that primarily supports it.
\end{enumerate}

\begin{center}
\begin{tabularx}{\linewidth}{|l|X|X|l|X|}
\hline
\rowcolor{popblueL}
\textbf{Level} & \textbf{Decision Example} & \textbf{Classification} & \textbf{IS Type} & \textbf{Justification} \\
\hline
Operational &  &  &  &  \\
\hline
Tactical &  &  &  &  \\
\hline
Strategic &  &  &  &  \\
\hline
\end{tabularx}
\end{center}

\medskip
\noindent\textbf{Task 2 -- Data Capture Design (25 marks)}

For each of the following processes, propose the \textbf{most appropriate data capture method(s)}. Justify your choice by referencing \textbf{speed, accuracy, cost, user skill level, and the Cameroonian context} (power, internet, literacy).
\begin{enumerate}
    \item \textbf{New Member Registration} (capturing: Photo, Fingerprint, National ID text, Signature, Form fields).
    \item \textbf{Daily Counter Transaction} (Member identification + Amount entry + Receipt printing).
    \item \textbf{Weekly Branch Report Submission} to Head Office.
\end{enumerate}
\emph{Hint: Consider technologies like OCR, OMR, Barcode/QR Code, Biometric (Fingerprint), Magnetic Stripe/Chip Card, Keyboard/Touchscreen, Mobile App, Sensor.}

\medskip
\noindent\textbf{Task 3 -- Validation Rules \& Verification Strategy (30 marks)}

The proposed digital \textbf{Member Registration Form} has the following fields. For \textbf{each field}, specify:
\begin{itemize}
    \item \textbf{Validation Rule(s)} (type + precise parameter, e.g., ``Range Check: 18--90 years'').
    \item \textbf{One example of INVALID data} that would be rejected.
    \item \textbf{The error message} displayed to the user.
\end{itemize}

\begin{center}
\begin{tabularx}{\linewidth}{|l|X|X|X|}
\hline
\rowcolor{popblueL}
\textbf{Field Name} & \textbf{Validation Rule(s)} & \textbf{Rejected Example} & \textbf{Error Message} \\
\hline
Full Name &  &  &  \\
\hline
National ID No. (13 digits) &  &  &  \\
\hline
Date of Birth &  &  &  \\
\hline
Phone Number (Cameroon format) &  &  &  \\
\hline
Initial Deposit Amount (FCFA) &  &  &  \\
\hline
\end{tabularx}
\end{center}

\medskip
\noindent\textbf{Verification Method:} Propose \textbf{one verification method} for the registration form data entry. Explain \textbf{how} it works and \textbf{why} it is preferred over the alternative for this specific context.

\medskip
\noindent\textbf{Task 4 -- IS Architecture \& Data Flow Diagram (25 marks)}
\begin{enumerate}
    \item Draw a \textbf{Context Diagram (Level 0 DFD)} for the proposed ``Sawa Savings Information System'' (SSIS). Show the system as a single process, all external entities, and major data flows.
    \item Draw a \textbf{Level 1 DFD} decomposing SSIS into \textbf{4--5 major processes} reflecting the TPS $\rightarrow$ MIS $\rightarrow$ DSS/ESS architecture. Label data stores (e.g., \texttt{D1: Member Master File}, \texttt{D2: Transaction Log}). Show data flows between processes and stores.
    \item In a short paragraph, explain how data flows \textbf{from a single deposit transaction at the TPS level} up to a \textbf{strategic dashboard at the ESS level}. Mention \cle{batch processing} vs \cle{real-time processing} where relevant.
\end{enumerate}

\medskip
\noindent\textbf{Task 5 -- The Pitch (Presentation \& Defence) -- Not written, performed.}

Prepare a 5-minute presentation (3 slides max + DFD) covering:
\begin{itemize}
    \item The \textbf{core problem} solved.
    \item The \textbf{proposed architecture} (TPS/MIS/DSS/ESS) and \textbf{key capture/validation choices}.
    \item \textbf{One major risk} (technical, human, or financial) and your mitigation strategy.
\end{itemize}
Be ready to answer 2 questions from the ``Board'' (teacher/class).

\courssub{Model Answer}

\textbf{Task 1 -- Organisational Levels \& Decision Mapping}

\begin{center}
\begin{tabularx}{\linewidth}{|l|X|X|l|X|}
\hline
\rowcolor{popblueL}
\textbf{Level} & \textbf{Decision Example} & \textbf{Classification} & \textbf{IS Type} & \textbf{Justification} \\
\hline
Operational & 1. Approve/Reject a cash withdrawal request against available balance. & Structured & TPS & Pre-defined rules (balance $\ge$ amount, daily limit); repetitive, high volume, algorithmic. \\
\cline{2-5}
 & 2. Assign a new membership number during registration. & Structured & TPS & Automatic sequential allocation; no human judgement required. \\
\hline
Tactical & 1. Set weekly cash float limits for each branch teller. & Semi-structured & MIS & Based on historical withdrawal trends (report) but requires manager judgement for market days and events. \\
\cline{2-5}
 & 2. Approve a loan application above 500,000 FCFA. & Semi-structured & MIS/DSS & A scoring rule exists (structured part) but collateral assessment and the character reference require human judgement. \\
\hline
Strategic & 1. Decide whether to open a 4th branch in Limbe next year. & Unstructured & ESS/DSS & No pre-defined model; relies on long-term trends, competitor analysis, Board vision, and regulatory changes. \\
\cline{2-5}
 & 2. Define the 5-year digital transformation roadmap (Mobile App vs USSD vs Agency Banking). & Unstructured & ESS & High uncertainty, external environment scanning, alignment with cooperative values; no single ``correct'' algorithmic answer. \\
\hline
\end{tabularx}
\end{center}

\medskip
\noindent\textbf{Task 2 -- Data Capture Design}

\begin{enumerate}
    \item \textbf{New Member Registration:}
    \begin{itemize}
        \item \textbf{Method:} \textbf{Biometric Fingerprint Scanner} + \textbf{National ID Card Scanner with OCR} + \textbf{Webcam (Photo)} + \textbf{Touchscreen/Keyboard} (for address, next of kin).
        \item \textbf{Justification:} \textbf{KYC compliance:} the fingerprint uniquely identifies the member, preventing ``ghost members'' and duplicate registration across branches. \textbf{Speed/Accuracy:} OCR scanning of the National ID card auto-fills Name, Date of Birth and ID Number, eliminating keyboard transcription errors (a major source of errors in the current manual system). \textbf{Context:} USB scanners and fingerprint readers are affordable and already familiar in Cameroonian banks; they work offline and data synchronises when the internet returns. A touchscreen is intuitive for younger cashiers; older staff can use a keyboard.
    \end{itemize}
    \item \textbf{Daily Counter Transaction:}
    \begin{itemize}
        \item \textbf{Method:} \textbf{Biometric Fingerprint Verification (one-to-one match)} + \textbf{Numeric Keypad/Touchscreen (Amount)} + \textbf{Thermal Receipt Printer}.
        \item \textbf{Justification:} \textbf{Security:} the fingerprint confirms the member's physical presence (non-repudiation), solving forged-receipt and impersonation fraud. \textbf{Speed:} a one-to-one match takes only a few seconds, versus several minutes searching a paper ledger. \textbf{Cost:} no need to issue and replace expensive smart cards for 12,000 members. \textbf{Power:} low-power USB peripherals run easily on the branch generator/inverter during Tiko outages. Thermal printers are fast and inkless (low running cost).
    \end{itemize}
    \item \textbf{Weekly Branch Report Submission:}
    \begin{itemize}
        \item \textbf{Method:} \textbf{Automated System Generation (scheduled batch job)} $\rightarrow$ \textbf{Encrypted File Transfer (HTTPS)} over 4G/fibre.
        \item \textbf{Justification:} \textbf{Elimination of human error:} the report is now a system-generated summary of the \texttt{Transaction Log} data store; there is no manual compilation, hence no arithmetic errors. \textbf{Timeliness:} it runs automatically every Friday at 18:00. \textbf{Context:} 4G is more than sufficient for small CSV/PDF files. Head Office receives clean data instantly for consolidation --- no more WhatsApp photos of notebooks.
    \end{itemize}
\end{enumerate}

\medskip
\noindent\textbf{Task 3 -- Validation Rules \& Verification Strategy}

\begin{center}
\small
\begin{tabularx}{\linewidth}{|l|X|X|X|}
\hline
\rowcolor{popblueL}
\textbf{Field Name} & \textbf{Validation Rule(s)} & \textbf{Rejected Example} & \textbf{Error Message} \\
\hline
Full Name & \textbf{Presence Check} (not empty). \textbf{Format Check:} letters, spaces, hyphens and apostrophes only. \textbf{Length Check:} 2--50 characters. & \texttt{John123} or \texttt{""} (empty) & ``Name must contain letters only and cannot be empty.'' \\
\hline
National ID No. & \textbf{Presence Check}. \textbf{Format Check:} exactly 13 digits. \textbf{Check Digit Check} (modulus algorithm on the 13th digit). \textbf{Existence/Uniqueness Check:} ID must not already exist in \texttt{Member\_Master\_File}. & \texttt{112345678901} (12 digits) or \texttt{1123456789012} (fails check digit) or an ID already registered & ``Invalid ID (13 digits with valid check digit required) or ID already registered.'' \\
\hline
Date of Birth & \textbf{Presence Check}. \textbf{Format Check:} \texttt{YYYY-MM-DD}, must be a real calendar date. \textbf{Range Check:} member must be at least 18 years old and at most 100. & \texttt{2010-05-15} (underage) or \texttt{1990-02-30} (impossible date) & ``Member must be at least 18 years old. Enter a valid date (YYYY-MM-DD).'' \\
\hline
Phone Number & \textbf{Presence Check}. \textbf{Format Check:} exactly 9 digits, starting with \texttt{6} (Cameroon mobile). \textbf{Uniqueness Check:} recommended if used for SMS alerts/OTP. & \texttt{67712345} (8 digits) or \texttt{377123456} (invalid prefix) & ``Enter a valid 9-digit Cameroon mobile number (e.g., 677123456).'' \\
\hline
Initial Deposit & \textbf{Presence Check}. \textbf{Type Check:} numeric. \textbf{Range Check:} $\ge$ 5,000 FCFA (minimum share capital) and $\le$ 1,000,000 FCFA (threshold above which extra KYC checks apply). & \texttt{-5000} or \texttt{2000} or \texttt{five thousand} & ``Deposit must be a number between 5,000 and 1,000,000 FCFA.'' \\
\hline
\end{tabularx}
\end{center}

\medskip
\noindent\textbf{Verification Method:} \textbf{Double Entry (Double Keying)}.
\begin{itemize}
    \item \textbf{How it works:} The registration form is keyed in \textbf{twice}, by \textbf{two different cashiers} (Cashier A and Cashier B). The system compares the two entries field by field. Any mismatch flags the specific field for correction by a supervisor before the record is committed to the \texttt{Member\_Master\_File}.
    \item \textbf{Why preferred over proofreading:}
    \begin{enumerate}
        \item \textbf{Accuracy:} Proofreading (comparing screen against paper) suffers from confirmation bias --- the eye sees what it expects to see. Double entry by a \emph{different} person forces independent cognitive processing, so the same slip is unlikely to be made twice.
        \item \textbf{Critical data:} The Member Master File is the \cle{single source of truth} for all downstream systems (savings, loans, regulatory reporting). An error here propagates everywhere (wrong ID $\rightarrow$ wrong credit record $\rightarrow$ legal risk).
        \item \textbf{Context:} Sawa Savings has 3 cashiers per branch. During registration peaks, Cashier 1 enters and Cashier 2 re-keys; during quiet periods the Branch Manager acts as the second entry. The cost is staff time (already on payroll) --- no extra hardware is needed.
    \end{enumerate}
\end{itemize}

\medskip
\noindent\textbf{Task 4 -- IS Architecture \& Data Flow Diagram}

\medskip
\noindent\textbf{4.1 Context Diagram (Level 0 DFD)}

\begin{popfigure}
\centering
\begin{tikzpicture}[
    entity/.style={rectangle, draw=popdark, thick, minimum width=2.5cm, minimum height=1.2cm, fill=popblueL, font=\small, align=center},
    process/.style={circle, draw=popdark, thick, minimum size=2.4cm, fill=poporange, font=\small\bfseries, align=center},
    flow/.style={->, thick, draw=popdark, font=\scriptsize, auto}
]
% Entities
\node[entity] (member) at (-5, 2.5) {Member};
\node[entity] (cashier) at (-5, 0) {Cashier / Teller};
\node[entity] (manager) at (-5, -2.5) {Branch Manager};
\node[entity] (board) at (5, 2.5) {Board / GM};
\node[entity] (regulator) at (5, 0) {COBAC / Regulator};
\node[entity, align=center] (mobile) at (5, -2.5){Mobile Money\\Partners};

% Central Process
\node[process] (sys) at (0,0) {SSIS};

% Flows
\draw[flow] (member) -- node[above, align=center] {Registration Details\\Transaction Request} (sys);
\draw[flow] (cashier) -- node[above, align=center] {Credentials\\Transaction Details} (sys);
\draw[flow] (manager) -- node[above, align=center] {Parameters\\Report Requests} (sys);
\draw[flow] (sys) -- node[above, align=center] {Receipts\\Statements\\Alerts} (member);
\draw[flow] (sys) -- node[above, align=center] {Consolidated Reports\\KPIs, Dashboards} (board);
\draw[flow] (sys) -- node[above, align=center] {Regulatory Returns\\Audit Logs} (regulator);
\draw[flow] (sys) -- node[above, align=center] {API Calls\\Settlement Files} (mobile);
\end{tikzpicture}
\legende{Context Diagram (Level 0 DFD) for the Sawa Savings Information System (SSIS)}
\end{popfigure}

\medskip
\noindent\textbf{4.2 Level 1 DFD}

\begin{popfigure}
\centering
\begin{tikzpicture}[
    entity/.style={rectangle, draw=popdark, thick, minimum width=2.2cm, minimum height=1cm, fill=popblueL, font=\scriptsize, align=center},
    process/.style={rectangle, rounded corners, draw=popdark, thick, minimum width=2.6cm, minimum height=1.2cm, fill=poporange, font=\small\bfseries, align=center},
    flow/.style={->, thick, draw=popdark, font=\scriptsize, auto},
    datastore/.style={rectangle, draw=popdark, thick, minimum width=2.4cm, minimum height=0.8cm, fill=popgreenL, font=\scriptsize, align=center}
]
% Entities (Left)
\node[entity] (member) at (-5.5, 3) {Member};
\node[entity] (cashier) at (-5.5, 1) {Cashier};
\node[entity] (manager) at (-5.5, -1) {Branch Mgr};
\node[entity] (board) at (-5.5, -3) {Board / GM};

% Processes (Center)
\node[process, align=center] (p1) at (0, 3){1.0\\Member\\Onboarding};
\node[process, align=center] (p2) at (0, 1){2.0\\Transaction\\Processing (TPS)};
\node[process, align=center] (p3) at (0, -1){3.0\\Branch\\Reporting (MIS)};
\node[process, align=center] (p4) at (0, -3){4.0\\Strategic\\Analysis (DSS/ESS)};

% Data Stores (Right)
\node[datastore] (d1) at (5.5, 3.5) {D1 Member Master};
\node[datastore] (d2) at (5.5, 2) {D2 Account Ledger};
\node[datastore] (d3) at (5.5, 0.5) {D3 Transaction Log};
\node[datastore] (d4) at (5.5, -1) {D4 Loan Portfolio};
\node[datastore] (d5) at (5.5, -2.5) {D5 MIS Warehouse};
\node[datastore] (d6) at (5.5, -4) {D6 Analytics Cube};

% External Entities Right
\node[entity] (reg) at (9.8, 0.5) {COBAC};
\node[entity] (momo) at (9.8, -1) {MoMo Partners};

% Flows: inputs
\draw[flow] (member) -- node[above] {Reg Form, ID, Bio} (p1);
\draw[flow] (cashier) -- node[above] {Tx Details, Auth} (p2);
\draw[flow] (manager) -- node[above] {Params, Approvals} (p3);
\draw[flow] (board) -- node[above] {Queries, Policies} (p4);

% Process to Stores
\draw[flow] (p1) -- node[above] {New Member} (d1);
\draw[flow] (p2) -- node[above] {Update Balance} (d2);
\draw[flow] (p2) -- node[above] {Log Tx} (d3);
\draw[flow] (p2) -- node[above, pos=0.4] {Loan Updates} (d4);
\draw[flow] (p3) -- node[above] {ETL Summaries} (d5);
\draw[flow] (p4) -- node[above] {Aggregates} (d6);

% Store to Process (reads)
\draw[flow] (d1.west) to[bend right=15] node[below, pos=0.6] {Profile} (p1.south east);
\draw[flow] (d2.west) -- node[below, pos=0.55] {Balance} (p2.east);
\draw[flow] (d3.west) -- node[below, pos=0.55] {Raw Tx} (p3.east);
\draw[flow] (d4.west) -- node[below, pos=0.55] {Loan Status} (p3.east);
\draw[flow] (d5.west) -- node[below, pos=0.55] {Clean Data} (p4.east);
\draw[flow] (d6.west) -- node[below, pos=0.55] {KPIs} (p4.east);

% Outputs
\draw[flow] (p1) -- node[above, pos=0.6] {Card, Receipt} (member);
\draw[flow] (p2) -- node[above, pos=0.6] {Receipt, SMS} (member);
\draw[flow] (p3) -- node[above, pos=0.6] {Weekly Reports} (manager);
\draw[flow] (p4) -- node[above, pos=0.6] {Dashboards} (board);
\draw[flow] (p3) -- node[above] {Returns} (reg);
\draw[flow] (p2) -- node[above, pos=0.4] {Settlement} (momo);
\end{tikzpicture}
\legende{Level 1 DFD showing the TPS (2.0), MIS (3.0) and DSS/ESS (4.0) decomposition}
\end{popfigure}

\medskip
\noindent\textbf{4.3 Data Flow Narrative: From Deposit to Dashboard}

\begin{enumerate}
    \item \textbf{TPS (real-time processing):} A member deposits 50,000 FCFA at Molyko. The cashier scans the member's fingerprint $\rightarrow$ Process \texttt{2.0} reads \texttt{D1 (Member Master)} and \texttt{D2 (Account Ledger)} $\rightarrow$ validates the request $\rightarrow$ updates \texttt{D2 (new balance)} and writes a permanent record to \texttt{D3 (Transaction Log)} $\rightarrow$ prints a thermal receipt and sends an SMS alert. \textbf{Latency: a few seconds.} Each transaction is processed immediately and completely (all-or-nothing), so the ledger can never be half-updated.
    \item \textbf{MIS (batch processing):} Every night, an \textbf{ETL (Extract--Transform--Load)} job runs (Process \texttt{3.0}). It reads \texttt{D3 (Transaction Log)} and \texttt{D4 (Loan Portfolio)}, aggregates the data by branch, product, teller and date, and loads clean summaries into \texttt{D5 (MIS Data Warehouse)}. The Branch Manager views the ``Daily Flash Report'' the next morning, and the weekly report to Head Office is produced automatically.
    \item \textbf{DSS/ESS (on-demand analysis):} The General Manager opens the ``Strategic Dashboard'' (Process \texttt{4.0}). The tool queries \texttt{D5} (or the dedicated \texttt{D6 Analytics Cube}) with ad-hoc questions such as: \emph{total loan disbursements per branch for the current month}. The ESS visualises the result as charts and a map with branch bubbles sized by loan volume. The Board sees \emph{this month's} figures --- for example Limbe's potential versus Molyko's saturation --- \emph{before} the quarterly meeting.
\end{enumerate}
\textbf{Key Distinction:} the TPS captures \textbf{atomic events} in real time; the MIS provides \textbf{structured periodic summaries} (``What happened?''); the DSS/ESS enables \textbf{ad-hoc analysis and prediction} (``Why did it happen? What if we\ldots?'').

\medskip
\noindent\textbf{Task 5 -- Pitch Outline (Guidance for Students)}

\begin{methode}[Slide Structure for the 5-Minute Pitch]
\begin{enumerate}
    \item \textbf{Slide 1: The Problem (30 sec).} ``Sawa Savings loses money every year to errors and fraud, and the Board decides blindly on 4-week-old figures. We propose SSIS: a 3-tier architecture (TPS/MIS/ESS) built on biometric capture, double-entry verification, and automated nightly reporting.''
    \item \textbf{Slide 2: The Architecture \& Key Choices (3 min).} Show the Level 1 DFD. Highlight: \textbf{Biometric ID} (solves ghost members and KYC), \textbf{Double Entry} (protects master data integrity), \textbf{Nightly ETL} (solves stale reports), \textbf{Offline-first TPS} (survives Tiko power cuts).
    \item \textbf{Slide 3: Risk \& Mitigation (1 min).} \textbf{Risk:} cashier resistance to biometrics and fear of job losses. \textbf{Mitigation:} phased rollout (Buea Town first), ``super-user'' champions among the cashiers, and clear communication that the system \emph{helps} them (no more balancing ledgers late at night).
    \item \textbf{Defence:} Anticipate questions: ``Cost?'' (an open-source stack --- e.g., PostgreSQL plus Android tablets --- costs far less than the annual losses from errors and fraud). ``Internet down?'' (a local database on each branch machine synchronises when 4G returns).
\end{enumerate}
\end{methode}

\begin{attention}[Common Pitfalls in This Activity]
\begin{itemize}
    \item \textbf{Confusing Validation and Verification:} Validation asks ``Is the data \emph{sensible and reasonable}?'' (range, format, presence, check digit). Verification asks ``Is the data \emph{exactly what was on the source document}?'' (double entry, proofreading). Never list ``range check'' as a verification method.
    \item \textbf{Generic DFDs:} Do not draw a generic ``library system'' DFD. Your external entities \textbf{must} be Member, Cashier, Branch Manager, Board/GM, COBAC and MoMo Partner; your data stores \textbf{must} reflect microfinance (Member Master, Loan Portfolio, Transaction Log).
    \item \textbf{Ignoring the Context:} Do not propose fibre-to-the-desktop in Tiko, or smart cards for 12,000 low-income members (card issuance alone would cost millions of FCFA). Justify every technology choice against the \textbf{Constraints} listed in the situation.
    \item \textbf{Weak Classifications:} ``Setting the interest rate'' is \emph{not} operational --- it is tactical/strategic. Operational means ``apply the \emph{current} rate to \emph{this} loan''.
    \item \textbf{Missing the Check Digit:} A ``format check: 13 digits'' alone cannot catch transposition errors (the number one error in the scenario). You \textbf{must} add a \textbf{check digit check} for the National ID number.
\end{itemize}
\end{attention}

\begin{savaistu}[Cameroon Context: Microfinance \& Technology]
The Cameroonian microfinance sector is supervised at the regional level by \textbf{COBAC} (Commission Bancaire de l'Afrique Centrale), which requires microfinance institutions to identify their customers rigorously (\cle{KYC} --- Know Your Customer) as part of the fight against money laundering across the CEMAC zone. Biometric identification is now common in Cameroonian banks and larger microfinance institutions. Many cooperatives run core-banking software (open-source systems such as \textit{Mifos}, or commercial packages) on local servers, with Android tablets used by field agents, and rely heavily on mobile money (MTN MoMo, Orange Money) as a liquidity channel. The ``Sawa Savings'' scenario mirrors the real digitalisation journey of many cooperatives in the Southwest and Northwest regions, where paper records proved fragile and manual consolidation delayed decisions. The analysis and design skills you practise in this activity --- data capture, validation, verification, TPS/MIS/DSS layering and DFD modelling --- are exactly what employers look for in a junior systems analyst.
\end{savaistu}

\newpage

\coursec{Methods and worked examples}

\coursec{Supporting decision-making in organisations}

\begin{definition}[Information System]
An \cle{information system} (IS) is an organised set of resources — \cle{people}, \cle{hardware}, \cle{software}, \cle{data} and \cle{procedures} — that collects, stores, processes and disseminates information to support the operations and the \cle{decision-making} of an organisation.
\end{definition}

\begin{aretenir}[The five components of an IS]
\begin{itemize}
\item \textbf{Hardware:} physical equipment (computers, servers, scanners, printers, network).
\item \textbf{Software:} programs (OS, DBMS, application packages).
\item \textbf{Data:} raw facts stored and processed (records, transactions, readings).
\item \textbf{People:} users and specialists (clerks, managers, administrators).
\item \textbf{Procedures:} rules and instructions governing use (login, backup, validation rules).
\end{itemize}
\end{aretenir}

\begin{aretenir}[The four functions of an IS]
\begin{itemize}
\item \textbf{Input (capture):} getting data into the system.
\item \textbf{Processing:} transforming data into information.
\item \textbf{Storage:} holding data for future use.
\item \textbf{Output (dissemination):} presenting information to users.
\end{itemize}
\end{aretenir}

\courssub{Method 1 — Identifying the Components of an Information System}

\begin{methode}[Analysing an information system into its components]
When an examination question describes an organisation and asks you to \emph{identify} or \emph{describe} its information system, proceed as follows:
\begin{enumerate}
\item \textbf{Recall the definition.} An information system is an organised set of resources (people, hardware, software, data, procedures) that collects, stores, processes and disseminates information to support the activities and decisions of an organisation.
\item \textbf{Find the five components} in the scenario, one by one:
\begin{itemize}
\item \cle{Hardware}: the physical equipment (computers, servers, barcode scanners, printers, network cables).
\item \cle{Software}: the programs (operating system, database management system, payroll application).
\item \cle{Data}: the raw facts stored and processed (student marks, stock quantities, customer records).
\item \cle{People}: the users and specialists (data-entry clerks, managers, database administrator).
\item \cle{Procedures}: the rules and instructions governing use (login rules, backup schedule, validation rules).
\end{itemize}
\item \textbf{Trace the four functions} of the IS in the scenario: \cle{input} (capture), \cle{processing}, \cle{storage} and \cle{output} (dissemination).
\item \textbf{Conclude} by linking the system to its purpose: supporting operations and \cle{decision-making}.
\end{enumerate}
\textbf{Reflex:} in a ``describe the IS'' question, one mark is usually awarded per correctly named \emph{and} illustrated component — always quote an element of the scenario, never give a bare list.
\end{methode}

\begin{exoresolu}[Analysing the information system of a school]
\textbf{Statement (GCE style).} Bilingual Grammar School Mfoundi uses a computerised system to manage examinations. Teachers enter marks on laptops connected to a server; a program called \emph{MarkMaster} computes averages and ranks; results are stored in a database; the principal consults termly summaries to decide on awards, and clerks print report cards. Rules require each teacher to log in with a personal password and marks to be backed up every evening.
\begin{enumerate}
\item[a)] Define an information system. \hfill (2 marks)
\item[b)] Identify the five components of the school's information system, giving one example of each from the text. \hfill (5 marks)
\item[c)] State the four functions of an information system and illustrate each from the scenario. \hfill (4 marks)
\end{enumerate}
\tcblower
\textbf{Solution.}
\begin{enumerate}
\item[a)] An information system is an organised set of resources — people, hardware, software, data and procedures — that collects, stores, processes and disseminates information in order to support the operations and the decision-making of an organisation.
\item[b)]
\begin{itemize}
\item \textbf{Hardware:} the laptops, the server and the printers used for report cards.
\item \textbf{Software:} the \emph{MarkMaster} program that computes averages and ranks (plus the database software).
\item \textbf{Data:} the marks, averages, ranks and student records stored in the database.
\item \textbf{People:} the teachers (data entry), the clerks (printing), the principal (decision-maker).
\item \textbf{Procedures:} personal password login and the compulsory evening backup of marks.
\end{itemize}
\item[c)]
\begin{itemize}
\item \textbf{Input / capture:} teachers enter marks on laptops.
\item \textbf{Processing:} \emph{MarkMaster} computes averages and ranks.
\item \textbf{Storage:} results are kept in the database on the server.
\item \textbf{Output / dissemination:} termly summaries for the principal and printed report cards.
\end{itemize}
\end{enumerate}
\textbf{Examiner's tip:} in (b), do not write ``hardware: computers'' only — anchor every answer in the scenario to earn full marks.
\end{exoresolu}

\courssub{Method 2 — Choosing an Appropriate Data Capture Method}

\begin{definition}[Data Capture]
\cle{Data capture} is the process of obtaining data from its source and entering it into an information system. It may be \cle{manual} (human keying) or \cle{automated} (direct capture at source without human transcription).
\end{definition}

\begin{aretenir}[Common automated capture technologies]
\begin{itemize}
\item \cle{OCR} (Optical Character Recognition): printed or typed documents.
\item \cle{OMR} (Optical Mark Recognition): shaded bubbles — multiple-choice sheets, lotteries.
\item \cle{Barcode / QR code reader}: product identification in shops, warehouses.
\item \cle{MICR} (Magnetic Ink Character Recognition): bank cheques.
\item \cle{RFID}: contactless identification (tags on goods, toll badges).
\item \cle{Biometric devices}: fingerprints, face — attendance, security.
\item \cle{Sensors / data loggers}: temperature, humidity, traffic — automatic, continuous capture.
\end{itemize}
\end{aretenir}

\begin{methode}[Selecting a data capture technique for a given situation]
To choose and justify a method:
\begin{enumerate}
\item \textbf{Identify the nature of the source data}: printed text, barcode, magnetic stripe, handwritten answers, images, voice, continuous physical quantity.
\item \textbf{Match the data to the technology}:
\begin{itemize}
\item \cle{Keyboard / form}: small volumes of typed data.
\item \cle{OCR}: printed/typed documents.
\item \cle{OMR}: shaded bubbles (MCQ sheets).
\item \cle{Barcode / QR}: product IDs.
\item \cle{MICR}: cheques.
\item \cle{RFID}: contactless tags.
\item \cle{Biometrics}: fingerprints, face.
\item \cle{Sensors}: continuous physical quantities.
\end{itemize}
\item \textbf{Weigh the factors}: volume, speed required, accuracy (automated capture removes transcription errors), cost, environment (dust, distance, mobility).
\item \textbf{Justify} your choice with at least two advantages linked to the scenario, and mention one limitation if asked to evaluate.
\end{enumerate}
\textbf{Key idea:} automated capture is \emph{faster} and \emph{more accurate} at the source, but requires investment; manual capture is cheap but slow and error-prone.
\end{methode}

\begin{attention}[Common trap]
Do not confuse \cle{OCR} (reads \emph{characters}) with \cle{OMR} (reads \emph{marks/bubbles}). OCR is for text; OMR is for multiple-choice answers.
\end{attention}

\begin{savaistu}[Cameroon context]
In Cameroon, \cle{MICR} is used by all commercial banks for cheque clearing through the GIMAC (Groupement Interbancaire Monétique de l'Afrique Centrale) platform. \cle{OMR} is used by the GCE Board for marking multiple-choice papers. \cle{RFID} tags are increasingly used for livestock tracking in the North Region and for toll collection on the Yaoundé–Douala highway.
\end{savaistu}

\begin{exoresolu}[Choosing capture methods for a supermarket and an examination board]
\textbf{Statement.} 
\begin{enumerate}
\item[a)] A new supermarket in Douala wants fast checkouts and automatic stock updates. Recommend a data capture method for products and justify your choice with two advantages. \hfill (3 marks)
\item[b)] The GCE Board must process 60\,000 multiple-choice answer sheets in two weeks. Recommend a capture method, name the device, and give two reasons why it is preferable to manual entry. \hfill (3 marks)
\item[c)] A bank wants to process cheques securely and quickly. Name the appropriate technology and explain briefly how it works. \hfill (2 marks)
\end{enumerate}
\tcblower
\textbf{Solution.}
\begin{enumerate}
\item[a)] \textbf{Method: barcode scanning (with a laser barcode reader at the till).}
Advantages:
\begin{itemize}
\item \textbf{Speed:} the cashier passes the product in front of the reader; the code is read in a fraction of a second, much faster than typing a product code, so queues are shorter.
\item \textbf{Accuracy and integration:} the barcode identifies the product unambiguously; the till automatically retrieves the price and decrements the stock level in the database, eliminating typing errors and keeping stock information up to date in real time.
\end{itemize}
\item[b)] \textbf{Method: Optical Mark Recognition (OMR)}, using an OMR scanner.
Reasons:
\begin{itemize}
\item \textbf{Volume and speed:} an OMR scanner reads hundreds of sheets per minute; 60\,000 sheets can be processed within the deadline, which would be impossible by manual keying.
\item \textbf{Reliability:} OMR detects shaded bubbles directly, removing the transcription errors (misread digits, wrong candidate numbers) inevitable with manual entry of such a volume.
\end{itemize}
\item[c)] \textbf{MICR — Magnetic Ink Character Recognition.} The characters at the bottom of the cheque (account number, cheque number, branch code) are printed with a special magnetic ink. A reader-sorter magnetises and reads these characters at high speed; the magnetic ink can still be read even if the cheque is stamped or slightly dirty, which makes processing fast and resistant to fraud and alteration.
\end{enumerate}
\end{exoresolu}

\courssub{Method 3 — Distinguishing Verification from Validation}

\begin{definition}[Verification]
\cle{Verification} is the process of checking that data entered into the computer \emph{matches the original source document} — i.e.\ data has been \emph{copied correctly}. It does \emph{not} check that the data is sensible.
\end{definition}

\begin{definition}[Validation]
\cle{Validation} is the automatic checking, by the program, that data is \emph{reasonable, complete and within acceptable limits} — i.e.\ data is \emph{sensible}.
\end{definition}

\begin{aretenir}[Verification techniques]
\begin{itemize}
\item \cle{Double entry}: keying data twice and comparing the two versions.
\item \cle{Proofreading / visual check}: comparing screen display with source document.
\item \cle{Sending data back to sender} for confirmation (e.g.\ email confirmation link).
\end{itemize}
\end{aretenir}

\begin{aretenir}[Validation checks — classification]
\begin{itemize}
\item \cle{Presence check}: field not empty.
\item \cle{Type check}: correct data type (numeric, date, etc.).
\item \cle{Range check}: value between min and max.
\item \cle{Format / picture check}: matches a pattern (e.g.\ \texttt{AB1234}).
\item \cle{Length check}: exact or maximum number of characters.
\item \cle{Lookup / membership check}: value exists in a list/table.
\item \cle{Check digit}: extra digit computed from others to detect typing errors.
\item \cle{Consistency (cross-field) check}: two fields agree (e.g.\ DOB vs age).
\end{itemize}
\end{aretenir}

\begin{methode}[Answering ``verification vs validation'' questions]
\begin{enumerate}
\item \textbf{Memorise the definitions.} Verification = ``copied correctly''; Validation = ``sensible''.
\item \textbf{Classify any given check}: if it compares entered data with the source $\Rightarrow$ verification; if it applies a rule about allowed values $\Rightarrow$ validation.
\item \textbf{Give the standard conclusion}: verification and validation are \emph{complementary}; neither guarantees that the original data was true (\emph{garbage in, garbage out}).
\end{enumerate}
\textbf{Reflex:} the phrase ``data is copied correctly but may still be wrong'' signals verification; ``data is rejected because it is impossible or out of range'' signals validation.
\end{methode}

\begin{attention}[Examiner's trap]
A question may describe a \emph{check digit} and ask whether it is verification or validation. A check digit is a \textbf{validation} check (it validates the internal consistency of the number), not a verification technique.
\end{attention}

\begin{exoresolu}[Verification or validation?]
\textbf{Statement.} A hospital in Bamenda computerises patient records. For each situation below, state whether it is \emph{verification} or \emph{validation}, and justify in one sentence.
\begin{enumerate}
\item[a)] Two clerks type the same patient's date of birth from the admission form and the system compares the two versions. \hfill (2 marks)
\item[b)] The system rejects a body temperature of $95\,^{\circ}\mathrm{C}$ entered for a living patient. \hfill (2 marks)
\item[c)] A clerk reads the screen aloud while a colleague checks it against the paper file. \hfill (2 marks)
\item[d)] The system refuses a date of birth of 30/02/2019. \hfill (2 marks)
\item[e)] Explain why, even after both verification and validation, a record may still contain wrong information. \hfill (2 marks)
\end{enumerate}
\tcblower
\textbf{Solution.}
\begin{enumerate}
\item[a)] \textbf{Verification} — double entry: the two typed versions are compared to detect a \emph{copying} (transcription) error between the source form and the computer.
\item[b)] \textbf{Validation} — a range check: $95\,^{\circ}\mathrm{C}$ is outside the plausible range for a living patient (roughly $35$–$42\,^{\circ}\mathrm{C}$), so the program rejects it as unreasonable.
\item[c)] \textbf{Verification} — proofreading/visual check: the entered data is compared with the original source document.
\item[d)] \textbf{Validation} — a data-type/format and range check on the date: February never has 30 days, so the date is invalid.
\item[e)] Verification only ensures the data was \emph{copied faithfully} and validation only ensures it is \emph{plausible}. If the source document itself contained an error (e.g.\ the patient gave a wrong birth year such as 1990 instead of 1980 — plausible and correctly typed), no check can detect it: \emph{garbage in, garbage out}.
\end{enumerate}
\end{exoresolu}

\courssub{Method 4 — Designing Validation Checks}

\begin{methode}[Choosing and writing the correct validation check]
Given a field and its constraints, select the appropriate check:
\begin{enumerate}
\item \textbf{List the constraints} on the field: must it be present? of a fixed type? within limits? from a fixed list? of a fixed pattern? self-consistent?
\item \textbf{Match each constraint to a check}:
\begin{itemize}
\item \cle{Presence check}: the field must not be left empty.
\item \cle{Type check}: the data is of the correct type (numeric, date\ldots).
\item \cle{Range check}: value lies between a minimum and a maximum.
\item \cle{Format / picture check}: data matches a pattern (e.g.\ matricule \texttt{AB1234}).
\item \cle{Length check}: exact or maximum number of characters.
\item \cle{Lookup / membership check}: value exists in a list or table (e.g.\ class code).
\item \cle{Check digit}: an extra digit computed from the others detects typing errors in long numbers (account numbers, barcodes).
\item \cle{Consistency (cross-field) check}: two fields agree (e.g.\ date of birth vs age).
\end{itemize}
\item \textbf{Write the rule precisely}, e.g.\ ``range check: $0 \le \text{mark} \le 20$''.
\item \textbf{Give a test example}: one value accepted, one value rejected, with reasons.
\end{enumerate}
\textbf{Reflex:} for an identifier typed by humans (bank account, barcode), always think of the \emph{check digit}.
\end{methode}

\begin{exoresolu}[Validating a student registration form]
\textbf{Statement.} A school management system records each student with the fields below. For each field, name the most appropriate validation check(s), write the rule, and give one value that would be rejected.
\begin{enumerate}
\item[a)] \textbf{Mark} (out of 20): numeric, compulsory. \hfill (3 marks)
\item[b)] \textbf{Gender}: must be \texttt{M} or \texttt{F}. \hfill (2 marks)
\item[c)] \textbf{Matricule}: exactly 6 characters, two letters followed by four digits (e.g.\ \texttt{CM0451}). \hfill (3 marks)
\item[d)] \textbf{Date of birth} and \textbf{Age} are both entered. \hfill (2 marks)
\end{enumerate}
\tcblower
\textbf{Solution.}
\begin{enumerate}
\item[a)] \textbf{Presence check} (the mark must be entered) \textbf{+ type check} (must be numeric) \textbf{+ range check}: $0 \le \text{mark} \le 20$. \emph{Rejected values:} an empty field (presence), \texttt{"twelve"} (type), $23$ or $-5$ (range).
\item[b)] \textbf{Lookup / membership check} against the list $\{\texttt{M}, \texttt{F}\}$ (often implemented with a drop-down list, which prevents wrong entry altogether). \emph{Rejected value:} \texttt{"Male"} or \texttt{"X"}.
\item[c)] \textbf{Length check}: exactly 6 characters; \textbf{format (picture) check}: pattern \texttt{LLNNNN} (two letters then four digits). \emph{Rejected values:} \texttt{CM451} (length 5), \texttt{1M0451} (first character not a letter).
\item[d)] \textbf{Consistency (cross-field) check}: the age entered must equal the age computed from the date of birth and the current date. \emph{Rejected case:} date of birth 12/03/2008 with age entered as 25 — the two fields contradict each other, so at least one is wrong.
\end{enumerate}
\textbf{Examiner's tip:} always \emph{write the rule} (e.g.\ $0 \le \text{mark} \le 20$), not just the name of the check — this is what earns the second mark.
\end{exoresolu}

\courssub{Method 5 — Computing and Using a Check Digit}

\begin{propriete}[Modulo-11 weighted check digit]
A \cle{check digit} is an extra digit appended to a code to detect transcription errors. In a \textbf{modulo-11 weighted scheme}:
\begin{enumerate}
\item Assign weights to the digits (usually from right to left: $2, 3, 4, 5, \ldots$).
\item Multiply each digit by its weight and sum the products.
\item Divide the sum by 11 and keep the \textbf{remainder} $r$.
\item Compute the check digit: $c = 11 - r$ (if $c = 11$ use $0$; if $c = 10$ use the letter \texttt{X}).
\item To \textbf{validate} a complete number: recompute including the check digit; the total must be a multiple of 11 (remainder $0$).
\end{enumerate}
A single mistyped digit or a transposition of two adjacent digits always changes the remainder $\Rightarrow$ the error is detected.
\end{propriete}

\begin{methode}[Working with the modulo-11 check digit]
\begin{enumerate}
\item Draw a table: digits, weights, products.
\item Compute the sum, divide by 11, find remainder $r$.
\item Check digit $c = 11 - r$ (with conventions for 10 and 11).
\item For validation: recompute on the full code including check digit; remainder must be 0.
\item Always show the full calculation table in your answer.
\end{enumerate}
\end{methode}

\begin{exoresolu}[Applying a modulo-11 check digit]
\textbf{Statement.} A water-supply company in Yaoundé identifies each customer by a 5-digit account number plus a check digit computed with modulo 11, weights $2,3,4,5,6$ applied from right to left to the five digits.
\begin{enumerate}
\item[a)] Compute the check digit for account number \texttt{40572}. \hfill (4 marks)
\item[b)] A clerk enters the number \texttt{40527-6}. Using the check digit, show that this entry is wrong. \hfill (3 marks)
\item[c)] State one type of error, other than a wrong digit, that a check digit detects. \hfill (1 mark)
\end{enumerate}
\tcblower
\textbf{Solution.}
\begin{enumerate}
\item[a)] Digits of \texttt{40572}, weights from right to left $2,3,4,5,6$:
\begin{center}
\begin{tabular}{|c|c|c|c|c|c|}
\hline
\rowcolor{popblueL} Digit & 4 & 0 & 5 & 7 & 2 \\
\hline
\rowcolor{popblueL} Weight & 6 & 5 & 4 & 3 & 2 \\
\hline
Product & 24 & 0 & 20 & 21 & 4 \\
\hline
\end{tabular}
\end{center}
Sum $= 24+0+20+21+4 = 69$. Divide by 11: $69 = 6 \times 11 + 3$, so $r = 3$.
Check digit: $c = 11 - 3 = \mathbf{8}$. The full account number is \texttt{40572-8}.
\item[b)] For the entered number \texttt{40527-6}, recompute the sum on the five digits \texttt{40527}:
\[ 4\times 6 + 0\times 5 + 5\times 4 + 2\times 3 + 7\times 2 = 24+0+20+6+14 = 64. \]
$64 = 5\times 11 + 9$, so $r = 9$ and the correct check digit would be $11-9 = 2$, not $6$. The entered check digit disagrees $\Rightarrow$ the entry is rejected. (In fact the clerk transposed the last two digits: \texttt{40527} instead of \texttt{40572}.)
\item[c)] A \textbf{transposition error} — swapping two adjacent digits (e.g.\ typing \texttt{40527} for \texttt{40572}) — is detected, because the weights are different for each position.
\end{enumerate}
\end{exoresolu}

\courssub{Method 6 — Matching Information Systems to Organisational Levels}

\begin{definition}[Organisational levels of decision-making]
\begin{itemize}
\item \cle{Strategic level} (top management: general manager, board): long-term decisions (3--5 years), unstructured, e.g.\ opening a new branch. Needs summarised, external and internal information; uses \cle{ESS/EIS} (Executive Support System) and \cle{DSS} (Decision Support System).
\item \cle{Tactical level} (middle management: heads of department): medium-term decisions (months--1 year), semi-structured, e.g.\ planning budgets, setting sales targets. Needs summaries and exception reports; uses \cle{MIS} (Management Information Systems) and DSS.
\item \cle{Operational level} (supervisors, clerks): day-to-day decisions, structured, repetitive, e.g.\ recording a sale, restocking. Needs detailed, real-time data; uses \cle{TPS} (Transaction Processing Systems).
\end{itemize}
\end{definition}

\begin{aretenir}[Data flow up the pyramid]
TPS (operational) $\rightarrow$ capture detailed transactions $\rightarrow$ feed MIS (tactical) $\rightarrow$ summarise into reports $\rightarrow$ feed ESS (strategic) $\rightarrow$ high-level summaries, external data, ``what-if'' models.
\end{aretenir}

\begin{aretenir}[DSS ingredients]
A \cle{DSS} (Decision Support System) combines:
\begin{itemize}
\item a \textbf{database} (internal data),
\item a \textbf{model base} (simulation, optimisation, ``what-if'' analysis),
\item a \textbf{user interface} (dialogue management).
\end{itemize}
It supports \emph{semi-structured} decisions where human judgement is needed.
\end{aretenir}

\begin{methode}[Linking the three management levels to their information systems]
\begin{enumerate}
\item \textbf{Draw the decision pyramid} and recall the three levels (see definitions above).
\item \textbf{For any scenario}, ask: \emph{Who decides? Over what horizon? With what detail of information?} Then place the decision at the correct level.
\item \textbf{Remember the data flow}: TPS at the operational level \emph{feed} the MIS, which \emph{feed} the ESS — information is summarised more and more as it goes up the pyramid.
\item \textbf{Know the DSS ingredients}: database, model base, user interface — helping \emph{semi-structured} decisions.
\end{enumerate}
\textbf{Reflex:} ``daily, detailed, repetitive'' $\Rightarrow$ operational/TPS; ``monthly reports, comparisons'' $\Rightarrow$ tactical/MIS; ``long-term, external, unstructured'' $\Rightarrow$ strategic/ESS.
\end{methode}

\begin{exoresolu}[Classifying decisions in a brewery company]
\textbf{Statement.} A brewery company with headquarters in Douala has several depots across Cameroon. Classify each decision below as \emph{strategic}, \emph{tactical} or \emph{operational}, name the type of information system that supports it, and justify briefly.
\begin{enumerate}
\item[a)] The board decides to build a new factory in Bafoussam within five years. \hfill (3 marks)
\item[b)] The sales manager compares monthly sales per depot to fix next quarter's distribution plan. \hfill (3 marks)
\item[c)] A depot clerk records each delivery of crates and updates the stock. \hfill (3 marks)
\item[d)] Explain how the system used in (c) contributes to the decision in (b). \hfill (2 marks)
\end{enumerate}
\tcblower
\textbf{Solution.}
\begin{enumerate}
\item[a)] \textbf{Strategic decision} — long-term (5 years), taken by top management, unstructured. Supported by an \textbf{ESS} (Executive Support System), possibly with a \textbf{DSS} for ``what-if'' simulations (costs, market forecasts), using summarised internal data and external data (market trends, competitors).
\item[b)] \textbf{Tactical decision} — medium-term (a quarter), taken by middle management. Supported by an \textbf{MIS} (Management Information System), which produces periodic summary and comparison reports (monthly sales per depot) from the company's data.
\item[c)] \textbf{Operational decision/activity} — daily, detailed, repetitive, taken at supervisor/clerk level. Supported by a \textbf{TPS} (Transaction Processing System): each delivery is a \emph{transaction} recorded immediately, keeping the stock file accurate and up to date.
\item[d)] The TPS captures \emph{detailed, accurate, real-time} data on every delivery and sale. The MIS \emph{aggregates and summarises} these transactions into monthly sales figures per depot. Without reliable TPS data at the operational level, the tactical reports used by the sales manager would be wrong — decisions at the top depend on correct data captured at the bottom.
\end{enumerate}
\end{exoresolu}

\courssub{Method 7 — Designing a Complete Data Processing Solution (Integration)}

\begin{methode}[Tackling a synthesis (integration) question]
Integration activities combine the whole chapter. Structure your answer in four blocks:
\begin{enumerate}
\item \textbf{Capture}: choose an input method for each data source and justify it (volume, speed, accuracy, cost).
\item \textbf{Control}: propose at least one \emph{verification} technique and two different \emph{validation} checks, each written as a precise rule.
\item \textbf{Processing and storage}: state what is computed, where data is stored (database), and one security measure (passwords, backups, access rights).
\item \textbf{Decision support}: identify who uses the outputs, at which organisational level, and which type of system (TPS, MIS, DSS/ESS) produces them.
\end{enumerate}
\textbf{Reflex:} read the mark allocation and give one developed idea per mark; use the vocabulary of the chapter (\emph{capture, verification, validation, TPS, MIS, DSS}) explicitly.
\end{methode}

\begin{exoresolu}[Integration — computerising a cooperative of cocoa farmers]
\textbf{Statement.} A cocoa cooperative in the Centre Region buys cocoa from 400 farmers. Each bag is weighed and graded (A, B or C) at the depot; payment depends on weight and grade. The cooperative wants a computerised system to record purchases, pay farmers correctly and help its managers plan.
Design the solution by answering:
\begin{enumerate}
\item[a)] Propose a data capture method for the weight of each bag and for the farmer's identity, justifying each choice. \hfill (4 marks)
\item[b)] Give one verification technique and two validation checks (with rules) for the recorded data. \hfill (3 marks)
\item[c)] State the type of information system that records the daily purchases, and the type that produces monthly reports for the manager. Justify. \hfill (3 marks)
\end{enumerate}
\tcblower
\textbf{Solution.}
\begin{enumerate}
\item[a)] \textbf{Weight:} an \textbf{electronic weighing scale connected directly to the computer} (automatic capture by sensor). Justification: the weight is transmitted without being retyped, which removes transcription errors and speeds up the depot queue during the harvest season (high volume of bags).
\textbf{Farmer's identity:} a \textbf{membership card with a barcode (or RFID tag)} read by a reader. Justification: identification is fast and unambiguous — no risk of confusing two farmers with similar names — and the card number can include a \emph{check digit} so that a damaged or misread card is detected immediately.
\item[b)] \textbf{Verification:} the grade and weight displayed on screen are \emph{read back and confirmed against the weighing ticket} by a second clerk (visual check / proofreading) before saving; alternatively, the ticket is printed, given to the farmer, and any dispute triggers a re-weighing (confirmation with the source).
\textbf{Validation checks:}
\begin{itemize}
\item \textbf{Range check on weight:} $1 \le \text{weight} \le 200$ kg per bag — a value such as $0$ or $950$ kg is rejected.
\item \textbf{Lookup check on grade:} grade must belong to the list $\{\texttt{A}, \texttt{B}, \texttt{C}\}$ — any other character is rejected.
\item (Also acceptable: \textbf{presence check} — farmer number, weight and grade must all be entered before the record is saved.)
\end{itemize}
\item[c)] \textbf{Daily purchases:} a \textbf{TPS} (Transaction Processing System): each purchase of a bag is a transaction recorded in real time at the operational level, updating the purchase database and computing each farmer's payment ($\text{payment} = \text{weight} \times \text{price per kg of the grade}$).
\textbf{Monthly reports:} a \textbf{MIS} (Management Information System): it summarises the TPS data (total kilograms bought per village, per grade, payments per farmer) into periodic reports used by the manager at the tactical level to plan purchases, negotiate prices and prepare the budget. A DSS could further simulate ``what-if'' price scenarios before the buying season (strategic support).
\end{enumerate}
\textbf{Conclusion:} reliable decisions at the top of the cooperative rest on a chain: accurate \emph{capture} $\rightarrow$ rigorous \emph{verification and validation} $\rightarrow$ faithful \emph{transaction processing} $\rightarrow$ summarised \emph{management information}.
\end{exoresolu}

\begin{exercice}[Practice questions for Chapter CSC13]
\begin{enumerate}
\item Define the term \emph{information system} and list its five components. \hfill (3 marks)
\item Distinguish between \emph{verification} and \emph{validation}. Give one example of each. \hfill (4 marks)
\item A school uses a 6-digit student ID with a modulo-11 check digit (weights 2,3,4,5,6,7 from right to left). Compute the check digit for ID \texttt{123456}. \hfill (4 marks)
\item Classify the following decisions as strategic, tactical or operational, and name the supporting IS type:
\begin{itemize}
\item The CEO decides to merge with a competitor.
\item The HR manager plans next year's recruitment budget.
\item A cashier processes a customer's payment.
\end{itemize}
\hfill (6 marks)
\item Design a validation rule for a \emph{date of admission} field that must be a valid date not later than today. \hfill (2 marks)
\end{enumerate}
\end{exercice}

\begin{corrige}[Exercise 1]
An information system is an organised set of resources (people, hardware, software, data, procedures) that collects, stores, processes and disseminates information to support organisational operations and decision-making. The five components are: hardware, software, data, people, procedures.
\end{corrige}

\begin{corrige}[Exercise 2]
Verification checks that data has been copied correctly from the source document (e.g.\ double entry). Validation checks that data is reasonable and within allowed limits (e.g.\ range check on age 0–120). Verification ensures accurate transcription; validation ensures sensible content.
\end{corrige}

\begin{corrige}[Exercise 3]
Digits: 1 2 3 4 5 6. Weights from right to left: 7 6 5 4 3 2.
Products: $1\times7=7$, $2\times6=12$, $3\times5=15$, $4\times4=16$, $5\times3=15$, $6\times2=12$.
Sum = 77. $77 \div 11 = 7$ remainder $0$. Check digit $c = 11 - 0 = 11 \rightarrow 0$. Full ID: \texttt{1234560}.
\end{corrige}

\begin{corrige}[Exercise 4]
\begin{itemize}
\item Merger: strategic, ESS/DSS.
\item Recruitment budget: tactical, MIS.
\item Cashier payment: operational, TPS.
\end{itemize}
\end{corrige}

\begin{corrige}[Exercise 5]
Validation checks: \textbf{type check} (must be a valid date format), \textbf{range check} (date $\le$ today's date). Rule: \texttt{admissionDate $\le$ currentDate()}. Rejected: 31/02/2024 (invalid date), 15/06/2026 (future date if today is 2025).
\end{corrige}

\newpage

\coursec{Exercises}

\courssub{I check my knowledge}

\begin{exercice}[Multiple choice questions]
For each question, choose the single correct answer.
\begin{enumerate}
\item An information system is best defined as:
\begin{enumerate}
\item a computer connected to the Internet;
\item an organised set of people, procedures, hardware, software and data that collects, processes, stores and distributes information;
\item a database management system only;
\item a network of printers in an office.
\end{enumerate}
\item Which of the following technologies is an \emph{automated} data capture method?
\begin{enumerate}
\item Typing names on a keyboard;
\item Filling a paper registration form;
\item Reading an examination answer sheet with an OMR scanner;
\item Copying figures from a notebook into a spreadsheet.
\end{enumerate}
\item Checking that a student's age lies between 10 and 25 is an example of:
\begin{enumerate}
\item a presence check;
\item a range check;
\item a format check;
\item proofreading.
\end{enumerate}
\item Which information system is most appropriate for a managing director studying five-year market trends?
\begin{enumerate}
\item TPS;
\item MIS;
\item ESS;
\item OMR.
\end{enumerate}
\end{enumerate}
\end{exercice}

\begin{exercice}[True or false — justify]
State whether each statement is \textbf{true} or \textbf{false}, and justify your answer in one or two sentences.
\begin{enumerate}
\item Verification checks that data has been copied from the source document into the system without transcription errors.
\item A validation check can detect that a mark of 12 has been wrongly typed as 21.
\item Middle managers use mainly tactical information to plan resources over the medium term.
\item A barcode reader captures data manually.
\end{enumerate}
\end{exercice}

\begin{exercice}[Key definitions]
Define each of the following terms in two or three sentences, giving one concrete example for each:
\begin{enumerate}
\item Information system;
\item Data capture;
\item Data verification;
\item Data validation;
\item Decision Support System (DSS).
\end{enumerate}
\end{exercice}

\begin{exercice}[Check-digit calculation]
A microfinance institution in Bamenda assigns each account a 6-digit number plus a check digit computed as follows: multiply the six digits successively by $1, 2, 3, 4, 5, 6$, add the products, and take the remainder of the division of the total by $11$; the check digit is this remainder (write X if the remainder is 10).
\begin{enumerate}
\item Compute the check digit for the account number $415203$.
\item The full number $3071427$ is presented. Determine, showing your working, whether it is valid.
\end{enumerate}
\end{exercice}

\courssub{I practise}

\begin{exercice}[Classifying capture technologies]
Copy the table below into your exercise book and complete it. For each of the eight technologies, state whether the capture is \textbf{manual} or \textbf{automated}, and give one realistic Cameroonian use.
\begin{center}
\begin{tabularx}{\linewidth}{|l|l|X|}
\hline
\rowcolor{popblueL} \textbf{Technology} & \textbf{Manual / Automated} & \textbf{Cameroonian use} \\
\hline
OMR & & \\
\hline
OCR & & \\
\hline
MICR & & \\
\hline
RFID & & \\
\hline
Barcode reader & & \\
\hline
Keyboard & & \\
\hline
Sensor & & \\
\hline
Biometric device & & \\
\hline
\end{tabularx}
\end{center}
\end{exercice}

\begin{exercice}[Validating a mark field]
A secondary school in Yaoundé enters each student's term mark, a whole number out of 20, into a computerised report-card system.
\begin{enumerate}
\item List \textbf{four} distinct validation checks that should be applied to the mark field, and for each one give an example of an entry it would reject.
\item Describe \textbf{one} verification method the school could use when typing the marks.
\item A mark of 12 is wrongly typed as 21. Explain, for each of your four checks, whether it would catch this error, and state the general lesson you draw about the limits of validation.
\end{enumerate}
\end{exercice}

\begin{exercice}[Verification versus validation]
Copy and complete the following comparison table.
\begin{center}
\begin{tabularx}{\linewidth}{|l|X|X|}
\hline
\rowcolor{popblueL} \textbf{Criterion} & \textbf{Verification} & \textbf{Validation} \\
\hline
Purpose & & \\
\hline
When performed & & \\
\hline
Who / what performs it & & \\
\hline
Example 1 & & \\
\hline
Example 2 & & \\
\hline
\end{tabularx}
\end{center}
\end{exercice}

\begin{exercice}[Where should the error have been caught?]
The computerised patient-record system of a hospital in Douala contains the following errors:
\begin{enumerate}
\item A patient's date of birth recorded as 30 February 2019;
\item A patient's weight recorded as 480 kg instead of 48 kg;
\item A name typed as ``NGONO'' although the hospital card shows ``NGONO'' correctly, but the admission clerk typed ``NGNO'' on the second entry screen;
\item A consultation fee of $-5000$ FCFA;
\item A patient whose sex is recorded as ``B''.
\end{enumerate}
For each error, state at which stage (\cle{capture}, \cle{verification} or \cle{validation}) it should have been detected, name the specific technique or check involved, and propose one corrective measure for the hospital.
\end{exercice}

\courssub{I prepare for the GCE}

\begin{exercice}[Structured question — supermarket chain (20 marks)]
SOCAMARKET is a supermarket chain with four branches in Douala. Every evening, each branch's tills send the day's sales to the head office. The chain is computerising its decision-making.
\begin{enumerate}
\item Define the term \emph{information system} and list its five components, illustrating each with an element of SOCAMARKET. \textbf{(5 marks)}
\item Identify, with one concrete example each, an \textbf{operational}, a \textbf{tactical} and a \textbf{strategic} decision taken in this chain. \textbf{(6 marks)}
\item Recommend one appropriate information system (TPS, MIS, DSS or ESS) for each of the three decisions in (b), justifying each choice by the characteristics of the system and of the decision. \textbf{(6 marks)}
\item State two automated data capture technologies suitable for the tills and one advantage of each over manual entry. \textbf{(3 marks)}
\end{enumerate}
\end{exercice}

\begin{exercice}[Structured question — data quality in a school (15 marks)]
A GCE examination centre captures candidates' subject entries from paper registration forms.
\begin{enumerate}
\item Distinguish clearly between \emph{verification} and \emph{validation}. \textbf{(4 marks)}
\item The candidate number is a 6-digit code ending with a check digit. Explain the purpose of a check digit and describe how it is verified at data entry. \textbf{(4 marks)}
\item Describe two verification techniques the centre could apply to the typed registration data, with one advantage and one limitation of each. \textbf{(4 marks)}
\item A candidate's date of birth is typed as 12/13/2005. Name the validation check that rejects this entry and explain why. \textbf{(3 marks)}
\end{enumerate}
\end{exercice}

\begin{exercice}[Essay question (25 marks)]
\emph{``Discuss how information systems support decision-making at the different levels of an organisation, illustrating your answer with a Cameroonian enterprise of your choice.''}

Your essay should:
\begin{enumerate}
\item define an information system and present the three classical organisational levels (operational, tactical, strategic), with the characteristics of the decisions taken at each level;
\item explain, for each level, which type of information system (TPS, MIS, DSS, ESS) is appropriate and why, using a concrete Cameroonian enterprise (bank, mobile-money operator, agro-industrial company, hospital, school, etc.) as a running example;
\item discuss the importance of data quality (capture, verification, validation) for the reliability of the decisions taken;
\item conclude on the limits of information systems in decision-making.
\end{enumerate}
\end{exercice}

\newpage

\coursec{Solutions to the exercises}

\begin{corrige}[Exercise 1 — Multiple choice questions]
\begin{enumerate}
\item \textbf{Answer: (b).} An information system is an \emph{organised} whole combining \cle{people}, \cle{procedures}, \cle{hardware}, \cle{software} and \cle{data} to collect, process, store and distribute information. Options (a), (c) and (d) each describe only a single technical component (a machine, a DBMS, a network), whereas the information system is the socio-technical whole.
\item \textbf{Answer: (c).} Reading an answer sheet with an \cle{OMR} scanner is automated: the machine detects the shaded bubbles directly, without any human keying. Options (a), (b) and (d) all require a person to write or type the data, so they are manual capture methods.
\item \textbf{Answer: (b).} Testing that a value lies between two bounds (here $10 \le \text{age} \le 25$) is precisely a \cle{range check}. A presence check tests that the field is not empty; a format check tests the shape of the data (e.g.\ two digits); proofreading is a verification technique, not a validation check.
\item \textbf{Answer: (c).} Studying five-year market trends is a \emph{strategic}, long-term, weakly structured activity of top management, so the appropriate system is the \cle{ESS} (Executive Support System), which summarises internal data and external information. A TPS handles daily transactions, an MIS produces periodic reports for middle management, and OMR is a capture technology, not an information system.
\end{enumerate}
\end{corrige}

\begin{corrige}[Exercise 2 — True or false — justify]
\begin{enumerate}
\item \textbf{True.} \cle{Verification} is exactly the control that the data entered into the system is identical to the data on the source document; it detects \emph{transcription} (copying) errors, for example by double entry or proofreading.
\item \textbf{False.} The statement claims that validation would catch a transposition error like typing 21 instead of 12. If 21 lies within the accepted range (e.g.\ marks out of 25), no validation check can detect it, because validation only tests \emph{plausibility}, not \emph{accuracy}. In general, validation cannot detect an error that produces a plausible value; only verification (comparison with the source) can.
\item \textbf{True.} Middle managers (heads of department, branch managers) work at the \cle{tactical} level: they plan resources, budgets and staffing over the medium term (months to about a year), using summarised internal information typically supplied by an MIS.
\item \textbf{False.} A barcode reader is an \emph{automated} capture device: the laser scans the bars and converts them directly into digits, with no human keying. Manual capture means a person types or writes the data (keyboard, paper form).
\end{enumerate}
\end{corrige}

\begin{corrige}[Exercise 3 — Key definitions]
\begin{enumerate}
\item \textbf{Information system.} An organised set of people, procedures, hardware, software and data that collects, processes, stores and distributes information to support the activities and decisions of an organisation. \emph{Example:} the report-card system of a college in Yaoundé, which gathers marks, computes averages and produces term reports.
\item \textbf{Data capture.} The operation of collecting raw data from its source and entering it into the information system, either manually (keyboard, paper form) or automatically (OMR, barcode reader, sensors). \emph{Example:} scanning candidates' answer sheets at a GCE examination centre.
\item \textbf{Data verification.} The control that the data entered is identical to the data on the source document; it detects transcription errors. \emph{Example:} two clerks type the same registration form twice (double entry) and the system flags any difference.
\item \textbf{Data validation.} The automatic control, performed by the program at entry, that data is plausible and conforms to defined rules (presence, format, range, type, check digit). \emph{Example:} rejecting a mark of 25 out of 20 by a range check.
\item \textbf{Decision Support System (DSS).} An interactive information system that helps managers solve semi-structured problems by combining data, analytical models and a user-friendly dialogue, allowing ``what-if'' simulations. \emph{Example:} a system used by a bank in Douala to simulate the effect of different interest rates on loan profitability.
\end{enumerate}
\end{corrige}

\begin{corrige}[Exercise 4 — Check-digit calculation]
\begin{enumerate}
\item For the account number $415203$, multiply the digits by weights $1,2,3,4,5,6$ respectively and add:
\begin{align*}
S &= 1\times 4 + 2\times 1 + 3\times 5 + 4\times 2 + 5\times 0 + 6\times 3\\
  &= 4 + 2 + 15 + 8 + 0 + 18 = 47.
\end{align*}
$47 = 11 \times 4 + 3$, so the remainder is $3$. The check digit is $\boxed{3}$ and the full account number is $4152033$.
\item For $3071427$, the six significant digits are $307142$ and the stated check digit is $7$.
\begin{align*}
S &= 1\times 3 + 2\times 0 + 3\times 7 + 4\times 1 + 5\times 4 + 6\times 2\\
  &= 3 + 0 + 21 + 4 + 20 + 12 = 60.
\end{align*}
$60 = 11 \times 5 + 5$, so the correct check digit is $5$, not $7$. The number $3071427$ is therefore \textbf{invalid}: at least one digit was wrongly captured.
\end{enumerate}
\end{corrige}

\begin{corrige}[Exercise 5 — Classifying capture technologies]
\begin{center}
\begin{tabularx}{\linewidth}{|l|l|X|}
\hline
\rowcolor{popblueL} \textbf{Technology} & \textbf{Manual / Automated} & \textbf{Cameroonian use} \\
\hline
OMR & Automated & Marking GCE multiple-choice answer sheets \\
\hline
OCR & Automated & Scanning printed birth certificates at a council registry \\
\hline
MICR & Automated & Reading the code line of cheques in a bank in Douala \\
\hline
RFID & Automated & Electronic toll or access badges at a company gate \\
\hline
Barcode reader & Automated & Checkout tills of a supermarket in Yaoundé \\
\hline
Keyboard & Manual & Typing students' marks into the report-card system \\
\hline
Sensor & Automated & Automatic weather-station readings for agriculture \\
\hline
Biometric device & Automated & Fingerprint attendance control of civil servants \\
\hline
\end{tabularx}
\end{center}
Note that the keyboard is the only manual device in the list: all the others transfer data to the system without human keying, which is precisely what makes them faster and less error-prone.
\end{corrige}

\begin{corrige}[Exercise 6 — Validating a mark field]
\begin{enumerate}
\item Four distinct validation checks on the mark field:
\begin{itemize}
\item \textbf{Presence check:} the field must not be empty — rejects a record where no mark was typed at all.
\item \textbf{Type (numeric) check:} the entry must be a number — rejects ``AB'' or ``12A''.
\item \textbf{Format / integer check:} the mark must be a whole number — rejects $12.5$ or ``12,5''.
\item \textbf{Range check:} the mark must satisfy $0 \le m \le 20$ — rejects $-3$ or $25$.
\end{itemize}
(One could also accept a \emph{length check}, e.g.\ at most two digits, rejecting $123$.)
\item \textbf{Verification method:} double entry — two different operators type the same mark sheet independently and the system compares the two versions, flagging every difference for correction. (Proofreading of the printed entry list against the original mark sheets is an acceptable alternative.)
\item If $12$ is typed as $21$:
\begin{itemize}
\item Presence check: \emph{not caught} — the field is filled.
\item Type check: \emph{not caught} — $21$ is numeric.
\item Format check: \emph{not caught} — $21$ is a whole number.
\item Range check: \textbf{caught} if the range is $0$--$20$, since $21 > 20$; but had the slip produced $11$ instead of $21$, no check would react.
\end{itemize}
\textbf{General lesson:} validation only guarantees that data is \emph{plausible}, not that it is \emph{accurate}. An erroneous value that respects all the rules passes every check; only verification against the source document can guarantee fidelity. Hence a quality data-entry procedure always combines validation \emph{and} verification.
\end{enumerate}
\end{corrige}

\begin{corrige}[Exercise 7 — Verification versus validation]
\begin{center}
\begin{tabularx}{\linewidth}{|l|X|X|}
\hline
\rowcolor{popblueL} \textbf{Criterion} & \textbf{Verification} & \textbf{Validation} \\
\hline
Purpose & Ensure entered data is identical to the source document (detect transcription errors) & Ensure entered data is plausible and conforms to rules (detect unreasonable data) \\
\hline
When performed & During or immediately after data entry, by comparison with the source & Automatically at the moment of entry, by the program \\
\hline
Who / what performs it & A human operator (proofreading) or a second operator (double entry) & The computer, through programmed checks \\
\hline
Example 1 & Double entry of a registration form & Range check on a mark out of 20 \\
\hline
Example 2 & Proofreading a printed list against paper forms & Format check on a date DD/MM/YYYY \\
\hline
\end{tabularx}
\end{center}
\end{corrige}

\begin{corrige}[Exercise 8 — Where should the error have been caught?]
\begin{enumerate}
\item \textbf{30 February 2019:} detected at \cle{validation} by a \emph{format / date-validity check} (February never has 30 days). \emph{Corrective measure:} use a calendar picker that only offers real dates.
\item \textbf{480 kg instead of 48 kg:} detected at \cle{validation} by a \emph{range check} (e.g.\ adult weight between 1 and 300 kg); if 480 had been typed as a plausible value, only \cle{verification} against the weighing record could help. \emph{Corrective measure:} set realistic range limits and capture weight directly from an electronic scale.
\item \textbf{``NGNO'' on the second entry screen:} detected at \cle{verification} by \emph{double entry}, which flags the mismatch between the two versions. \emph{Corrective measure:} keep double entry for names and display the source card beside the entry screen.
\item \textbf{Fee of $-5000$ FCFA:} detected at \cle{validation} by a \emph{range / sign check} (a fee must be $\ge 0$). \emph{Corrective measure:} program a non-negativity rule on all amount fields.
\item \textbf{Sex recorded as ``B'':} detected at \cle{validation} by a \emph{format / membership check} against the allowed list $\{$M, F$\}$. \emph{Corrective measure:} replace free typing by a drop-down list or radio buttons, which makes invalid values impossible at \cle{capture}.
\end{enumerate}
General remark: good interface design (drop-down lists, calendar pickers, connected instruments) prevents many errors at the capture stage itself, before verification and validation are even needed.
\end{corrige}

\begin{corrige}[Exercise 9 — Structured question — supermarket chain (20 marks)]
\textbf{Indicative mark scheme and model answer.}
\begin{enumerate}
\item \textbf{Definition (2 marks):} an information system is an organised set of people, procedures, hardware, software and data that collects, processes, stores and distributes information to support the organisation's activities and decisions. \textbf{Five components (3 marks, $\frac{1}{2}$ + $\frac{1}{2}$ illustration each):}
\begin{itemize}
\item \emph{People:} cashiers, branch managers, IT staff of SOCAMARKET;
\item \emph{Procedures:} the rule that every till transmits its sales each evening;
\item \emph{Hardware:} tills, barcode scanners, servers at head office;
\item \emph{Software:} the sales-management and stock-management programs;
\item \emph{Data:} product references, prices, daily sales figures.
\end{itemize}
\item \textbf{Decisions (2 marks each):}
\begin{itemize}
\item \emph{Operational:} a branch manager decides to reorder 50 cartons of milk because stock is low — short-term, routine, structured, daily.
\item \emph{Tactical:} the commercial director fixes next term's promotional discounts and each branch's monthly sales targets — medium-term, semi-structured.
\item \emph{Strategic:} the general manager decides to open a fifth branch in Bonabéri within five years — long-term, unstructured, based on trends and external information.
\end{itemize}
\item \textbf{Systems (2 marks each):}
\begin{itemize}
\item \emph{Operational $\to$ TPS:} the Transaction Processing System records every sale in real time and updates stock, so it directly triggers the reorder.
\item \emph{Tactical $\to$ MIS:} the Management Information System produces periodic (weekly/monthly) summaries of sales per branch, exactly what middle management needs to set targets and promotions.
\item \emph{Strategic $\to$ ESS:} the Executive Support System aggregates internal summaries with external market data and presents trends graphically, suiting long-term, weakly structured decisions. (A DSS would be accepted for the tactical decision if justified by simulation of promotion scenarios.)
\end{itemize}
\item \textbf{Capture technologies (1.5 marks each):}
\begin{itemize}
\item \emph{Barcode reader:} captures the product reference instantly; advantage over manual entry — much faster and eliminates typing errors on references.
\item \emph{RFID:} reads tagged items without line-of-sight scanning; advantage — several items read at once, reducing checkout time and omissions.
\end{itemize}
\end{enumerate}
\textbf{Examiner expectations:} precise definitions, illustrations drawn from SOCAMARKET (not generic), and justifications that explicitly link the \emph{characteristics of the decision} (horizon, structure, information needed) to the \emph{characteristics of the system}.
\end{corrige}

\begin{corrige}[Exercise 10 — Structured question — data quality in a school (15 marks)]
\textbf{Indicative mark scheme and model answer.}
\begin{enumerate}
\item \textbf{Verification vs validation (4 marks, 2 each).} \emph{Verification} checks that the data typed is identical to the source document; it is performed by people (proofreading, double entry) and detects transcription errors. \emph{Validation} checks that the data is plausible and respects defined rules (presence, type, format, range, check digit); it is performed automatically by the program at entry and detects unreasonable values. The two are complementary: validation cannot detect a plausible transcription error, and verification cannot judge plausibility.
\item \textbf{Check digit (4 marks).} \emph{Purpose (2):} a check digit is a redundant digit computed from the other digits of the code; it detects the most common capture errors (single wrong digit, transposition of two digits) on identifiers that are typed very often. \emph{Verification at entry (2):} the operator types the whole 6-digit number; the program recomputes the check digit from the first five digits using the agreed formula and compares it with the sixth digit typed; if they differ, the number is rejected and must be re-entered.
\item \textbf{Two verification techniques (2 marks each):}
\begin{itemize}
\item \emph{Double entry:} two operators type the same form and the system flags differences. \emph{Advantage:} detects nearly all transcription errors. \emph{Limitation:} costly and slow (work done twice), and useless if the source form itself is wrong.
\item \emph{Proofreading / visual check:} the operator compares the screen or a printed list with the paper forms. \emph{Advantage:} cheap and quick to organise. \emph{Limitation:} unreliable — a tired reader easily overlooks errors.
\end{itemize}
\item \textbf{12/13/2005 (3 marks).} This is rejected by a \emph{format / range check on the date}: in the DD/MM/YYYY format the second field is the month, and a month number must lie between 01 and 12; $13$ is impossible, so the entry is invalid (the candidate presumably meant 12/13 in the American order, i.e.\ 13 December).
\end{enumerate}
\textbf{Examiner expectations:} the distinction in (a) must mention \emph{who performs the control} and \emph{what kind of error} each detects; in (b) the recomputation-and-comparison mechanism must be explicit.
\end{corrige}

\begin{corrige}[Exercise 11 — Essay question (25 marks)]
\textbf{Indicative mark scheme:} introduction and definitions 4 marks; the three levels and their decisions 6 marks; matching systems to levels with a running example 8 marks; data quality 4 marks; conclusion on limits 3 marks.

\textbf{Model plan and expected content.}

\emph{Introduction.} Define an information system as the organised set of people, procedures, hardware, software and data that collects, processes, stores and distributes information. Announce the plan: decisions are taken at three levels, each level needs a different kind of information, hence a different type of system.

\emph{1. The three organisational levels.}
\begin{itemize}
\item \textbf{Operational level} (supervisors, cashiers, clerks): routine, highly structured, short-term (daily) decisions, based on detailed internal data. Example: a cashier or depot clerk of an agro-industrial company such as SOSUCAM deciding whether to issue stock.
\item \textbf{Tactical level} (middle managers): semi-structured, medium-term (months to a year) decisions about resource allocation, based on summarised internal data. Example: a regional manager setting quarterly sales targets.
\item \textbf{Strategic level} (top management): unstructured, long-term (years) decisions, based on aggregated internal data plus external information (market, competitors, regulations). Example: the board deciding to invest in a new factory.
\end{itemize}

\emph{2. Matching systems to levels (running example: a mobile-money operator).}
\begin{itemize}
\item \textbf{TPS} at the operational level: it records every transaction (money transfer, withdrawal) in real time and keeps balances up to date; without it no daily decision is possible.
\item \textbf{MIS} at the tactical level: it aggregates TPS data into periodic reports (transactions per agency, per week) used by middle managers to plan agent liquidity and commissions.
\item \textbf{DSS} for semi-structured problems: simulations (``what if the commission rate changes?'') help managers test scenarios.
\item \textbf{ESS} at the strategic level: dashboards combining internal summaries with external indicators (market share, regulation) support decisions such as launching a new service.
\end{itemize}
Justify each match by the \emph{nature of the information} (detailed vs summarised vs external) and the \emph{time horizon}.

\emph{3. Data quality.} All these systems are only as reliable as their input: ``garbage in, garbage out''. Careful \cle{capture} (automated where possible: barcode, OMR, sensors), \cle{verification} (double entry, proofreading) and \cle{validation} (presence, format, range, check digits) guarantee that the information on which decisions rest is accurate and plausible; a single wrong figure in a TPS can falsify an MIS report and mislead a strategic choice.

\emph{Conclusion — limits.} Information systems supply and organise information but do not decide: strategic decisions remain weakly structured and depend on human judgement, experience and factors the system cannot model (politics, social context). Systems can also fail (breakdowns, cyber-attacks, obsolete data). Hence information systems \emph{support} decision-making at every level but never replace the decision-maker.

\textbf{Examiner expectations:} a structured essay (introduction, developed paragraphs, conclusion), correct use of the vocabulary (TPS, MIS, DSS, ESS; operational, tactical, strategic), one coherent Cameroonian enterprise used throughout, and a balanced conclusion that mentions both the value and the limits of information systems.
\end{corrige}

\newpage

\coursec{Summary sheet}

\begin{popfigure}
\begin{tikzpicture}[
  mindmap, grow cyclic,
  every node/.style={concept, text=white, font=\small},
  root concept/.append style={concept color=popdark, font=\small\bfseries, minimum size=2.6cm},
  level 1 concept/.append style={sibling angle=51.4, level distance=3.4cm, minimum size=2.2cm},
  level 2 concept/.append style={sibling angle=45, level distance=2.2cm, minimum size=1.7cm, font=\scriptsize},
]
\node[root concept]{Information Systems \& Decision-Making}
  child[concept color=popblue]{ node {IS concepts \& components}
    child[concept color=popblue!75]{ node {Data, information, knowledge} }
    child[concept color=popblue!75]{ node {People, hardware, software, data, procedures, network} }
  }
  child[concept color=poporange]{ node {Data capture}
    child[concept color=poporange!75]{ node {Manual: forms, keyboards} }
    child[concept color=poporange!75]{ node {Automatic: barcode, RFID, OMR, OCR, MICR, sensors} }
  }
  child[concept color=popgreen]{ node {Verification \& validation}
    child[concept color=popgreen!75]{ node {Verification = accurate transcription} }
    child[concept color=popgreen!75]{ node {Validation = data is sensible \& correct} }
  }
  child[concept color=poppurple]{ node {Validation checks}
    child[concept color=poppurple!75]{ node {Range, type, length, presence, format, check digit} }
  }
  child[concept color=popgold]{ node {Organisational levels}
    child[concept color=popgold!75]{ node {Operational, tactical, strategic} }
  }
  child[concept color=poppink]{ node {IS per level}
    child[concept color=poppink!75]{ node {TPS, MIS, DSS, ESS} }
  }
  child[concept color=popmuted]{ node {Decision-making}
    child[concept color=popmuted!75]{ node {Structured to unstructured decisions} }
  };
\end{tikzpicture}
\legende{Mind map of the chapter: supporting decision-making in organisations.}
\end{popfigure}

\begin{bilan}
\textbf{Key definitions.}
\begin{itemize}
  \item \cle{Data}: raw, unprocessed facts and figures (numbers, text, images) with no context.
  \item \cle{Information}: data that has been processed, organised and given meaning so it is useful for decisions.
  \item \cle{Information system (IS)}: an organised combination of \textbf{hardware, software, data, people, procedures and networks} that collects, stores, processes and distributes information to support operations and decision-making in an organisation.
  \item \cle{Data capture}: the process of collecting data and entering it into a computer system, either manually or automatically.
  \item \cle{Verification}: checking that data entered matches the original source (accurate transcription). Methods: \textbf{double entry} (typing twice and comparing) and \textbf{proofreading / visual check}.
  \item \cle{Validation}: automatic checking by software that data is \emph{reasonable, sensible and within acceptable limits}. Validation does \textbf{not} guarantee correctness.
\end{itemize}

\textbf{Data capture technologies to know.}
\begin{itemize}
  \item \textbf{Manual}: keyboard, mouse, touch screen, paper forms.
  \item \textbf{Automatic / direct}: \cle{barcode reader} (shops, stock), \cle{QR code}, \cle{RFID} (tags read by radio, no line of sight), \cle{OMR} (optical mark recognition — multiple-choice answer sheets), \cle{OCR} (optical character recognition — printed text), \cle{MICR} (magnetic ink — bank cheques), \cle{magnetic stripe / chip cards}, \cle{biometric} devices (fingerprint), \cle{sensors} (automatic data logging).
  \item Advantages of automatic capture: faster, fewer transcription errors, cheaper per record in volume.
\end{itemize}

\textbf{Validation checks (learn them all — very frequent at the GCE).}
\begin{itemize}
  \item \textbf{Presence check}: a required field must not be empty.
  \item \textbf{Type check}: data is of the right type (e.g.\ a number, not letters).
  \item \textbf{Range check}: value lies between limits (e.g.\ a mark $0 \le m \le 20$).
  \item \textbf{Length check}: correct number of characters (e.g.\ a phone number of 9 digits).
  \item \textbf{Format check}: data follows a pattern (e.g.\ date \texttt{DD/MM/YYYY}, matricule code).
  \item \textbf{Check digit}: an extra digit computed from the others (modulus arithmetic, e.g.\ ISBN) to detect typing errors.
  \item \textbf{Lookup / existence check}: value must appear in an allowed list or table.
\end{itemize}

\textbf{Organisational levels and their information systems.}
\begin{center}
\begin{tabularx}{\linewidth}{|l|X|X|}\hline
\rowcolor{popblueL} \textbf{Level} & \textbf{Decisions} & \textbf{IS used} \\ \hline
Strategic (top management) & Long-term, unstructured (e.g.\ open a new branch in Douala) & ESS (Executive Support System), DSS \\ \hline
Tactical (middle management) & Medium-term, semi-structured (budgets, staffing) & MIS (Management Information System), DSS \\ \hline
Operational (supervisors, staff) & Daily, structured, repetitive (sales, payroll) & TPS (Transaction Processing System) \\ \hline
\end{tabularx}
\end{center}
\begin{itemize}
  \item \cle{TPS}: records routine transactions; fast, accurate, large volumes.
  \item \cle{MIS}: summarises TPS data into regular reports for middle managers.
  \item \cle{DSS}: interactive ``what-if'' analysis using models and data.
  \item \cle{ESS}: dashboards, external + internal data, trends for executives.
\end{itemize}

\textbf{Methods.}
\begin{itemize}
  \item To choose a capture method: consider \textbf{volume, speed, accuracy, cost, environment} of the data source.
  \item To answer a validation question: identify the field, then name the \emph{exact} check and justify it with an example of rejected data.
  \item To place an IS: ask \emph{who uses it} and \emph{what decision} it supports, then match to operational / tactical / strategic.
\end{itemize}

\textbf{Common mistakes.}
\begin{itemize}
  \item Confusing \textbf{verification} (data copied correctly) with \textbf{validation} (data is sensible). A wrong but plausible value passes validation.
  \item Saying validation ``ensures data is correct'' — false: $99$ in a field validated $0$--$100$ is valid but may be wrong.
  \item Mixing up OMR (marks) and OCR (characters); MICR is for cheques, not barcodes.
  \item Attributing strategic decisions to a TPS: a TPS only processes daily transactions.
\end{itemize}

\textbf{GCE tips.}
\begin{itemize}
  \item Learn one Cameroonian example for each level (e.g.\ a TPS at a Yaound\'e supermarket till, an MIS producing termly results for a school principal, an ESS for a bank's board).
  \item For ``state a suitable validation check'' questions, always \textbf{name the check + give the rule + an example}.
  \item Draw the three-level pyramid (operational at the base, strategic at the top) if asked to describe organisational structure: a labelled diagram earns marks.
  \item Define terms precisely before applying them; examiners reward correct vocabulary (TPS, MIS, DSS, ESS).
\end{itemize}
\end{bilan}

\findecours

\end{document}
